\documentclass[14pt]{article} 
\usepackage[utf8]{inputenc}
\usepackage{epsfig}
\usepackage[ngerman]{babel}
\usepackage{mathtools,extarrows}
\usepackage[thmmarks]{ntheorem}
\usepackage{amsmath,theorem}
\usepackage{amsfonts}
\usepackage{amssymb}
\usepackage{txfonts}
\usepackage{marvosym}
\usepackage{MnSymbol}
\usepackage{mathrsfs}
\usepackage{parallel}
\usepackage{tikz}
\usepackage{geometry}
\usepackage{titletoc}
\usepackage{url}
\usepackage{nicefrac}
\usepackage[pdftex,bookmarks=true,bookmarksnumbered=true,bookmarksopen=true,colorlinks=true,
		filecolor=black,linkcolor=blue,urlcolor=blue,plainpages=false,pdfpagelabels,
		citecolor=black,setpagesize=false]{hyperref}


\makeatletter
\renewcommand\l@section{\@dottedtocline{1}{0em}{2.5em}}
\renewcommand\l@subsection{\@dottedtocline{2}{2.5em}{2em}}
\makeatother

\titlecontents{chapter}[1.5em]{\addvspace{1em}\bfseries}
	{\contentslabel{1.5em}}{\hspace^{-1.5em}}
	{\hfill\contentspage}
\dottedcontents{section}[3.8em]{}{1.3em}{1pc}
\dottedcontents{subsection}[6.0em]{}{2.2em}{1pc}
\dottedcontents{subsubsection}[9.2em]{}{3.1em}{1pc}

%%%%theoremenvironments
\theoremstyle{plain} 

\theoremheaderfont{\bfseries}

\theorembodyfont{\it}
\newtheorem{theorem}{Theorem}[section]
\newtheorem{lem}[theorem]{Lemma}
\newtheorem{cor}[theorem]{Corollar}
\newtheorem{prop}[theorem]{Proposition}

\theorembodyfont{\rm}
%\newtheorem{remark}[theorem]{Bemerkung}
\newtheorem{defi}[theorem]{Definition}
\newtheorem{rul}[theorem]{Regel}

\theoremstyle{nonumberplain}

\theoremheaderfont{\bfseries}
\theorembodyfont{\normalfont}
\newtheorem{remark}{Bemerkung:}
\newtheorem{beispiel}{Beispiel:}

\theoremheaderfont{%
    \normalfont\scshape}
\theorembodyfont{\normalfont}
\theoremsymbol{\square}
\newtheorem{proof}{Beweis:}
\newtheorem{idea}{Beweisidee:}

\topmargin=-2mm \textwidth=155mm \textheight=230mm
\oddsidemargin=0mm \evensidemargin=0mm
\thispagestyle{empty}
\newcommand{\mb}{\mathbb}
\newcommand{\mc}{\mathcal}
\newcommand{\mf}{\mathfrak}
\newcommand{\mr}{\mathrm}
\newcommand{\ms}{\mathscr}
\newcommand{\K}{\mb{K}}
\newcommand{\C}{\mb{C}}
\newcommand{\R}{\mb{R}}
\newcommand{\Q}{\mb{Q}}
\newcommand{\Z}{\mb{Z}}
\newcommand{\N}{\mb{N}}
\newcommand{\ra}{\rightarrow}
\newcommand{\LRa}{\Leftrightarrow}
\newcommand{\Ra}{\Rightarrow}
\newcommand{\La}{\Leftarrow}
\newcommand{\eps}{\varepsilon}
\newcommand{\vt}{\vartheta}
\newcommand{\vp}{\varphi}
\newcommand{\qq}{\hspace{.25cm}}
\newcommand{\ol}{\overline}
\newcommand{\ul}{\underline}
\renewcommand{\d}{ \, \mr{d}}
\newcommand{\Res}{\mr{Res}}
\newcommand{\Const}{\mr{Const} \, }
\newcommand{\Id}{\mr{Id}}
\renewcommand{\star}{\filledstar}
\newcommand{\Gl}{\mr{Gl}}
\newcommand{\sign}{\mr{sign}\,}
\newcommand{\dist}{\mr{dist}}
\newcommand{\diag}{\mr{diag}}
\newcommand{\CP}{\mr{CP}}
\newcommand{\CV}{\mr{CV}}
\newcommand{\supp}{\mr{supp}\,}
\newcommand{\RV}{\mr{RV}}
\renewcommand{\div}{\mr{div}\,}
\newcommand{\conv}{\mr{conv}}
\newcommand{\diam}{\mr{diam}}
\newcommand{\CS}{\mr{CS}}
\newcommand{\Bild}{\mr{Bild}}
%\renewcommand{\span}{\mr{span}}
\newcommand{\lal}{\langle}
\newcommand{\ral}{\rangle}
\newcommand{\esssup}{\mr{ess\, sup}}
\geometry{hmargin=3cm,top=2cm,bottom=4cm}

\newcommand{\makelicense}{Dieser Mitschrieb steht unter der freien \texttt{CC-BY-SA-DE 3.0} 
Lizenz.
\begin{center}\includegraphics[height=31pt,width=88pt]{license.pdf}\end{center}
Für weitere Informationen besuchen Sie
\begin{center}\url{http://creativecommons.org/licenses/by-sa/3.0/deed.de}\end{center}}


\input xy
\xyoption{all}

\begin{document}


\begin{titlepage}
\[ \]
\begin{center}

\hrule\hspace*{1cm\\[3mm]}

\textsc{\Huge Nichtlineare Funktionalanalysis }\\[1cm]
\hrule \hspace*{1cm}\\[3mm]
{\sc Kompletter Mitschrieb zur gleichnamigen Vorlesung im Sommersemester 2011}\\[0.5cm]
    \vspace{1cm}
\begin{minipage}{0.4\textwidth}
\begin{flushleft} \large
\emph{Dozent:}\\
    Dr.~Patrick Dondl
\end{flushleft}
\end{minipage}
\begin{minipage}{0.4\textwidth}
\begin{flushright} \large
\emph{Stand:} \\
    \today
\end{flushright}
\end{minipage}
\[ \]\\[1.5cm]
Luis Felipe Müller
\[ \]\\[1.5cm]
Institut für Angewandte Mathematik\\
Universität Heidelberg
     
\end{center}
\vfill
\noindent\makelicense

\end{titlepage}

%\vfill
%\listoffigures
%\vfill
%\listoftables

\newpage

\pagenumbering{roman}

\section*{Organisatorisches}

\subsection*{Termine}
\begin{itemize}
    \item Vorlesung: Mo./Mi. 9-11ct. -104 Ang. Math.
    \item Übungsaufgaben: Mi. - Mi. vor der Vorlesung. Kasten in der Angew.Math
    \item Übungsgruppe: Fr. 16-18 Angew. Math. -101
    \item Website zur Vorlesung: \url{http://dondl.org/wiki/Sommersemester_11}
    \item Literatur: \begin{enumerate}
        \item Růžička M: Nichtlineare Funktionalanalysis, Eine Einführung
        \item Aubin-Ekeland: Applied nonlinear Analysis
        \item Deimling: Nonlinear Functional Analysis
        \item Schwartz: Nonlinear Functional Analysis
        \item Zeidler: Nonlinear Functional Analysis and its applications 
        \end{enumerate}
    \item Prüfungen: Zulassung mit etwa 50\% der Übungsaufgaben-Punkte. Prüfung ist mündlich, 
            beispielsweise am 29. Juli (Fr.) 
    \item Dozent: Patrick Dondl, Sprechstunde Mo./Mi. 11-12 in Raum 130 (Angew. Math.)
    \item Tutor: Julian Scheuer
\end{itemize}

\newpage

\tableofcontents

\newpage

\pagenumbering{arabic}

\chapter{Einleitung}

\section{Thema der Vorlesung}

In der linearen Funktionalanalysis haben wir eine Vielzahl von Methoden kennengelernt um
Ergebnisse aus der endlichdimensionalen linearen Algebra auf den unendlichdimensionalen Fall 
zu verallgemeinern. Ein Hauptaufgabe war dabei, die Lösbarkeit von Gleichungen der Form
\[
    Ax=y
\]
für lineare Operatoren $A$ auf $\infty$-dimensionalen Banchräumen zu zeigen.
\begin{description}
    \item{Ein Beispiel:}
    
    Sei $\Omega\subset\R^3$, gefüllt mit einer inkompressiblen, riskosen Flüssigkeit.
    $v_j(x)$ sei die Geschwindigkeit der Flüssigkeit an der Stelle $x\in \Omega$. 
    $p(x)$ ist der Druck an der Stelle $x$.
    \begin{description}
    \item{Randbedingungen:}
    $v_j(x)=0 \qq x\in \partial\Omega$
    \item{Inkompressibilität:}
    $\partial_jv_j(x)=0 \qq x\in \Omega$
    \item{Bewegungsungleichung:}
    Wir betrachten die Kräfte, die auf einen kleinen Würfel, eingeschlossen durch
    $(x_1,x_2,x_3)$, $(x_1+\Delta x_1, x_2+\Delta x_2, x_3+\Delta x_3)$. Druck auf eine Oberfläche
    des Würfels mit Normale $x$, ist $f_j^p=p\cdot \Delta x_2 \Delta x_3 \cdot \delta_{ij} $ auf die
    gegenüberliegende Seite wirkt
    \[
        -(p+\partial _i p \Delta x_i) \Delta x_2 \Delta x_3 \delta_{ij}
    \]
    \end{description}
Zusammen ergibt sich 
\[
    f=(\partial_j p)\Delta V
\]

Kraft durch Viskosität auf eine Oberfläche mit Normale $x_1$ ist 
\[
    f_j^{V,x_i}= -2\eta \Delta x_2\Delta x_3 \partial_1v_j
\]
mit einer Konstante $\eta$. Der gleiche Trick wie oben ergibt für die gegenüberliegende
Oberfläche
\[
    \eta \Delta x_2 \Delta x_3 \partial (v_j+\partial v_j\Delta x_1)
\]
Zusammen ergibt sich
\begin{align*}
    f_j^V&= \eta\Delta V \cdot \partial_i\partial_iv_j\\
    \text{Newton:} \qq \rho \Delta V \frac{\d}{\d t} v_j (t,x(t))
        &= \eta \Delta V  \partial _i\partial_iv_j- \Delta V (\partial _j p) + \Delta 
        V\kappa_j \qq \text{mit einer externen Konstante $\kappa_j$, bspw. Gravitation}
\end{align*}
Teilen durch $\Delta V$ und die Kettenregel ergibt
\begin{align*}
    \rho \partial_t v_j&= \eta \partial_i\partial_i v_j - \rho (v_i\partial_i) - \partial_j p + K_j 
        \qq \text{(Newton)}\\
    d_jv_j&=0
\end{align*}
Naurier-Stokes. Frage: Existiert eine eindeutige Lösung zur sationären Naurier-Stokes-Gleichung:
\begin{align}\label{1}
    \eta \partial_i\partial_iv_j-(v_i\partial_i)v_j+\partial_jp+K_j=0
\end{align}
\end{description}
Wir können die Gleichung etwas umschreiben. Sei $H$ ein Hilbertraum
\[
    H:=\ol{\left\{ u\in C^\infty_c (\Omega,\R^3), \text{ so dass } \partial_jv_j=0 \right\}}
        ^{W^{1,2}(\Omega, \R^3)}
\]
Ein Skalarprodukt auf $H$ ist gegeben durch
\begin{align*}
    (u,v)_H:= \int_{\Omega} \nabla u\cdot \nabla v \d x = \int_\Omega \sum_{i,j=1}^3 
    (\partial_iu_j)(\partial_iv_j)
\end{align*}
$\Omega$ beschränkt $\Ra$ $(\cdot, \cdot)_H$ ist äquivalent zum üblichen Skalarprodukt 
(mittels Poincaré).
Wir multiplizieren (\ref{1}) mit $\omega \in H $, integrieren und erhalten
\begin{align*}
    \int_\Omega (\eta \partial_i\partial_i v_j- (v,\partial_i) v_j + K_j) \cdot \omega_j 
    &= \int_\Omega (\partial _jp)\omega_j=0 \text{, da $\omega$ divergenzfrei.}\\
    (\ref{1}) \Ra \eta (v,\omega)_H-a(v,v,w)-\int_\Omega K\omega&=0
\end{align*}
Ebenso für $a$:
\begin{align*}
    a(u,v,w):= (\underbrace{B(u,v)}_{\text{bilinear.}},w)_H
\end{align*}
Also
\begin{align*}
    (\ref{1}) \Ra \, (\eta v - B(u,v)) - \tilde K , \omega)_H=0 \qq \forall \omega \in H
\end{align*}
somit
\[
    \eta v - B(v,v) = \tilde K
\]
Das ist eine Gleichung der Form
\begin{align}\label{2}
    Fv=\tilde K,\qq\text{mit $F$ einem Nichtlinearen Operator}
\end{align}
Im ersten Teil der Vorlesung beschäftigen wir uns mit der eindeutigen Lösbarkeit von
Gleichungen der Form
\[
    Fx=y, \qq F:X\ra Y, \qq X,Y\qq \text{Banachräume}
\]
und zum Abschluß zeigen wir mit Hilfe des \textit{Schauder'schen Fixpunktsatzes} die Existenz
und finden eine Lösung von (\ref{2}), also der schwachen Form der stationären 
Navier-Stokes-Gleichung.

Im zweiten Teil der Vorlesung beschäftigen wir uns mit Variationsrechnung (d.h. dem Finden
von Minimierern nichtlinearer Funktionalen)
\[
    W:X\ra \R \qq\text{mit} \qq X \qq \text{ein Banachraum}
\]
Finde 
\[
    x_0\in X: \, W(x_0)= \inf_{y\in X}W(y)
\]
Insbesonder treffen wir dort auf Probleme in der Elastizitätstheorie.

\section*{Aufbau der Vorlesung}
\begin{itemize}
    \item Abbildungsgrad $\ra$ Existenz von Lösungen von $Fx=y$
    \item Monotone Operatoren $\ra$ Eindeutigkeit von Lösungen von $Fx=y$; 
    	zeitabhängige Probleme.
    \item Variationsrechnung $\ra$ $\inf_{y\in X} W(y)$
\end{itemize}
   % 1 Einleitung \\ 1.1 Thema der Vorlesung
\subsection{Vorarbeiten}

\subsubsection{Ableitung in Banachräume und implizite Funktionen}

Es seien $X$ und $Y$ Banachräume, $\Omega\subset X$ offen, $F:\Omega \ra Y$, $x_0\in \Omega$

\begin{defi}[Gâteaux-Ableitung]
    Die Gâteaux-Ableitung $\d F(x_0,\psi)$ des Operators $F$ im Punkt $x_0$ in Richtung $\psi\in X$ 
    ist gegeben durch 
    \[
        \d F (x_0,\psi)=\lim_{s\ra 0} \frac{F(x_0+s\cdot \psi)-F(x_0)}{s}= \left. \frac{\d}{\d s} 
        F(x_0+s\psi)\right|_{s=0}
    \]
    falls der Limes existiert. Der Operator $F$ heißt in diesem Fall in $x_0$ Richtung $\psi$ 
    Gâteaux-differenzierbar.
\end{defi}
\begin{defi}[Fréchet-Ableitung]
    Der Operator $F$ heißt Fréchet-differenzierbar in $x_0\in \Omega$, falls ein beschränkter,
    linearer Operator 
    \[
        F'(x_0):X\ra Y 
    \]
    existiert, so dass 
    \begin{align}\label{3}
        \lim_{\|h\|\ra 0} \frac{\| F(x_0+h) - F(x_0) -F'(x_0)\cdot h \|  }{\|h\|} = 0
    \end{align}
    $F'(x_0)$ heißt dann Fréchet-Ableitung von $F$ in $x_0$.
\end{defi}

\begin{theorem}
    \begin{description}
        \item{i)} $F'(x_0)$ ist durch (\ref{3}) eindeutig bestimmt.
        \item{ii)} Falls $F$ stetig ist in $x_0$, so ist jeder lineare Operator, der (\ref{3}) erfüllt,
                ebenfalls stetig.
        \item{iii)} Ist $L:X\ra Y$ linear, so gilt
                \[
                    L'(x)=L \qq \forall x\in X
                \]
    \end{description}
\end{theorem}

\begin{proof}

\begin{description}

\item{i)} 
Es gelte (\ref{3}) auch für L. Dann haben wir
\begin{align*}
    \|Lh-F'(x_0)h\| \leq \eps \|h\| \qq \text{falls} \qq \|h\| < \delta = \delta(\eps, x_0)
\end{align*}
Für beliebiges $h$ folgt aber
\begin{align*}
    \|(L-F'(x_0)) (\delta\|h\|^{-1}\cdot h)\|&\leq \delta\eps\\
    \Ra \, \|(L-F'(x_0))h  \|&\leq \eps \|h\| \qq \forall \, h\in X, \, \forall \, \eps>0 \\
    \Ra \, \|L-F'(x_0)\|_{\ms L(X,Y)} &= 0
\end{align*}

\item{ii)}
(\ref{3}) wird umgeformt zu
\[
    \|F'(x_0) h \| \leq \eps \|h\| + \|F(x_0+h)-F(x_0)\|
\]
Mit $h\ra 0$ folgt die Stetigkeit (für $\|h\|\leq \delta$) von $F'(x_0)$ an der Stelle 0. 
Wegen Linearität von $F'(x_0)$ ist $F'(x_0)$ somit stetig. 

\item{iii)} (\ref{3}) gilt offensichtlich für $L'(x_0)=L$, mit i) folgt Eindeutigkeit.
\end{description}
\[ \]
\end{proof}


\begin{prop}
    Jeder Fréchet-differenzierbare Operator $F$ ist Gâteaux-differenziebar $\forall \, \psi\in X$ 
    und es gilt
    \[
        F'(x_0)\psi=\d F(x_0,\psi)
    \]
\end{prop}
\begin{proof}
    Übungsaufgabe \[\]
\end{proof}

\begin{defi}
    $F$ heißt (Fréchet-)differenzierbar auf $\Omega$, falls $\forall\, x\in X$ ein $F'(x)$ existiert,
    sodass $F'(x)$ stetig ist  und (\ref{3}) erfüllt. $F$ heißt stetig (Fréchet-)differenzierbar in
    $\Omega$, falls die Abbildung
    \[
        F': \Omega\ra \ms L(X,Y)
    \]
    stetig ist.
\end{defi}

\begin{prop}
    Existiert die Gâteaux-Ableitung $\d F(x,\psi) \, \forall x\in \Omega$, und ist sie linear und
    stetig in $\psi$ $\forall\, x\in \Omega$, so ist $F$ Fréchet-differenziebar auf $\Omega$ und
    es gilt 
    \[
        F'(x)\psi=\d F(x,\psi)
    \]
\end{prop}

\begin{proof} 
    Übungsaufgabe \[\]
\end{proof}

\begin{defi}
    Sei $F$ auf $\Omega$ stetig differenziebar, $x_0\in \Omega$. Falls ein stetiger linearer Operator
    \[
        F''(x_0):X\ra \ms L(X,Y)
    \]
    existiert mit
    \[
        \lim_{\|h\|\ra 0}\frac{\|F'(x_0+h)-F'(x_0)-F''(x_0)h\|_{\ms L(X,Y)}}{\|h\|}=0
    \]
    dann heißt $F$ in $x_0$ zweimal (Fréchet-)differenzierbar und $F''(x_0)$ heißt zweite Ableitung
    von $F$ in $x_0$. Höhere Ableitungen entsprechend.
\end{defi}

\begin{remark}
    Es gilt die Kettenregel:
    Seien $X$, $Y$, $Z$ Banachräume, $\Omega_X\subset X$ offen, $x_0\in \Omega_X$,
    \[
        F:\Omega_X\ra Y, \qq F(x_0)=y_0\in \Omega_Y\subset Y \qq \text{offen}
    \]
    \[
        G:\Omega_Y\ra Z
    \]
    Falls $F'(x_0)$ und $G'(y_0)$ existiert, so ist 
    \[
    (G\circ F)'(x_0)=G'(y_0)\circ F'(x_0)
    \]
\end{remark}

\begin{defi}[Partielle Ableitung]
    Seien $X,$ $Y$, $Z$ Banachräume, $\Omega_X\subset X$ offen, $x_0\in \Omega_X$, $\Omega_Y\subset Y$ 
    offen, $y_0\in \Omega_Y$. Der Operator
    \[
        F: \Omega_X \times \Omega_Y \ra Z 
    \]
    heißt partiell in $(x_0,y_0)$ nach dem zweiten Argument (nach $y$) differenzierbar, falls die 
    Ableitung 
    \[
        F(x_0, \cdot) : \Omega_Y \ra Z 
    \]
    differenziebar ist. Wir nennen den linearen Operator $F_Y(x_0,y_0): Y\ra Z$, der
    \[
        \lim_{\|h\|\ra 0} \frac{\| F(x_0,y_0+h) - F(x_0,y_0) - F_Y( x_0,y_0)h \|}{\|h\|} = 0 
    \]
    erfüllt, die partielle Ableitung von $F$ in $(x_0,y_0)$ nach dem zweiten Argument.
\end{defi}

\begin{prop}
    Seien $X,Y$ Banachräume, $\Omega \subset X $ offen und konvex mit $x_0, x_1\in \Omega$. 
    $F: \Omega \ra Y$ sei stetig Fréchet-differenzierbar auf $\Omega$. Dann gilt
    \[
        F(x_1)-F(x_0)=\int_{0}^1 F'(x_0+t(x_1-x_0))(x_1-x_0)\d t
    \]
    Das Integral ist als Limes der entsprechenden Riemannsumme zu verstehen und dieser existiert.
\end{prop}

\begin{proof}
    Übungsaufgabe \[\]
\end{proof}

Ähnlich dem endlichdimensionalen Fall geben uns die Ableitungen im Banachraum
hinreichende Bedingungen um Operatoren implizit zu definieren. Die Fragestellung ist die folgende:
Seien $X,\, Y,\, Z$ Banachräume, $U$ eine Umgebung von $x_0\in X$. $V$ eine Umgebung von $y_0\in Y$.
Wir suchen zu $F:U\times V \ra Z$ einen Operator
\[
    T:U_0\subset U \ra V, 
\]
sodass gilt
\[
    F(x,Tx)=F(x_0,y_0) \qq \forall \, x\in U_0.
\]
Durch eine einfache Verschiebung ist es ausreichend, den Fall
\[
    F(x_0,y_0)=0
\]
zu untersuchen.

\begin{prop}
    Sei $X$ ein Banachraum, $\Id: X\ra X$, $x\mapsto x$ die Identität auf $X$. Es sei
    \[
        R:B_r(0)\subset X \ra X, 
    \]
    eine $k$-Kontraktion mit $k<1$, d.h. $\|R(x)-R(y)\|\leq k \|x-y\|$, und es gelte
    \[
        \|R(0)\| < r(1-k)
    \]
    Dann existiert genau ein $x\in B_r(0)$ mit
    \[
        (\Id + R)x = 0
    \]
\end{prop}

\begin{proof}
    Sei $S=-R$, wir suchen also einen Fixpunkt von $S$.
    \begin{description}
        \item{1. Eindeutigkeit:} 
        Seien $Sx=x$ und $Sx'=x'$, damit gilt
        \[
            \|x-x'\|=\|Sx-Sx'\| \leq k \|x-x'\|.
        \]
        Mit $k<1$ folgt $x=x'$.
        
        \item{2. Existenz:}
        Sei $x\in B_r(0)$, es gilt
        \[
            \|Sx\|\leq \|Sx-S(0)\|+\|S(0)\|\leq k\|x\|+\|S(0)\|< kr + r(1-k) = r
        \]
        Sei $x_p=Sx_{p-1}$, $x_0=0$. Es gilt (wie auch im Banach'schen Fixpunktsatz, siehe ÜB 1), dass
        \[
            \|x_{n+p}-x_n\| \leq k^n(1-k)^{-1}\|x\|,
        \]
        damit ist $(x_p)_{p\in \N}$ eine Cauchy-Folge und konvergiert gegen $x\in X$.
        Wir haben weiter, dass
        \[
            \|x\| \leq \underbrace{\|x-x_{p+1}\|}_{\ra 0} + \|x_{p+1}\| \qq \text{mit} 
            \qq \|x_{p+1}\|< (1-k) \|S(0)\|=r\\
            \Ra \, \|x\|<r
        \]
        Wegen $x_{p+1}=Sx_p$ gilt dass $x=Sx$, somit ist $x$ der gesuchte Fixpunkt. 
    \end{description}
    \[  \]
\end{proof}
   % 1.2 Vorbereitung \\ 1.2.1 Ableitung in Banachräume und implizite Funktion
\begin{theorem}[Satz über implizite Funktion]\label{1.12}
    Seien $X$, $Y$, $Z$ Banachräume, $U\subset X$ Umgebung von $x_0\in X$, $V\in Y$ Umgebung von 
    $y_0 \in Y$. Sei weiter
    \[F:U\times V\ra Z\]
    stetig und stetig differenzierbar nach der zweiten Variablen. $F_Y(x_0,y_0)$ sei eine Bijektion
    von $Y$ nach $Z$ und es gelte
    \[
        F(x_0,y_0)=0
    \]
    Dann existiert $B_\delta(x_0)\subset U$, $B_r(y_0)\subset V$ und genau ein Operator 
    $T:B_\delta(x_0)\ra B_r(y_0)$, so dass $T(x_0)=y_0$ und $F(x,Tx)=0$ $\forall \, x\in B_\delta(x_0)$.
    $T$ ist stetig.
\end{theorem}

\begin{proof}
    Ohne Einschränkung sei $x_0=y_0=0$. Sei $L:= F_Y(0,0)$, $\Id: Y\ra Y$ die Identität auf $Y$. Es 
    sei $S(x,y):= L^{-1}F(x,y)-y$ für $(x,y)\in U\times V$. Somit gilt
    \[
        F(x,y)=0 \qq \LRa \qq y+S(x,y)=0.
    \]
    $S$ ist stetig differenzierbar nach dem zweiten Argument mit
    \[
        S_Y=L^{-1} F_Y(x,y)-\Id
    \]
    Damit gilt
    \[
        S_Y(0,0)=0
    \]
    Sei $k\in (0,1)$. Wegen Stetigkeit von $S_Y$ existiert $r>0$ mit 
    \[
        \|S_Y(x,y)\|\leq k \qq \forall (x,y)\in B_r(0)\times B_r(0)
    \]
    Sei nun $x\in B_r(0)$, $y, \tilde y \in B_r(0)$
    Es gilt nach Proposition \ref{1.10}, dass
    \[
        \|S(x,y)-S(x,\tilde y)\|=\left\|\int_0^1 S_Y(x,\tilde y +t(y-\tilde y)) (y-\tilde y) \d t 
        \right\|\leq k\cdot\|y-\tilde y\|
    \]
    Wegen $S(0,0)$ und Stetigkeit von $S$ existiert $\delta \leq r$, so dass
    \[
        \|S(x,0)\|\leq r(1-k) \qq \forall \, x \in B_\delta (0)
    \]
    Sei also $x\in B_\delta(0)$. Nach Proposition \ref{1.11} existiert genau ein $y\in B_r(0)$ mit
    $y+S(x,y)=0$. Wir setzen
    \[
        Tx=y, \qq T:B_\delta(0)\ra B_r(0)
    \]
    Es gilt $T(0)=0$ wegen $0+S(0,0)=T(0)+S(0,T(0))=0$ und der Eindeutigkeit von $T$. Es bleibt die
    Stetigkeit von $T$ zu zeigen: Seien $x,x'\in B_\delta (0)$, damit gilt
    \[
        0=Tx+S(x,Tx)=T(x'+S(x,Tx'))
    \]
    also
    \begin{align*}
        \|Tx-Tx'\|&\leq \|S(x',Tx')-S(x,Tx')\|+\|S(x,Tx)-S(x,Tx)\|\\
            &\leq \|S(x',Tx') - S(x,Tx')\|+k\|Tx-Tx'\|\\
            &= (1-k)\|Tx-Tx'\|\\
            &\leq \|S(x',Tx')-S(x,Tx')\| \ra 0 \qq \text{für $x\ra x'$} 
    \end{align*}
    Somit ist $T$ stetig.
    \[\]
\end{proof}

\begin{remark}
    Ist $F$ $r$-mal stetig differenziebrar, so gilt das auch für $T$.
\end{remark}

\begin{proof}
    Übungsaufgabe \[ \]
\end{proof}

\begin{theorem}\label{1.13}
    Seien $X$, $Y$ Banachräume, $U\subset X$ eine Umgebung von $x_0$. Es sei $F:U\ra Y$ stetig 
    differenzierbar und $F'(x_0)$ sei eine lineare Bijektion von $X$ nach $Y$.
    Dann existiert eine Umgebung $U_0\subset U$ von $x_0$, so dass
    \[
        F|_{U_0}: U_0\ra F(U_0)\ni y_0 -F(x_0)
    \]
    ein Homöomorphismus (bistetige Abbildung) ist.
\end{theorem}

\begin{proof}
    Wir wenden Satz $\ref{1.12}$ auf
    \[
        \tilde F(x,y):= F(x)-y
    \]
    an.
    \[  \]
\end{proof}

\begin{remark}
    Ist $F$ $r$-mal stetig differenzierbar, so gilt das auch für $F^{-1}$ ($F$ ist ein 
    $r$-Diffeomorphismus).
\end{remark}

\begin{proof}
    Übungsaufgabe\[ \]
\end{proof}

\begin{defi}[Zusammenhände Mengen]\label{1.14}
    \begin{description}
    \item{-}
    Sei $X$ ein (topologischer metrischer, normierter) Raum. Eine Menge $\Omega \subset X$ heißt
    zusammenhängend, falls es keine zwei abgeschlossenen (offenen) $\Omega_1$, $\Omega_2$ gibt
    mit
    \[
        \Omega\subset\Omega_1\cup\Omega_2, \qq \Omega\cap \Omega_1 \cap \Omega_2 = \varnothing, \qq
        \Omega\cap\Omega_{1,2}\neq \varnothing
    \]
    \item{-}
    Eine Menge $\Omega\subset X$ heißt wegzusammenhängend, falls sich je zwei Punkte in $\Omega$ durch
    eine stetige, in $\Omega$ verlaufende Kurve verbinden lassen.
    \item{-}
    Eine Menge $\ol\Omega \subset \Omega$ heißt Zusammenhangskomponente von $\Omega$, falls
    $\ol \Omega \subset \Omega$ maximal, zusammenhängend.
    \end{description}
\end{defi}

\begin{remark}
    Wegzusammenhängend $\Ra$ Zusammenhängend.

    \noindent Offen, zusammenhängend $\Ra$ Wegzusamenhängend
\end{remark}

\begin{theorem}[Mittelwert]\label{1.15}
    Seien $X$, $Y$ Banachräume, $F:X\ra Y$ stetig differenzierbar.
    \begin{description}
    \item{i)}
    Falls $\Omega$ konvex ist, so gilt
    \[
        \|F(x)-F(y)\|\leq M\|x-y\|,
    \]
    wobei
    \[
        M=\max_{0\leq t \leq 1}\| F'((1-t)x + ty) \|
    \]
    \item{ii)}
    Umgekehrt gilt: Falls
    \[
        \| F(x)-F(y) \|\leq M\|x-y\| \qq \forall \, x,y\in \Omega
    \]
    Dann gilt
    \[
        \sup_{x\in \Omega}\|F'(x)\|\leq M
    \]
    \end{description}
\end{theorem}

\begin{proof}
    Sei $f(t):= F((1-t)x+ty)$, $0\leq t\leq1$. Nach Kettenregel gilt
    \[
        f'(t)=F'((1-t)x+ty)(x-y)  
    \]
    \[
        \Ra \, \|f'(x)\|\leq\tilde M:= M\|x-y\|  
    \]
    \begin{description}
    \item{i)}
    Sei $\phi(t):=\|f(t)\|$ für $\delta>0$. Wir wollen zeigen, dass $\phi(t)\leq0$ $\forall \, 
    \delta>0$, $0\leq t\leq1$. Sei also (zum Widerspruch)
    \[
        t_0:=\max\{t\in [0,1]\, | \, \phi(s)\leq 0 \, \forall \, s\leq t\}.
    \]
    Dann gilt
    \begin{align*}
        \phi(t_0+\eps)&=\|f(t_0+\eps)-f(t_0)+f(t_0)-f(0)\|-(\tilde M + \delta)t\\
        &\leq \|f(t_0+\eps)-f(t_0)\|-(\tilde M +\delta) -\phi(t_0)\\
        &\leq \|f'(t_0)\eps+ \mc o (1)\|-(\tilde M +\delta)\eps\\
        &\leq (-\delta + \mc o (1))\eps 
    \end{align*}
    \item{ii)}
    Angenommen, es existiert $x_0$ mit $\|F'(x_0)\|\geq M+2\delta$, $\delta>0$. Dann existiert
    $e\in X$, $\|e\|=1$, $\| F'(x_0)e \|\geq M+\delta$. Somit gilt
    \begin{align*}
        M\eps&\geq \|F(x_0+\eps e) - F(x_0)\|= \|F'(x_0)(\eps e)+\mc o (\eps)\|\\
         &\geq (M+\delta)\eps- \mc o(\eps) > M\eps
    \end{align*}
    Das ist ein Widerspruch.
    \end{description}
    \[\]
\end{proof}

\begin{cor}\label{1.16}
    Sei $\Omega \subset X$ offen, (weg-)zusammenhängend, $F$ stetig differenzierbar auf $\Omega$.
    Es gilt
    \[
        F= \Const \qq \LRa \qq F'=0
    \]
\end{cor}
\begin{remark}
    Wir schreiben wie im endl. dim. $C(\Omega)= C^0(\Omega), \, C^1(\Omega)\, …$
\end{remark}

\noindent \textbf{Anwendungen}: Lokale Existenz und Eindeutigkeit Banachraum-wertiger
Differenztialgleichunegen. Sei $X$ Banachraum, $\Omega\subset X$ offen, $I\subset \R$
kompaktes Intervall. Es sei $C_b(I,\Omega)$ der Banachraum der beschränkten, stetigen Abbildungen
von $I$ nach $\Omega$, versehen mit der $\sup$-Norm.

\begin{lem}
    Sei $Y$ ein Banachraum, $f\in C(\Omega,Y)$ und sei die Funktion
    \[
        f_\star: C_b(I,\Omega)\ra C_b(I,Y)
    \]
    definiert als
    \[
        (f_\star x)(t)= f(x(t))
    \]
    Es gilt $f_\star\in C^r$
\end{lem}

\begin{proof}
    \begin{description}
    \item{$r=0$:}
    Sei $x_0\in C_b(I,\Omega)$, $\eps>0$ $\forall \, t \in I$ existiert $\delta(t)>0$, so dass
    \[
        \|f(\chi)-f(x_0(t))\|\leq \frac{\eps}{2} \qq \forall \, \chi\in U
    \]
    Die offenen Kugeln 
    \[
        \{ B_{\delta(t)}(x_0(t)) \}_{t\in I}
    \]
    sind eine offene Überdeckung vom $\{ x_0(t) \}_{t\in I}$. Diese Menge ist als stetiges Bild einer
    kompakten Menge kompakt, und somit existiert endliche Teilüberdeckung
    \[
        \{ B_{\delta(t_j)}(x_0)(t_j) \}_{1\leq j \leq N}
    \]
    Sei nun $x\in C_b(I,\Omega)$ mit
    \[
        \|x-x_0\|\leq \delta := \min_{1\leq j \leq N}\delta(t_j)
    \]
    Somit existiert $\forall \, t\in I$ ein $t_j$, so dass $\|x_0(t)-x_0(t_j)\|\leq \delta(t_j)$,
    und deshalb gilt
    \[
        \|f(x(t))- f(x_0(t))\|\leq \|\underbrace{f(x(t))-f(x_0(t_j))}_{\leq 2\delta}\|
        +\|\underbrace{f(x_0(t_j))-f(x_0(t))}_{\leq\delta}\|,
    \]
    denn
    \[
        \|x(t)-x_0(t_j)\|\leq \|x(t)-x_0(t)\|+\|x_0(t)-x_0(t_j)\|\leq 2\delta(t_j).
    \]
    Somit folgt die Steigkeit……
    \item{$r=1$:}
    Wir müssen zeigen, dass
    \[
        \sup_{t\in I} \|f(x_0(t))+x(t)- f(x_0(t)) - f'(x_0(t))x(t)\|\leq \eps \sup_{t\in I}\|x(t)\|
    \]
    denn
    \[
        (f'_\star(x_0)x)(t)= f'(x_0(t))x(t)
    \]
    Übungsaufgabe. Folgt wie Stetigkeit durch Kompaktheit von $I$.
    \end{description}
    \[  \]
\end{proof}
   %
\section{Der Brouwer'sche Abbildungsgrad}

\subsection*{Motivation}
\begin{description}
    \item[Ziel:]
    $f(x)=0$ zu lösen für $f:U\subset X \ra X$, $X$ Banachraum.
    \item[Frage:]
    Existenz/Anzahl der Lösungen
    \item[Rückblick auf Funktionentheorie:]
    Sei $z_0\in \C$
    \[
        n(\gamma,z_0) = \frac{1}{2\pi i} \int_\gamma \frac1{z-z_0} \d z
    \]
\end{description}
$\leadsto$ Verallg.: $f\in \mc H(\C)$. $0\nin f(\gamma)$
\[
    n(f(\gamma),0)=\frac1{2\pi i} \int_\gamma \frac{f'}{f} \d z = \sum_k n(\gamma,z_k) \alpha_k
\]
wobei $f(z_n)=0$, $\alpha_k$ Vielfachheiten.

\noindent Ziel: Verallg. des Begriffs „Umlaufzahl“ für Abb. $f:U\subset \R^n\ra \R^n$

\subsection{Die Determinantenformel}

\subsubsection{Notation}

$U\Subset \R^n$ offen, $f\in C^1(U,\R^n)$
\begin{align*}
    J_f(x)&:=\det \d f(x)\\
    \RV(f)&:=\{ y\in \R^n \,| \, \forall \, x\in f^{-1}(y), \, J_f(x)\neq 0 \}\\
    \CV(f)&:= \R^n \setminus \RV(f)\\
    D_y^r(\ol U, \R^n)&:=\{f\in C^k(\ol U, \R^n)\,|\,y\nin f(\partial U)\}\\
    D_y(\ol U, \R^n)&:=D_y^0(\ol U, \R^n)
\end{align*}
$\tau(\R^n)$ bezeichne die Topologie auf $\R^n$.
\begin{defi}
    Eine Abbildung
    \[
        \deg: \bigcup_{U\in \tau(\R^n),\, y\in \R^n} (D_y(\ol U, \R^n)\times \{U\}\times \{y\}\ra \R,
    \]
    d.h.
    \[
        \deg=\deg(f,U,y) 
    \]
    heißt Gradabbildung, falls
    \begin{description}
        \item[(D1)]
        Translationsinvarianz: 
        \[
            \deg(f,U,y)=\deg(f-y,U,0)
        \]
        \item[(D2)]
        Normalisation:
        \[
            \deg(\Id,U,y)=1 \qq \forall y\in U
        \]
        \item[(D3)]
        Additivität: Seien $U_1,U_2\subset U$ offen und disjunkt, sodass
        \[
            y\nin f(\ol U\setminus(U_1\cup U_2)),
        \]
        dann gelte
        \[
            \deg(f,U,y)=\deg(f,U_1,y)+\deg(f,U_2,y)
        \]
        \item[(D4)]
        Homotopieinvarianz: 
        \[
            H(t)=(1-t)f+tg \in D_y(\ol U,\R^n) \qq \forall \, t\in [0,1]\, \Ra \, \deg(f,U,y)=\deg(g,U,y)
        \]
\end{description}
\end{defi}

\begin{theorem}\label{2.1}
    Sei $\deg$ eine Gradabbildung. Dann gilt
    \begin{description}
        \item{(i)}
        $\deg(f,\varnothing,y)=0$ und
        \[
            \deg(f,U,y)=\sum_{i=1}^N \deg(f,U_i,y)
        \]
        falls 
        \[
            y\nin f(\ol U\setminus\bigcup_{i=1}^NU_i),
        \]
        wobei $U_i\subset U$ offen und disjunkt für $1\leq i \leq N$.
        \item{(ii)}
        $y\nin f(U) \, \Ra \, \deg(f,U,y)=0$
        \item{(iii)}
        $|f(x)-g(x)|<\dist(y,f(\partial U)) \, \forall \, x\in \partial U \, \Ra \, \deg(f,U,y)
        =\deg(g,U,y)$
    \end{description}
\end{theorem}

\begin{proof}
    \begin{description}
    \item{(i)}
    Sei $U_1=U$, $U_2=\varnothing$, einsetzen in \textbf{(D3)}
    \[
        \Ra \, \deg(f,\varnothing,y)=0
    \]
    $i=1$: $U_2=\varnothing$
    \[
        \Ra \, \deg(f,U,y)=\deg(f,U_1,y)
    \]
    $i>1:$ Induktion mittels \textbf{(D3)}
    \item{(ii)}
    \begin{align*}
        y\nin f(U) \Ra y &\nin f(\ol U)\, \Ra \, y\nin f(\ol U\setminus \varnothing)\\
        \overset{(i)}{\Ra } \, \deg (f,U,y)&=0 \qq (i=1, U_1=\varnothing)
    \end{align*}
    \item{(iii)}
    Sei $H(t,x):=(1-t)f(x)+tg(x)$ und sei $x\in \partial U$
    \begin{align*}
        \Ra \, |H(t,x)-y|&= |f(x)-y+t(g(x)-f(x))|\\
            &\geq |f(x)-y| - |g(x)-f(x)|\\
            &\geq \dist(y,f(\partial U))- |g(x)-f(x)| >0\\
        \Ra \, y&\nin H(t,\partial U) \qq \forall \, t\,\Ra \,H(t)\in D_y(\ol U,\R^n)\\
        &\overset{(\textbf{D4})}{\Ra} \text{ Behauptung.}
    \end{align*}
    \end{description}
    \[ \]
\end{proof}

\begin{theorem}\label{2.2}
    \begin{description}
        \item{i)} $\deg(\cdot,U,y)$ ist lokal konstant in $D_y(\ol U,\R^n)$
        \item{ii)} $\deg(f,U,\cdot)$ ist lokal konstant in $\R^n\setminus f(\partial U)$
        \item{iii)} Seien $H:[0,1]\times \ol U\ra \R^n$ und $y:[0,1]\ra \R^n$ stetig (d.h. $H$ ist eine
            Homotopie zwischen $H(0)=H(0,\cdot)$ und $H(1)$), so gilt
        \[
            \deg(H(0),U,y(0))=\deg(H(1),U,y(1)),
        \]
        falls $H(t)\in D_{y(t)}(\ol U, \R^n)$ $\forall t \in [0,1]$
    \end{description}
\end{theorem}

\begin{proof}
    Beachte: $D_y(\ol U, \R^n)$ ist offen in $C^0(\ol U, \R^n)$
    \begin{description}
    \item{i)}
        $\|f-g\|_{C^0,\ol U} <\eps \, \Ra \, |f(x)-g(x)|<\eps \, \forall \, x\in\partial U$ mit
        \[
            \eps:= \dist(y,f(\partial U)) \, \Ra \, \deg (f,U,y)=\deg (g,U,y)
        \]
    \item{ii)}
        Sei $y_0\nin f(\partial U)$ und $y\in B_{\dist(y_0,f(\partial U))}(y_0)\subset \R^n\setminus
                f(\partial U)$
        \begin{align*}
            \Ra \, \|(f-y)-f\|&<\dist(y_0,f(\partial U))\\
                \overset{i)}{\Ra} \, \deg(f-y,U,y_0) &= \deg(f,U,y_0)\\
            \overset{\textbf{(D1)}}{\Ra } \, \deg (f,U,y_0+y)&= \deg(f,U,y_0)
        \end{align*}
    \item{iii)}
        $H$ ist gleichmäßig stetig
        \[
            \Ra \, H:[0,1]\ra C^0(\ol U,\R^n)\qq t\mapsto H(t,\cdot)
        \]
        ist auch stetig. $H$ ist ein stetiger Weg in $D_y(\ol U,\R^n)$. Sei $y$ fest
        \[
            \Ra \, \deg(H(t,\cdot),U,y)=\Const,
        \]
        weil $\deg(\cdot,U,y)$ konst. auf Zusammenhangskomponenten ist. Für $y=y(t):$
        \begin{align*}
            \deg(H(0),U,y(0))&=\deg(H(0)-y(0),U,0)=\deg(H(t)-y(t), U, 0) \qq \forall t\in [0,1]\\
            &= \deg(H(1),U,y(1)).
        \end{align*}
    \end{description}
    \[ \]
\end{proof}

\begin{lem}\label{1.3.5}
    Zwei nichtsinguläre Matrizen $M_1,M_2\in \Gl(n)$ sind genau dann homotop in $\Gl(n)$, falls
    \[
        \sign \det M_1 = \sign\det M_2
    \]
\end{lem}

\begin{proof}
    \begin{description}
    \item{„$\Ra$“} Sei $M\in \Gl(n)$. Wegen der Linearität von $\det $ in Zeilen können elementare
    Zeilenumformungen mit Hilfe stetiger Deformationen in $\Gl(n)$ erzeugt werden.
    \[
        M \, \leadsto\diag(m_1,…,m_m), \qq \text{mit} \qq |m_i|=1
    \]
    weil $\sign \det M_1=\sign \det M_2$.
    \[
        H(t):= \begin{pmatrix} \pm \cos (\pi t)& \mp \sin(\pi t)\\ \sin(\pi t) & \cos(\pi t)\end{pmatrix}
    \]
    ist eine Homotopie in $\Gl(n)$ von $\diag(\pm 1,1)$ nach $\diag(\mp 1,-1)$. $i=n$ transformiere
    $\diag(\mp 1,-1)$.

    \noindent $i=n$ transformiere $\diag(m_i,…,m_i)$ nach $\diag(\pm1,1)$

    \[
        \leadsto \begin{pmatrix} \sign\det M& 0& \cdots&0\\ 0& 1 &\cdots&0\\\vdots&&\ddots&0\\ 0&\cdots&
        \cdots& 1\end{pmatrix}
    \]

    \end{description}
\end{proof}

\begin{theorem}
    Sei $f\in D_y^1(\ol U, \R^n)$, $y\nin \CV(f)$ und $\deg$ eine Gradabbildung. Dann gilt
    \[
        \deg(f,U,y)=\sum_{x\in f^{-1}(y)} \sign J_f(x),
    \]
    wobei die Summe endlich ist.
\end{theorem}

\begin{proof}
    O.B.d.A. $y=0$ (nach \textbf{D1}). Alle $x\in f^{-1} (0)$ sind isolierte Punkte in $U$
    (Homöomorphiesatz).
    $f^{-1}(0)$ hat höchstens am Rand einen Häufungspunkt, aber $0\nin f(\partial U)$.
    \[
        \Ra \, f^{-1}(0)= \{ x^i \}_{i=1}^N
    \]
    Wähle $\delta >0$ so klein, dass $B_\delta(x^i)$ paarw. disjunkt.
    \[
        \overset{\text{\ref{2.1} i)}}{\Ra} \, \deg(f,U,0)=\sum_{i=1}^N \deg(f,B_\delta(x^i),0).
    \]
    Beachte $0\nin f(\ol U \setminus \bigcup_{i=1}^N B_\delta(x^i))$.
    \begin{align*}
        f\in C^1 \, \Ra \, f(x)&= \d f(x)(x-x^i) + |x-x^i| r(x-x^i) \qq \text{mit} \qq r\in C^0(B_\delta(x^i),\R^n), 
        \qq r(0)=0.\\
           H(t,x) &:= \d f (x^i) (x-x^i)+ (1-t) (x-x^i)+(x-x^i)
    \end{align*}
    Zeige $0\nin H(t,\partial B_\delta(x^i))$, für alle $t\in [0,1]$.
    \[
        J_f(x^i)\neq 0 \, \Ra \, \exists\, \lambda >0:\, |\d f(x^i)(x-x^i)|\geq \lambda|x-x^i|
    \]
    O.B.d.A. sei $\delta$ so klein, dass $|r(x-x^i)|<\lambda$ in $B_\delta (x^i)$.
    \begin{align*}
        \Ra\, | H(t,x) |&>|\d f(x^i)(x-x^i)|-(1-t)(x-x^i)r(x-x^i)\geq \lambda \delta - \delta|r|>0
        \qq \forall \, x\in \partial B_\delta (x^i)\\
        \overset{(\textbf{D4})}{\Ra} \, \deg(f,U,0)&=\sum_{i=1}^N\deg (\d f (x^i)(\cdot-x^i),
                B_\delta(x^i),0)
    \end{align*}
    Lemma \ref{1.3.5}
    \[
        \Ra \, \deg(\d f (x^i)(\cdot- x^i),B_\delta (x^i),0)= \deg(\diag(\sign J_f(x^i),1,…,1),
            B_\delta(x^i),0)
    \]
    Falls $J_f(x^i)>0\, \overset{\textbf{(D2)}}{\Ra}\,\deg(I(\cdot-x^i),B_\delta (x^i),0)=1$. Es genüngt
    also, $\deg(M(\cdot-x^i),B_1(x^i),0)$. Zu berechnen, wobei $M=\diag(-1,1,…,1)$ ist.
    \begin{align*}
        U_1&:= B_1 (x^i) = \{ \max_{1\leq k\leq n} |x_k-x_k^i|<1 \}\\
        U_2&:= U_1+ (2,0,\cdot,0)\\
        g(r)&=2-|r-1|, \, h(r)=1-r^2\\
        f_1(x)&:= (1-g(x_1-x_1^i)h(x_2-x_2^(i))…h(x_n-x_n^{(i)}),…, 1)\\
        f_2(x)&:= (1,x_2-x_2^{(i)}, … , x_n-x_n^{(i)})\\
        f_1^{-1}(0)&= \{ y,z \} \qq y=x^i, z = x^i+(2,0,…,0)\\
        f_1|_{\partial U}&= f_2|_{\partial U} \qq \Ra \qq \deg(f_1,U_2,0)=0\\
        \Ra \, \deg (f_1,U,0)&= \deg(f_1,U_1,0)+ \deg(f_2,U_2,0) \qq (\star)\\
        \Ra \, \deg(M,B_1(x^i),0)&=\deg(\partial f_1(y),B_1(x^i), 0)\\
        &=\deg(f_1, U_1,0) \overset{(\star)}{=} -\deg (f_1,U_2,0)\\
        &= -\deg(\partial f_1(z),U_2,0)=\deg(\Id, U_2,0)=-1
    \end{align*}
    \[ \]
\end{proof}
   % 2 Der Brouwer'scher Abbildungsgrad \\ 2.1 Die Determinantenformel
\subsection{Verallgemeinerung der Determinantenformel}

Wir haben einen Kandidaten für $\deg$ identifiziert; für $f\in D'_Y(\ol\Omega,\R^n)$, $y\in \R^n\cap
RV(f)$ gilt
\[
    \deg(f,\Omega, g)=\sum_{x_j\in f^{-1}(y)} \sign J_f(x_j).
\]

\noindent \textbf{Problem:} Nur für glatte $f$, und reguläre Punkte zu definieren. Programm zum 
Existenzbeweis: Verallgemeinerung der Determinanten-Formel auf
\begin{description}
    \item{-} kirtische Werte $y$.
    \item{-} nur stetige Funktionen $f$.
\end{description}

\noindent Drei Beispiele zur Illustration:
\begin{enumerate}
    \item $f(x)=x^2$, $U=(-1,1)$
    \item $f(x)=x^2$, $U=(-1,2)$
    \item $f(x)=x+2\sin(x)$, $U=(-10,10)$
\end{enumerate}
Wie aus dem Beispielen ersichtlich ist, haben wir noch (kleine) Schwierigkeiten, den Abbildungsgrad
an kritischen Werten von $f$ zu definieren. Stetige Fortsetzungen liegt aber nahe. Das funktioniert
aber nur falls es von diesen nicht „zu viele“ gibt.

\subsubsection*{Schritt 1:} Kritische Werte von $f$:

\begin{lem}[Sard]\label{2.5}

Es sei $f\in C^1(\Omega, \R^n)$, $\Omega\subset\R^n$ ist offen und beschränkt. Dann ist
$CV(f)$ eine Lebesgue-Nullmenge.
\end{lem}
    
\begin{proof}
    Klar, falls $f$ eine affine Abbildung ist. (Dimension des Bildraumes der (konstanten) Ableitung!)
    Wir linearisieren. 

    \noindent Es sei $\CP(f):=\{x\in \Omega\,|\, J_f(x)=0\}$ die Menge der kritischen Punkte von $f$.
    Sei $\{ Q_j \}_{j\in \N}$ eine abzählbare, offen Überdeckung von $\Omega$ bestehend aus Würfeln,
    so dass
    \[
        \ol \Omega_i\subset \Omega, \, i\in \N.
    \]
    Es gilt 
    \[
        CV(f)=f(\CP(f))=\bigcup_{j\in \N} f(\CP(f)\cap Q_j)
    \]
    Es reicht also zu zeigen, dass
    \[
    |f(\CP(f)\cap Q_j)|
    \]
    für alle $j\in \N$ verschwindet. Sei nun $Q$ ein solcher Würfel, $\rho$ dessen Kantenlänge.
    Sei $\eps >0$, und sei $Q$ unterteilt in $N^n$ Würfel $Q^i$ der Seitenlänge $\frac{\rho}{N}$, so
    dass
    \begin{align}\label{4}
        |f(x)-f(\tilde x)-f'(\tilde x)(x-\tilde x)|&\leq \int_0^1|f'(\tilde x+t(x-\tilde x))-
        f'(\tilde x)|\cdot|x-\tilde x|\d x\leq \frac{\eps\rho}{N} \qq \text{für alle} 
            \qq x,\tilde x \in Q^i
    \end{align}
    So ein $N$ existiert, nachdem $f'$ auf $Q$ gleichmäßig stetig ist. Nun enthalte $Q^i$ einen
    kirtischen Punkt $\tilde x_i \in \CP(f)$, ohne Einschränkung sei $\tilde x_i=0$, $f(\tilde x_i)=0$,
    und $M=f'(\tilde x _i)$. Wegen $\det M =0$ existiert eine ONB $\{b^j\}_{j=1}^n$ mit
    \[
    b^n \bot R(M)
    \]
    Weiter gilt 
    \begin{align*}
        Q^i&\subset \left\{ \sum_{j=1}^n \lambda _j b^j \, \big| \, \| \lambda \|_2\leq \sqrt{n} 
            \frac{\rho}{N} \right\}\subset\left\{\sum_{j=1}^n\lambda _jb^j \, \big|
            \, |\lambda_j|\leq \sqrt{n}\frac{\rho}{M} \, \forall 1\leq j\leq n\right\}
    \end{align*}
    Damit existiert $C>0$, (unabhängig von $i$), so dass
    \[
        MQ^i\subset \left\{\sum_{j=1}^{n-1}\lambda_jb^j\, \big| \, |\lambda_j|
            \leq C\cdot \frac{\rho}{N} \right\}
    \]
    mit $C=\sqrt{n} \max_{x\in \ol Q}|f'(x)| $. Damit gilt nach (\ref{4}) sogar, dass
    \[
        f(Q^i)\subset\left\{ \sum_{j=1}^n \lambda_jb^j\, \big| \, |\lambda_j|\leq 
        (C+\eps)\frac{\rho}{N} \, \forall 1\leq j\leq n-1, \, |\lambda_n|\leq \frac{\eps\rho}{N}\right\}
    \]
    Es folgt
    \[
        |f(Q^i)|\leq \frac{\tilde C\eps}{N^n},
    \]
    falls in $Q^i$ ein kritischer Punkt liegt. Es gibt maximal $N^n$ Unterwürfel $Q^i$ mit kritischen
    Punkten. Somit gilt
    \[
        |f(Q\cap \CP(f))|\leq C\cdot \eps
    \]

\end{proof}

\noindent Dank Lemma \ref{2.5} ist $\R^n\setminus \CV(f)$ dicht in $\R^n$. Das reicht leider (?)
noch nicht.

\noindent Wir brauchen $d_1=d_2$, um den Abbildungsgrad sinnvoll durch die Determinanten-Formel
definieren zu können. (denn $\deg$ soll konstant sein auf Zusammenhangskomponenten, unabhängig
von krit. Werten.)

\textbf{Idee:} Umschreiben der Determinantenformel als \textit{Integral}.
Es sei im Weiteren $\eta_\eps$ ein Standard-Mollifier (Standard-Dirarcfolge), d. h.
\[
    \eta_\eps \in C^\infty_c(B_\eps(0)\subset \R^n), \qq \int_{\R^n} \eta_\eps=1,\qq \eta_\eps\geq 0
\]

\begin{lem}\label{2.6}
    Sei $f\in D^1_y(\ol \Omega, \R^n)$, $y\nin \CV(F)$. Dann gilt
    \begin{align}\label{5}
        \deg(f,\Omega,y)&=\sum_{x_j\in f^{-1}(y)} \sign J_f(x_j)\\
                        &= \int_\Omega \eta_\eps (f(x)-y) J_f(x)\d x
    \end{align}
    für alle $\eps$ hinreichend klein, d. h.
    \[
        \eps_0>\eps>0, \qq \text{mit} \qq \eps_0=\eps_0(f,y)
    \]
    Es gilt $\supp (\eta_\eps(f(\cdot)-y))\subset \Omega$ für $\eps < \dist(y,f(\partial \Omega))$
\end{lem}

\begin{proof}
    \begin{description}
        \item{1)} 
        Falls $f^{-1}(y)=\varnothing$, dann sei $\eps_0=\dist(y,f(\partial \Omega))$
        \item{2)}
        Falls $f^{-1}(y):= \{x^i\}_{1\leq i \leq N}$, sei $\eps_0>0$, so dass
        \[
            f^{-1}(B_{\eps_0}(y)) = \bigcup_{i=1}^N U_i
        \]
        mit $U_i\cap U_j=\varnothing$ für $i\neq j$.
        Aus dem \textit{Satz über die imliziete Funktion} (war ja klar, dass wir den mal brauchen)
        folgt für evtl. noch kleiners $\eps_0$, dass
        \[
            f\big|_{U_i(x_i)} \qq \text{bijektiv,} \qq J_f(x)\neq 0 \qq \forall x \in U_i(x_i).
        \]
        Wieder gilt $\eta _{\eps}(f(x)-y)=0$ für
        \[
            x\in \ol \Omega \setminus \bigcap _{i=1}^N U_i(x^i)
        \]
        Damit gilt
        \begin{align*}
            \int_\Omega \eta_\eps(f(x)-y) J_f(x) \d x&= \sum_{i=1}^N \sign J_f(x^i)\cdot
            \int_{B_{\eps_0(x)}}\eta_\eps (\tilde x ) \d \tilde x\\
            &=\sum_{x^i\in f^{-1(y)}} \sign  J_f(x^i)
        \end{align*}
        mit $f(x)=\tilde x$.
    \end{description}
    \[ \]
\end{proof}
Die Integraldarstellung ergibt auch für kritische Punkte Sinn. Aber:
Wegen Abhängigkeit von $\eps_0$ von $f,y$ ist. Auch der Wert des Integrals nicht a priori stetig in 
$f,y$, denn $\eps_0$ hängt ab von $f$ und $y$.

   % 2.2 Verallgemeinerung der Determinantenformel \\ Schritt 1: Kritische Werte von f
Dieses Problem können wir beseitigen, wenn wir $f\in C^2$ fördern. Wir brauchen zunächst einen Hilfssatz.

\begin{prop}\label{2.7}
    \begin{description}
        \item{a)}
        Sei $u\in C^1(\R^n)$, $\supp u \subset K$, K kompakt. Sei $\vp=\div u$. Dann gilt
        \[
            \int_{\R^n} \vp (x) \d x=0
        \]
        \item{b)}
        Sei $\vp \in C^1(\R^n)$, und sei $z\in \R^n$. Dann gilt
        \[
            \vp (x+z)-\vp(x)=\div \left( z\cdot \int_0^1 \vp (z+tz) \d t\right)
        \]
        \item{c)}
        Seien $u$, $K$, $\vp$ wie in $a)$, $\Omega \subset \R^n$ offen, beschränkt.
        Sei $f\in C^2(\ol \Omega)$ mit $K\cap f(\partial \Omega)=\varnothing$. Dann existiert
        $v\in C^1(\ol \Omega)$, mit $\supp v \subset \Omega$, so dass
        \[
            \vp (f(x))J_f(x)=\div v(x) \qq \text{auf} \qq \Omega
        \]
    \end{description}
\end{prop}

\begin{proof}
    \begin{description}
    \item{a)}
    Klar!
    \item{b)}
    Sei
    \[
        \eta (x):= \int_0^1 \vp (x-sz) \d s.
    \]
    \begin{align*}
        (z\cdot \eta(x))&=\frac{\d}{\d t } \eta (x+tz)\big|_{t=0}=\int_0^1 \frac{\d}{\d t} \vp(x+sz+tz)
       \big|_{t=0} \d s\\
        &= \int_0^1 \frac{\d}{\d t} \vp(x+sz)\d s=\vp(x+z)-\vp(x)
    \end{align*}
    \item{c)}
    Es sei $\d _{ij}$ der $(i,j)$- Eintrag der Kofaktormatrix von $(f')_{ij}= \partial _if_j(x)$.
    Es sei $v_j(x)=\sum_{j=1}^n u_j(f(x))\d _{ij}(x),i=1,…,n$. Es gilt $K\cap f(\partial \Omega) = \Phi$,
    $f\in C(\ol \Omega)$, damit existiert $\delta>0$ mit $\dist(K,f(\ol \Omega \setminus 
                f(\Omega_\delta)))>0$.
    Somit gilt $\supp v\subset \Omega_\delta\subset \Omega$.
    Wir rechnen:
    \begin{align*}
        \partial_iv_i&= \sum_{\d _{ij}(x)\partial_ku_j(f(x))}\partial _i f_k.+\sum_{j=1}^nu_j(f(x))
        \partial_i \d_{ij}x
    \end{align*}
    \textit{Behauptung 1:}
    \[
        \sum_{i=1}^n \partial_i\d_{ij}(x)=0
    \]
    \textit{Behauptung 2:}
    \[
        \sum_{i=1}^n \d_{ij} (x) \partial_if_k(x) =\delta_{ij}\cdot J_f(x)
    \]
    Somit gilt:
    \begin{align*}
        \div u(x)&=\sum_{k,j}\partial _k u_j(f(x))\cdot \left( \sum_{i=1}^n \d_{ij}(x) \partial f_k(x) 
                \right) +\sum_{i=1}^nu_j(f(x)) \left( \sum \partial _i \d_{ij}(x) \right)\\
        =&\sum_{k,j}\partial _k u_j(f(x))\delta _{jk} J_f(x)=\vp(f(x))J_f(x)
    \end{align*}
    Behauptungen: Siehe Übungsblatt 2.
    Wir wollen nun zeigen, dass der über die Determinantenformel definierte Abbildungsgrad konstant
    ist auf Zusammenhangskomponente. Zu zeigen: $\deg(f,\Omega,y_1)=\deg(f,\Omega,y_2)$.

    \begin{lem}\label{2.8}
        Sei $f\in C^2(\ol \Omega,\R^n),$ $ y_0\nin f(\partial \Omega) $,
        \[
            \rho:=\dist(y_0,f(\partial \Omega))
        \]
        Dann ist $\deg(f,\Omega,\cdot)$ (Definiert durch die Determinantenformel) konstant auf
        \[
            B_\delta (y_0)\cap \RV(f)
        \]
    \end{lem}

    \begin{proof}
        Sei $y^j\in B_\delta (y_0)\cap \RV(f),$ $j=1,2$, sei $\delta := \rho-\max_{j=1,2}\|y^j-y_0\|$.
        Sei $\eps>0$, so dass
        \[
            \deg(f,\Omega,y^j)=\int \eta_\eps(f(x)-y^j)J_f(x) \d x \qq \text{(nach Lemma $\ref{2.6}$)}
        \]
        mit Proposition $\ref{2.7}$ b) gilt
        \begin{align*}
            \eta_\eps(x-y^2)-\eta_\eps(x-y^1)&= \eta_\eps (x-y^1+(y^1-y^2))-\eta_\eps(x-y^1)\\
            &= \div w(x) \qq\text{mit}\\
            w(x)&= (y^1-y^2)\int_0^1 \eta_\eps (x-y^1+t(y^1-y^2))\d t
        \end{align*}
    \textit{Behauptung:} Es gilt $\supp w\subset B_\delta(y_0)$
    \begin{proof}
        Sei $x\in \supp w$. Damit existert $t\in [0,1]$ mit $|x-y^1+t(y^1-y^2)|<\eps$
        \begin{align*}
            \Ra \, |x-y_0|&\leq \eps + |y^1-t(y^1-y^2)-y_0|\leq \eps +|(1-t)(y^1-y_0)+t(y^2-y_0)|\\
                \eps + \rho-\delta&<\rho
        \end{align*}
        Damit gilt aber, dass $f(\partial \Omega)\cap \supp w=\varnothing$. Mit Proposition 
        $\ref{2.7}$ e) existiert
        \[
            v\in C^1(\Omega),\, \supp v \subset \Omega.
        \]
        und
        \[
            \left( \eta_\eps(f(x)-y^2)-\eta_\eps (f(x)-y^1) \right)J_f(x)=\div v(x)
        \]
        Mit der Proposition $\ref{2.7}$ a) folgt die Behauptung. \[  \]
    \end{proof}

    \end{proof}
    \end{description}
\end{proof}

\begin{defi}\label{2.9}
    Sei $f\in C^2(\ol \Omega,\R^n)$, $y\in f(\partial \Omega)$. Wir setzen
    \[
        \deg (f,\Omega,y)= \deg(f,\Omega,\tilde y)=\sum_{x\in f^{-1}(\tilde y)} \sign J_f(x),
    \]
    wobei $\tilde y \in \RV(f)$ mit
    \[
        |\tilde y- y| < \dist(y,f(\partial \Omega))
    \]
\end{defi}

\noindent\textit{Behauptung:} $\deg$ ist somit nach dem vorherstehenden Überlegungen wohldefiniert.
\subsubsection*{2. Schritt: Nur stetige Funktionen $f$}

Auch hier die Idee: Sei $f\in D_y(\ol \Omega, \R^n)$. Wir suchen $\tilde f \in C^2 (\ol \Omega, \R^n)$
hinreichend nahe an $f$ (in derselben Zusammenhangskomponente von $D_y(\ol \Omega, \R^n)$) und übertragen
den Wert von $\deg(\cdot , \Omega, y)$ von $\tilde f$ auf $f$. Es bleibt zu zeigen, dass die so
definierte Funktion stetig ist.

\begin{lem}\label{2.10}
    Sei $f\in D_y^2(\ol \Omega, \R^n)$, sei $g\in C^2(\ol \Omega, \R^n)$. Dann existiert $\eps >0$, so 
    dass
    \[
        \deg (f+tg,\Omega,y)= \deg(f,\Omega,y) \qq \forall t\in (-\eps,\eps)
    \]
\end{lem}
\begin{proof}
    \begin{description}
        \item{1)}
        $f(y)=\varnothing \, \Ra \, (f+tg)^{-1}(y)=\varnothing$ falls
        \[
            |t| < \frac{\dist(y,f(\ol \Omega))}{\|g\|_\infty}
        \]
        \item{2)}
        $y\in \RV (f) \, \Ra \, f^{-1}(y)=\{x\}_{i=1}^N$. Mithilfe des \textit{Satzes über die 
            implizite Funktion} finden wir
            \[
                U(x^i)=:U^i
            \]
            disjunkte Umegbungen, so dass eindeutige Lösungen $x^i(t)\in U^i$ existieren von
            \[
                (f+tg)(x)=y \qq \forall |t|<\eps
            \]
            Wir können (zumindest auf evtl. noch kleineren $U^i$) annehmen, dass das Vorzeichen
            von $J_{f+tg}$ konstant ist auf $U^i$. Sei 
            \[
                \eps _2= \frac{\dist(y,f(\Omega\setminus \bigcup_{i=1}^NU^i))}{\|g\|_\infty}.
            \]
            Dann gilt $\dist(y,(f+tg)(\partial \Omega))>0$ $\forall |t|<\eps_2$. Das Lemma gilt somit
            für $\eps = \min(\eps,\eps_2)$
        \item{3)}
            $y\in \CV(f)$. Sei dann $\tilde y\in \RV(f)$.
            \[
                |y-\tilde y|< \frac{\rho}{3} = \frac13 \dist(y,f(\partial \Omega)).
            \]
            Nach Definition \ref{2.9} gilt
            \[
                \deg (f,\Omega,y)=\deg(f,\Omega,\tilde y).
            \]
            Sei $\tilde \eps>0$ mit $\deg(f,\Omega,\tilde y)=\deg(f+tg,\Omega,\tilde y)$ für $|t|<\tilde
            \eps$ (Nach Schritt 2).
            Mit $\eps=\min(\tilde \eps , \frac{\rho}{3\|g\|_\infty})$ gilt:
            \[
                |\tilde y-(f+tg)(x)|\geq \frac{\rho}{3} \qq \forall x \in \partial \Omega
            \]
            Somit gilt
            \[
                |\tilde y-y|< \dist (\tilde y, (f+tg)(\partial \Omega))
            \]
            also
            \[
                \deg(f+tg,\Omega,\tilde y)= \deg (f+tg,\Omega,y)
            \]
            Es folgt
            \[
                \deg(f,\Omega, y)= \deg(f+tg,\Omega,y)
            \]
    \end{description}
    \[ \]
\end{proof}

\begin{theorem}[Existenz und Eindeutigkeit des Brouwer'schen Abbildungsgrades]\label{2.11}
    Es existiert eine eindeutige Abbildung
    \[\deg(f,\Omega,y)\]
    mit den Eigenschaften \textbf{(D1)}-\textbf{(D4)}. Weiter gilt
    \[
        \deg(\cdot, \Omega, y): D_y(\ol \Omega, \R^n)\ra \Z
    \]
    ist konstant auf Zusammenhangskomponente von $D_y(\ol \Omega, \R^n)$. Für $f\in D_y(\ol \Omega,\R^n)$
    ist $\deg(f,\Omega,y)$ gegeben durch
    \[
        \deg(f,\Omega,y)=\sum_{x\in \tilde f ^{-1}(y)} \sign J_{\tilde f}(x),  
    \]
    wobei
    \[
        \tilde f \in D_y^2(\ol \Omega, \R^n)
    \]
    beliebig aus derselben Komponente von $D_y(\ol \Omega,\R^n)$ wie $f$ gewählt werden kann mit $y
    \in \RV(\tilde f)$.
\end{theorem}
   % Schritt 2: Nur stetige Funktionen f
\begin{proof}

Aus Lemma \ref{2.8} und \ref{2.10} folgt, dass $\deg$ wohldefiniert ist und lokal konstant mit 
Werten in $\Z$. Die Abbildung $\deg$ auf Zusammenhangskomponenten.
\begin{description}
\item[(D2)] $\ra$ klar.
\item[(D1)] gilt nachdem diese Bedingung per Konstruktion für die Determinatenformel gilt.
\item[(D3)] Wir wählen $\|f-\tilde f\|_\infty<\dist(y,f(\ol \Omega\setminus(\Omega_1\cap \Omega_2)))$.
    \textbf{D3} gilt per Konstruktion für $\tilde f$, sonst für $f$.
\item[(D4)] folgt aus der Konstruktion von $f$ auf Zusammenhangskomponente.
\end{description}
\[ \]
\end{proof}

\subsubsection*{Beispiele:}
\begin{description}
    \item{1)}
    Sei 
    \[
        f(x_1,x_2)=\begin{pmatrix} x_1-2x_2+\cos(x_1+x_2)\\x_2+2x_1+\sin(x_1+x_2) \end{pmatrix}
    \]
    und
    \[
        g(x_1,x_2):=\begin{pmatrix} x_1-2x_2 \\ x_2+2x_1 \end{pmatrix}
    \]
    Es gilt $|g(x)|=\sqrt{5} |x|$ und $|f(x)-g(x)|=\sqrt{\sin^2(x_1+x_2)+\cos^2(x_1+x_2)}=1$.
    Sei $h(t)=(1-t)g+tf=g+t(f-g)$.
    \begin{align*}
        |h(t)|&\geq |g|-t|f-g|>0 \qq \text{für} \qq |x|>\frac1{\sqrt{5}}\\
        \Ra \, \deg(f,B_r(0),0)&= \deg(g,B_r(0),0)=1 \qq \text{für} \qq r>\frac1{\sqrt{5}}
    \end{align*}
    Somit existiert eine Lösung $x:f(x)=0$.

    \item{2)}
    \begin{theorem}
        Ein stetiges Vektorfeld im $\R^n$, das auf einer Kugeloberfläche überall nach außen zeigt, muss
        auf einem Punkt im Innern der Kugel verschwinden. 
        
        Anders formuliert, sei $f:\ol{B_R(0)} \ra \R^n$ stetig, so dass $f(x)\cdot x>0$ $\forall |x|=R$. 
        Dann existiert ein $x_0\in B_R(0)$ mit $f(x_0)=0$.
    \end{theorem}

    \begin{proof}
        Wir haben $\deg(\Id, B_R(0),y)=1$ für $y\in B_R(0)$. Angenommen, $f(x)\neq 0$ für alle
        $x\in B_R(0)$. Dann gilt $f^{-1}(0)\cap B_R(0)=\varnothing$.
        \[
            \Ra \, \deg(f,B_R(0),0)=0
        \]
        Sei $H(t)=(1-t)\Id+tf$.
        \begin{align*}
            \deg(H(0),B_R(0),0)&=1\\
            \deg(H(1),B_R(0),0)&=0.
        \end{align*}
        Es existiert ein $t_0\in (0,1)$, so dass $H(t_0)\nin D_0(\ol{B_R(0)},\R^n)$.
        \[
            \exists x_0\in \partial B_R(0)\cdot (H(t_0))(x_0)=0.
        \]
        \begin{align*}
            0&=H(t_0)(x_0)\\
            \Ra \, 0&= H(t_0)(x_0)\cdot x_0= \underbrace{(1-t_0)R^n}_{>0}+\underbrace{t_0f(x_0)x_0}_{>0}
            >0
        \end{align*}
        Widerspruch!
        \[ \]
    \end{proof}
\end{description}

\subsection{Der Brouwer'sche Fixpunktsatz}

Der \textit{Brouwer'sche Fixpunktsatz} ist eine Folgerung aus den Eigenschaften des Abbildungsgrades.
Er besagt, dass stetige Abblidungen, die kompakte, konvexe Mengen im $\R^n$ (oder Mengen, die dazu
homöomorph sind) in sich selbst abbilden, einen Fixpunkt besitzen.

\subsubsection{Einschub:}

\begin{description}
        \item[Zerlegung der Eins:]
        Sei $\{U_i\}_{i\in I}$ eine Überdeckung von $X$, $X$ ein topologischer Raum.
        $\{\vp_\lambda\}_{\lambda\in\Lambda}$ ist eine $\{U_i\}_{i\in I}$ untergeordnete lokal endliche
        Zerlegung der Eins, falls gilt
        \begin{itemize}
        \item[-] $\vp_\lambda \in C(U,[0,1]) \,\forall\, \lambda$.
        \item[-] $\forall \lambda\in \Lambda\, \exists \, i \in I: \, \supp \vp_\lambda\subset U_i$
        \item[-] $\sum_{\lambda \in \Lambda} \vp _\lambda=1$
        \item[-] $\forall x \in X \, \exists \, U(x)$ Umgebung von $x$, so dass nur
        \textit{endlich viele} $\lambda \in \Lambda$ existieren, mit $U(x)\cap \supp \vp_\lambda\neq 
        \varnothing$ (lokal endlich).
        \end{itemize}
        \item[Konstruktion] aus einer lokal endlichen Überdeckung $(V_\lambda)_{\lambda\in\Lambda}$ 
        ($\forall x\in X \, \exists \, U(x)$ Umgebung: $U(x)\cap V_\lambda\neq \varnothing$ nur für
         endlich viele $\lambda\in \Lambda$.) in metrischen Räumen $(X,d)$.
        \begin{itemize}
            \item[-] Sei $\alpha(x)=\sum_{\lambda\in \Lambda} \dist(x,X\setminus V_\lambda)>0 \, \forall
            x \in X$, mit $\dist(x,\varnothing):=1$.
            \item[-] $\vp _\lambda(x):=\frac{1}{\alpha(x)}\dist(x,X\setminus V_\lambda) \, \in [0,1]$,
            $\vp_\lambda=0$ für $x\nin V_\lambda$.
        \end{itemize}
        \item[„Konstruktion“] einer lokal endlichen Überdeckung $\{V_\lambda\}_{\lambda\in\Lambda}$, die
        eine offen Überdeckung $\{U_i\}_{i\in I}$ verfeinert (d.h. $\forall \, \lambda\in \Lambda \, 
        \exists i \in I, \, V_\lambda\subset U_i$).
        \item{Die Existenz} einer solchen lokal endlichen verfeinernde Überdeckung zu jeder offenen
        Überdeckung ist die Definition der \textit{Parakompaktheit}.
        \begin{theorem*}
                    Jeder metrische Raum ist parakompakt.
        \end{theorem*}

        \begin{idea}
        Angenommen $I=\N$ also wohlgeordnet. (jede Teilmenge besitzt ein eindeutiges kleinstes Element.)
        Wir setzen für $i\in I, \, n\in \N$.
        \[
            D_{in}=\bigcup_{x\in \Phi(i,n)} B_{2^{-n}}(x)
        \]
        \[
            \Phi(i,n):=\left\{x\in X \, \Big| \, \text{mit}\, \begin{cases}\text{$i$ ist die kleinste
            Zahl, so dass $x\in U_i$}\\ \text{$x\nin D_{jm}$ für $m<n$}\\ B_{3\cdot2^{-n}}\subset U_i  
            \end{cases}\right\}
        \]
        \begin{description}
        \item klar: $\{ D_{in} \}$ verfeinert $\{U_i\}$
        \item[Überdeckung:] Klar, denn für jedes $x\in X$ finden in ein kleinstes $i\in I=\N$ und ein
        $n$ hinreichend groß.
        \item[lokale Endlichkeit:] (s.Artikel).
        Ist $I$ nun eine beliebige Indexmenge, so läßt sich diese wohlordnen (\textit{Zorn'sches Lemma}),
        und die Wohlordnung ist die einzige Eigenschaft von $I$, die wir benutzt haben.
        \end{description}
        \[ \]
        \end{idea}
    \end{description}
\begin{theorem}[Fortsetzung stetiger Funktionen]\label{2.12}
    Sei $X$ ein metrischer Raum, $Y$ normierter Raum. Sei $K\subset X$ abgeschlossen. Sei $F\in C(K,Y)$.
    Dann besitzt $F$ eine stetige Fortsetzung $G:X\ra Y$, so dass
    \[
        G(X)\subset\conv\, F(K)
    \]
\end{theorem}

\begin{proof}
    Wir betrachten die offene Überdeckung
    \[
        \left\{ B_{\rho(x)}(x) \right\}_{x\in X\setminus K} \qq \text{von} \qq X\setminus K
    \]
    mit $\rho(x)=\frac12 \dist(x,K)$.
    Wir wählen nun eine lokal endliche Zerlegung der Eins $\{\vp_\lambda\}_{\lambda\in \Lambda}$, welche
    der offenen Überdeckung untergeordnet ist.
    Sei
    \[
    G(x):= \begin{cases} F(x)& \text{für} \qq x\in K\\
           \sum_{\lambda \in \Lambda } \vp_\lambda(x) \cdot F(x_\lambda) & \text{für} 
           \qq x\in X\setminus K\end{cases},
    \]
    wobei $x_\lambda$ beliebig in $K$ so gewählt ist, dass 
    \[
        \dist(x_\lambda,\supp \vp_\lambda)\leq 2\cdot \dist(K,\supp \vp_\lambda).
    \]
   % 2.3 Der Brouwer'sche Fixpunktsatz
    $G$ ist offensichtlich stetig auf $X\setminus\partial K$ und
    \[
        G(X) \subset \conv\,G(K).
    \]
    Sei also $x_0\in\partial K$ und sei $\eps>0$. Wir wählen $\delta >0$, so dass
    \[
        \|F(x)-F(x_0)\|\leq\eps \qq \forall x\in K \qq \text{mit} \qq d(x,x_0)\leq 9\delta.
    \]
    Es bleibt zu zeigen, dass
    \[
        \|G(x)-F(x_0)\|\leq \eps \qq \forall x\in X\setminus K \qq \text{mit} \qq d(x,x_0)\leq \delta
    \]a
    Sei also $x\nin K$, dann gilt
    \[
        \|G(x)-F(x_0)\|\leq \sum_{\lambda\in \Lambda} \vp_\lambda (x) \|F(x_\lambda)-F(x_0)\|
    \]
    Nach Konstruktion liegt $x_\lambda$ nicht weit von $x$ entfernt, damit nicht weit von $x_0$, falls
    $x\in\supp \, \vp _\lambda$. In der Tat gilt für $x\in \supp\,\vp_\lambda$:
    \begin{align*}
        d(x,x_\lambda)&\leq \dist(x_\lambda,\supp\,\vp_\lambda)
        +\underbrace{\diam(\supp\, \vp_\lambda)}_{=\sup_{x,y\in\supp\,\vp_\lambda} d(x,y)}\\
            &\leq 2\dist(K,\supp \, \vp _\lambda)+\diam(\supp\, \vp _\lambda)
    \end{align*}
    Nachdem $\{\vp _\lambda\}_{\lambda\in\Lambda}$ der Überdeckung $\{B_{\rho(x)}\}_{x\in X\setminus K}$
    untergeordnet ist, existiert $\tilde x\in X\setminus K$, so dass
    \begin{align*}
        \supp\, \vp _\lambda &\subset B_{\rho(\tilde x)}(\tilde x)
    \end{align*}
    \begin{align*}
        \Ra \, \diam (\supp \, \vp _\lambda) &\leq 2\rho(\tilde x)=\dist(\tilde x,K)
        \leq2\dist(K,B_{\rho(\tilde x)}(\tilde x))\leq 2 \dist(K,\supp \vp_\lambda)\\
        \Ra \, d(x_0,x_\lambda)&\leq 4 \dist(K,\supp\, \vp_\lambda).
    \end{align*}
    Und es folgt
    \begin{align*}
        d(x_0,x_\lambda) &\leq d(x_0,x)+d(x,x_\lambda)\\
                &\leq d(x_0,x)+4\dist(K,\supp\,\vp_\lambda)\\
                &\leq d(x_0,x)+8\dist(K,\supp\, \vp_\lambda)\\
                &\leq d(x_0,x)+8d(x_0,x)=9d(x_0,x)\leq \delta.
    \end{align*}
    Nach Wahl von $\delta$ gilt
    \[
        \|F(x_{\lambda})-F(x_0)\|\leq \eps \qq \forall \lambda\in\Lambda:\, \vp_\lambda(x)\neq 0.
    \]
    Somit gilt
    \[
        \|G(x)-G(x_0)\|\leq \eps \qq \text{für} \qq d(x,x_0)\leq \delta.
    \]
\end{proof}

\begin{remark}
    Mit Hilfe dieses Satzes und dem Abbildungsgrad lässt sich der sog. „Igelsatz“ zeigen, der besagt,
    dass man einen Igel nicht stetig kämmen kann. Übungsaufgabe. 
\end{remark}

\begin{theorem}[Brouwerscher Fixpunktsatz]\label{2.13}
    Sei $K$ ein topologischer Raum, der zu einer kompakten konvexen Teilmenge des $R^n$ homöomorph ist.
    Sei $f\in C(K,K)$. Dann besitzt $f$ einen Fixpunkt.
\end{theorem}

\begin{proof}
    \begin{description}
    \item{1.}
    $K=\ol{B_r(0)}\subset\R^n$. Falls ein Fixpunkt auf dem Rand existiert, dann sind wir fertig.
    Ansonsten gilt für $H(t)=\Id-tf$, dass $0\nin H(t)(\partial B_r(0))$, nachdem
    \[
        |H(t)(x)|\geq |x|-t|f(x)|\geq (1-t)r>0 \qq \text{für} \qq 0\leq t< 1, \qq x\in\partial B_r(0).
    \]
    Nach der Annahme der Nichtexistenz eines Fixpunktes auf $\partial B_r(0)$ ist auch $H(1)(x)\neq 0$
    $\forall x \in \partial B_r(0)$.
    \[
        \Ra \, \deg(\Id-f,B_r(0),0)=\deg(\Id,B_r(0),0)=1.
    \]
    Somit existiert $x\in B_r(0)$ mit $x=f(x)$.
    \item{2.}
    Sei nun $K\subset \R^n$ konvex, kompakt. Für ein $\rho>0$ gilt $K\subset B_\rho(0)$ und gemäß Satz
    \label{2.12} können wir $f$ stetig durch $g$ auf $\ol{B_{\rho}(0)}$ fortsetzen mit
    \[
        g\left( \ol{B_\rho (0)} \right)\subset \conv \, K=K.
    \]
    Nach 1. finden wir $x\in \ol{B_\rho(x)}$ mit $x=g(x)$.
    Es gilt $g(x)\in K$, somit folgt $x\in K$ und wir haben $x \in K$ mit $f(x)=x$.
    \item{3.}
    Sei $K$ homöomorph zu $K^\star \subset \R^n$ kompakt, konvex. Sei $h$ die entsprechende Homöomorphie.
    nach 2. besitzt $h\circ f \circ h^{-1}$ einen Fixpunkt $x^\star\in K^\star$. Damit ist aber $x=h^{-1}
    (x^\star)\in K$ ein Fixpunkt von $f$.
    \end{description}
    \[ \]
\end{proof}

\begin{remark}
    \begin{description}
    \item{1.}
    Die Bedingungen sind tatsächlich notwendig. Gegenbeispiele siehe Übung.
    \item{2.}
    Es existiert auch eine stetige Abbildung $f\in C(\ol{B_1(0)},\ol{B_1(0)})$, $B_1\subset X$ separabler
    $\infty$-dimensionaler Hilbertraum ohne Fixpunkt.
    \begin{description}
    \item{$\ra$}
    Beispiel von Kakutani später.
    \item{$\ra$}
    \textit{Schauderscher Fixpunktsatz}.
    \end{description}
    \item{3.}
    Im eindimensionalen Fall ist \ref{2.13} nichts anderes als der Zwischenwertsatz angewendet auf
    $x-f(x)$.
    \item{4.}
    Vergleich mit dem \textit{Banach'schen Fixpunktsatz}: Wesentlich geringere Anforderungen an den
    Operator (nur Stetigkeit), dafür hohe Anforderungen an den Raum (endlichdim., kompakt, konvex).
    Wir bekommen keine Eindeutigkeitsaussage.
    \end{description}
\end{remark}

\subsubsection*{Ein Anwendugsbeispiel}

Existenz positiver (bzw. nichtnegativer) Eigenwerte und Eigenfunktionen. Sei $A=(a_{ij})_{i,j=1}^n$ eine
Matrix und sei $a_{ij}\geq 0$ $\forall (i,j)$. Dann existiert $\lambda\geq 0$, $x=(x_i)_{i=1}^n$,
$x_i\geq 0$ $\forall i$ mit
\begin{align}\label{6}
    Ax=\lambda x
\end{align}

\begin{proof}
    Sei 
    \[
        K=\{ x\in \R^n \, |\, x\geq 0, \, \sum_{i}x_i=1 \}
    \]
    kompakt und konvex.
    \begin{description}
        \item{1)}
        Falls $Ax=0$ für ein $x\in K$, gilt \ref{6} mit $\lambda=0$.
        \item{2)}
        Sei $Ax\neq 0$ $\forall x\in K$. Dann existiert $\alpha>0$ mit
        \[
            \sum_i(Ax)_i\geq \alpha \qq \forall x\in K.
        \]
        Es sei also
        \[
            f:x\mapsto \frac{Ax}{\sum_i (Ax)_i}; \qq (f(x))_i \geq 0 \qq \forall i
        \]
        und es gilt
        \[
            \sum_i (f(x))_i=1 \qq \forall x\in K.
        \]
        Dann ist $f(K)\subset K$ und nach Satz $\ref{2.13}$ existiert ein Fixpunkt $x\in K$ mit $x=f(x)$,
        d.h.
        \[
            Ax=\lambda x \qq \text{mit} \qq \lambda =\sum_i(Ax)_i >0
        \]
    \end{description}\[ \]
\end{proof}

\subsection{Weitere Eigenschaften des Abbildungsgrades}

Sei $m<n$, wir identifizieren im Folgenden den $\R^m$ mit dem Unterraum
\[
    \{x\in \R^n \, | \, x_{m+n}=…=x_n=0\}\subset \R^n.
\]

\begin{theorem}[Reduktionseigenschaften des Abbildungsgrades]\label{2.14}
    Sei $\Omega \subset \R^n$ offen und beschränkt, $\Omega \cap \R^m \neq \varnothing$ und sei
    \[
        f:\ol \Omega \ra \R^m\, \text{stetig,} \qq g=\Id-f
    \]
    Sei $y\in \R^m$, $y\nin g(\partial \Omega)$. Dann gilt
    \[
        \deg(g,\Omega,y)=\deg(g|_{\ol{\Omega}\cap\R^m},\Omega|_{R^m},y).
    \]
\end{theorem}

\begin{proof}
    Sei $\Omega_m=\Omega\cap\R^m\neq \varnothing$ offen, beschränkt im $\R^m$, und
    $g_m:=g|_{\ol{\Omega_m}}$. Es gilt
    \[
        \partial \Omega_m=\partial\Omega\cap\R^m.
    \]
    und
    \[
        g_m(\ol{\Omega_m})\subset\R^m,\, y\nin g_m(\partial\Omega_m).
    \]
    $\Ra$ $\deg(g_m,\Omega_m,y)$ ist definiert (als Abbildungsgrad im $\R^m$).
    \begin{description}
        \item{1)}
        Sei $g\in C^1(\ol \Omega)$ und $y\in \RV(g)$ und $x\in \Omega$, $g(x)=y$.
        \[
        \Ra \, x=y=f(x)\in \R^m \qq \Ra \qq x\in \Omega_m.
        \]
        Somit haben wir
        \[
            g^{-1}(y)=g_m^{-1}(y).
        \]
        Nach der Determinantenformel genögt es nun zu zeigen, dass
        \[
            J_{g_m}(x)=J_g(x) \qq \text{für} \qq x\in \Omega_m.
        \]
        Es sei $\Id_k$ die $k\times k$-Einheitsmatrix. Es gilt
        \begin{align*}
            J_{g_m}(x)&=\det(\Id_m-f'(x)) \qq (f|_{\Omega_m}')\\
            J_g(x)&= \det\begin{pmatrix} \Id_m-(\partial_jf_i(x)) & -\partial_jf_i\\
            0& \Id_{n-m}\end{pmatrix} \qq \text{wegen} \, f(\ol\Omega)\subset\R^m.
        \end{align*}
        Die gewünschte Aussage folgt sofort durch Entwicklung der Determinante nach den letzten
        $n-m$ Zeilen.
        \item{2)}
        Der allgemeine Fall folgt durch Wahl von $\tilde g=\Id-\tilde f$ hinreichend nah an $g$, so dass
        $\tilde g$ die in 1. geforderten Eigenschaften besitzt.
    \end{description}
    \[ \]
\end{proof}
   % 2.4 Weiter Eigenschaften des Abbildungsgrades
Es sei nun wieder $\Omega \subset \R^n$ offen, beschränkt, $f\in C(\ol \Omega,\R^n)$. Nach Satz \ref{2.2}
gilt, dass $\deg(f,\Omega,\cdot)$ konstant ist auf Zusammenhangskomponenten von $\R^n\setminus f(\partial 
\Omega)$. Wir bezeichnen diese Zusammenhangskomponenten mit $\{G_j\}_{j\in I}$ und schreiben $\deg(f,
\Omega,y)=\deg(f,\Omega,G_j)$ für $y\in G_j$

\begin{theorem}[Produktregel]\label{2.15}
    Sei $\Omega\subset\R^n$ offen, beschränkt. Seien $G_j$ die Zusammenhangskomponenten von $\R^n
    \setminus f(\partial \Omega)$ für $f\in C(\ol\Omega, \R^n)$ und sei
    \[
        g\circ f \in D_y(\ol \Omega,\R^n), \qq y\in \R^n
    \]
    Dann gilt
    \begin{align}\label{9}
        \deg(g\circ f, \Omega, y)=\sum_j\deg(f,\Omega,G_j)\cdot \deg(g,G_j,y)
    \end{align}
\end{theorem}

\begin{proof}
    (Endlichkeit der Summe) Es gilt $f(\ol\Omega)$ ist kompakt, somit existiert $r>0$ mit $f(\ol \Omega)
    \subset B_r(0)$. Nachdem $g^{-1}(y)$ abgeschlossen ist, gilt $g^{-1}(y)\cap B_r(0)$ kompakt ist.
    $\{G_j\}_{j\in I}$ ist eine offene Überdeckung dieser Menge, es genügen somit endlich viele $\{  
    G_j \}_{j=1}^N$ um $g^{-1}(y)\cap B_r(0)$ zu überdecken.
    \begin{description}
    \item{1)}
    Wir nehmen an, dass $f\circ g\in C^1(\ol \Omega)$, $y\in \RV(g\circ f)$. Es gilt nach der
    Kettenregel, dass
    \[
        (g\circ f)'(x)=g'(f(x))\circ f'(x)
    \]
    Die Behauptung folgt dann durch eine längere Rechnung.
    \begin{align*}
        \deg(g\circ f,\Omega,y)&=\sum_{x\in(g\circ f)^{-1}(y)}\sign J_{g\circ f}(x)=\sum_{x\in (g\circ f)
        ^{-1} (y)}(\sign J_g(f(x)))\cdot(\sign J_f(x))\\
        &= \sum_{z\in g^{-1}(y)} \sign J_g(z)\cdot \sum_{x\in f^{-1}(z)} \sign J_f(x)
        =\sum_{z\in g^{-1}(y)} \sign J_g(z)\cdot \deg(f,\Omega,z)
    \end{align*}
    Mit der Überdeckung $\{G_j\}_{j\in I}$ von $g^{-1}(y)$ gilt
    \begin{align*}
        \deg(g\circ f,\Omega,y)&=\sum_{j=1}^m\sum_{z\in g^{-1}(y)\cap G_j}\sign J_gi(z)\cdot \deg(f,\Omega,z)
        =\sum_{j=1}^N \deg(f,\Omega,G_j) \cdot \sum_{z\in g^{-1}(y)\cap G_j} \sign J_g(z)\\
        &= \sum_{j=1}^N \deg(f,\Omega,G_j)\cdot \deg(g,G_j,y)
    \end{align*}
    Diese Formel gilt nun natürlich auch für $y\in \CV(g\circ f)$ und für $g$ nur stetig. Etwas
    problematischer ist der Fall, dass $f$ nur stetig ist, da sich bei der Modifikation von $f$ die
    Mengen $G_j$ ändern.

    Es sei
    \[
        L_l=\{z\in \R^n\setminus(\partial \Omega)\, | \, \deg(f,\Omega,z)=l\}
    \]
    Für $l\neq 0$ gilt, dass $L_l$ aus einer Vereinigung von Mengen $G_j$ bestehen muss. Sei nun
    $\tilde f\in C^1 (\ol \Omega)$, so dass $|f(x)-\tilde f(x)|< \frac 12 \dist(g^{-1}(y),
    f(\partial\Omega))$ für $x\in \ol \Omega$. Wir definieren $\tilde G_j$ und $\tilde L_l$ entsprechend
    für die Funktion $\tilde f$.
    Es gilt
    \[
        L_l\cap g^{-1}(y)=\tilde L_l\cap g^{-1}(y).
    \]
    nach Satz \ref{2.1} (iii). ($\|f(x)-\tilde f(x)\|<\dist(y,f(\partial \Omega)\qq \forall
    x\in \partial \Omega)\qq \Ra \qq \deg(f,\Omega,y)=\deg(\tilde f,\Omega,y)$.) Es gilt somit
    \begin{align*}
        \deg(g\circ f, \Omega, y)&=\deg(g\circ\tilde f,\Omega,y)\\
        &=\sum_{j}\deg(f,\Omega,\tilde G_j)\cdot \deg(g,\tilde G_j,y)=\sum_{l\neq0} l\cdot\deg( g,\tilde
        L_l,y)=\sum_{l\neq 0} l\cdot \deg(g,L_l,y)\\
        &=\sum_j \deg(f,\Omega,G_j)\cdot \deg(g,G_j,y)
    \end{align*}
    \end{description}
\end{proof}

Eine wichtige Anwendung der Produktformel ist das folgende
\begin{theorem}[Jordanische Kurvensatz]\label{2.16}
Es seien $C_1$ und $C_2$ zwei kompakte, zueinander Homeomorphe Teilmengen des $\R^n$. Dann besitzen
$\R^n\setminus C_1$ und $\R^n\setminus C_2$ die selbe Anzahl von Zusammenhangskomponenten.
\end{theorem}
 % Nach dem Einschub geht es weiter mit 2.4 Eigenschaften des Abbildungsgrades.
\section{Der Fixpunktsatz von Kakutani und eine Anwendung in der Spieltheorie}

Im Folgenden beweisen wir die sogenannten \textit{Nash-Gleichgewichte} in n-Personen.
Als Vorbereitung ist es notwendig, den \textit{Brouwer'schen Fixpunktsatz} auf mengenwertige Funktionen
zu verallgemeinern.

Für $K\subset \R^n$ konvex, kompakt, bezeichnen wir mit $\CS(K)$ die Menge der konvexen, nichtleeren
(Teil-)Mengen von $K$.

\begin{theorem}[Kakutani]\label{2.16}
    Sei $K\subset \R^n$ kompakt, konvex $f:K\ra \CS(K)$. Falls die Menge
    \[
        \Gamma := \{(x,y)\in K\times K\, |\, y\in f(x)\}
    \]
    in $K\times K$ abgeschlossen ist, dann existiert $x\in K$ mit $x\in f(x)$.
\end{theorem}

\begin{proof}
    In einer Dimension: $K=[v_1,v_2]$, wir wählen $y_j\in f(v_j)$ und definieren
    \[
        f^1(x)=\sum_{j=1}^{2} \lambda_j(x)y_j,
    \]
    wobei $\lambda _j(x)$ die baryzentrischen Koordinaten von $x$ bezeichne, d.h.
    \[
        \lambda_j(x)\geq 0, \qq \sum_{i=1}^2 \lambda_1(x)=1,\qq x=\sum_{j=1}^2 \lambda_j(x)v_j
    \]
    Nach Konstruktion ist $f^1:K\ra K$ stetig, besitzt mit Brouwer also einen Fixpunkt. Das hilf
    noch nicht viel. Wir betrachten die $k-1$-te baryzentrische Unterteilung der Menge $K$ gegeben
    durch die Vertizes 
    \[
        \{v_i\}_{i=1}^{2^{k-1}+1}
    \]
    Wir wählen wieder $y_j\in f(v_j)$ und definieren $f^k(v_j)=g_j$ interpolieren auf den unterteilten
    Intervallen wie gehabt durch baryzentrische Koordintaten. Es gilt
    \[
        f^k: K\ra K
    \] 
    ist stetig, also mit Fixpunkt
    \begin{align}\label{7}
        x^k=\sum_{i=1}^2 \lambda_i^kv_i^k= \sum_{i=1}^2 \lambda_i^ky_i^k \qq \text{mit} \qq 
        y_i^k=f^k(v_i^k)
    \end{align}
    in einem Teilintervall  (in einem Simplex der $k$-fachen baryzentrischen Koordinaten). Es gilt:
    \[
        \left(x^k,\lambda_1^k,\lambda_2^k,y_1^k,y_2^2\right)\in K\times[0,1]^2\times K^2
    \]
    Somit konvergiert eine Teilfolge in $K$ gegen
    \[
        \left(x^0,\lambda_1^0,\lambda_2^0,y_1^0,y_2^0\right)
    \]
    Nachdem die Teilintervalle zu Punkten degenerieren gilt:
    \[
        v_i^k\ra x^0 \qq i=1,2
    \]
    Wir haben 
    \[
        \left(v_i^k,y_i^k\right)\in\Gamma\ra \left(v_i^0,y_i^0\right)\in \Gamma \qq 
        \text{wegen Abgeschlossenheit von }\Gamma
    \]
    Somit gilt $y_i^0\in f(x^0)$, $i=1,2$. Mit (\ref{7}) folgt
    \[
        x_0 =\lim_{k\ra\infty} \sum_{i=1}^2\lambda_i^ky_i^k=\sum_{i=1}^2 \lambda^0_iy_i^0 
        \in f\left(x^0\right),
    \]
    $f\left(x^0\right)$ ist konvex. 

    In $n$ Dimensionen nehmen wir zunächst an, dass $K$ ein Simplex (Konvexe Hülle von $n+1$
    nichtdegenerierten Vertizes) ist. Wir wenden wiederholt die Baryzentrische Unterteilung an,
    der Durchmesser der Simplizes schrumpft mit einem Faktor $\frac{n}{n+1}$.
    Zur $k$-ten baryzentrischon Unterteilung mit Vertizes $\{v_j\}_{j=1}^{N_k}$. Wählen wir wieder
    $y_j\in f(v_j)$ mit $f^k(x)$ durch Interpolation mittels baryzentrischen Koordinaten in jdn.
    Unter-Simplex. Wir erhalten wieder Fixpunkte $x^k$, und eine Teilfolge von
    \[
        (x^k,\lambda_1^k,…,\lambda_{n+1}^k,y_1^k,…,y_{n+1}^k)
    \]
    konvergiert im „richtigen“ Unter-Simplex. Der Rest des Beweises folgt exakt wie im eindimensionalen
    Fall. Ist $K$ kein Simplex, so können wir $f^k$ jeweils auf einen Simplex $\tilde K\supset K$ stetig
    fortsetzen und verfahren wie im Beweis des \textit{Brouwer'schen Fixpunktsatzes}.
    \[ \]
\end{proof}

\subsection*{Spieltheorie:}

Ein $n$-Personen Spiel besteht aus $n$ Spielern, wobei der $i$-te Spieler $m_i$ mögliche Aktionen
ausführen kann. Die Menge der möglichen Aktionen des $i$-ten Spielers bezeichnen wir mit
$\Phi_i=\{ 1,…,m_i \}$. Nachdem jeder Spieler eine Aktion ausgeführt hat, wird abgerechnet. Der $i$-te
Spiele bekommt den Payoff
\[
    R_i(\vp), \qq \vp=(\vp_1,…,\vp_n)\in \Phi=\prod_{j=1}^n \Phi_j
\]
Wir betrachten den Fall, dass das Spiel häufig ausgeführt wird, und jeder Spieler seine Aktion nach
einer Wahrscheinlichkeitsverteilung (einer Strategie) auf $\Phi_i $ wählt. Diese bezeichenen wir mit
\[
    S_i=\{s_i^1,…,s_i^{m_i}\},\qq s_i^k\geq 0, \qq \sum_{k=1}^{m_i}s_i^k=1 \qq \forall i
\]
Die Menge aller möglichen Strategien für den $i$-ten Spieler nennen wir $S_i$. $S_i^k$ ist die
Wahrscheinlichkeit, dass der $i$-te Spieler die Aktion $k$ wählt. Die Wahrscheinlichkeit, dass jeder
Spieler eine bestimmte Aktion $s_i^k$ wählt ist
\[
    S(\vp)=\prod_{i=1}^n s_i(\vp), \qq s_i(\vp)=s_i^{k_i}
\]
mit $\vp=(k_1,…,k_n)\in \Phi$. Wir nehmen hier an, dass die Spieler ihre Aktionen unabhängig voneinander
wählen. Der mittlere Payoff für Spieler $i$ ist dann der Erwartungswert für $R_i(\vp)$. Wir nennen diesen
$R_i(s)$ und erhalten
\[
    R_i(s)=\sum_{\vp\in\Phi}s(\vp) K_i(\vp)
\]
Nach Konstruktion ist
\[
    R_i(s)
\]
stetig in $s$. Was ist nun die optimale Strategie für Spiel $i$? Falls alle anderen Spiele nach einer
bekannten  Strategie handeln, ist es optimal $\ol{s}_i$ so zu wählen, dass
\begin{align}\label{8}
    R_i(s\setminus \ol{s_i})=\max_{\tilde s_i \in S_i} R_i(s\setminus \tilde s_i)
\end{align}
$s\setminus \tilde s_i$: Strategiekomination, die sich ergibt, wenn man $s_i$ durch  $\tilde s_i$
ersetzt. Wir bezeichnen mit $B_i(s)$ die Menge aller Strategien $\ol s_i\in S_i$, die $(\ref{8})$ erfüllen.
Insbesondere gilt $\ol s_i\in B_i(s)$ genau dann wenn $\ol s_i^k=0$ falls
\[
    R_i(s\setminus (0,…,1,…)) < \max_{1<l<m_i}R_i(s\setminus(0,…,1,…)) \qq \text{mit der 1 als $l$-ter
    Eintrag}
\]
Insbesondere existiert immer eine reine (immer die selbe Aktion wird gewählt) optimale Strategie
zu gegebenem $S$. Weiter sieht man, dass $B_i(s)$ konvex ist und nicht leer. Es seien nun $s,\ol s\in S$.
Wir nennen $\ol s$ die optimale Antwort gegen $s$, falls $\ol s_i\in B_1(s)$ $i=1,…,n$.
Es sei $B(s)$ die Menge aller optimalen Antworten gegen $s$, wir haben
\[
    B(s)=\prod_{i=1}^n B_i(s).
\]

\begin{defi}[Nash-Gleichgewicht]
    Eine Strategiekombination $\ol s\in S$ heißt Nash-Gleichgewicht, falls gilt
    \[
        \ol s \in B(\ol s)
    \]
    d.h. $\ol s$ ist eine optimale Antwort gegen sich selbst. Anders gesagt, $\ol s$ ist ein 
    Nash-Gleichgewicht, falls kein Spieler durch eigene Änderung seiner Strategie seinen Payoff erhöhen
    kann.
\end{defi}

\subsection*{Beispiele:}

%\begin{description}
    \textbf{Das Gefangenendilemma:}
    Zwei Gefangene können entweder schweigen, oder mit der Polizei kooperieren, das heißt den jeweils
    anderen belasten.
    \[
        \begin{array}{c|cc||c|cc} R_1& k_2& s_2& R_2&k_2&s_2\\ \hline k_1&0&2&k_1&0&-1\\
            s_1&-1&1&s_1&2&1\end{array}
    \]
%\end{description}
\begin{theorem}\label{2.18}
    Jedes $n$-Personen Spiel besitzt ein Nash-Gleichgewicht.
\end{theorem}

\begin{proof}
    Wir betrachten $s\mapsto B(s)$, $S$ ist kompakt, konvex, $B(s)\subset S$ ist nicht leer, konvex für
    alle $s$. Abgeschlossenheit des Graphen folgt mit Folgenkompaktheit.
    \[
        s^n\in S, \qq \ol s^m\in B(s^m), \qq s^m\ra s, \qq \ol s^k\ra \ol s
    \]
    Wir haben $R_i(s^m\setminus \tilde s_i)\leq R_i(s^m\setminus \ol s_m)$ für alle $\tilde s_i \in S_i$
    mit Stetigkeit von $R_i$ folgt sofort
    \[
        R_i(s\setminus \tilde s_i) \leq R_i(s\setminus \ol s) \qq \forall s_i\in S_i
    \]
    Somit ist $\ol s \in B(s)$ und $\Gamma= \{ (s,B(s))\subset S^2 \}$ ist abgeschlossen.
    Der \textit{Fixpunktsatz von Kakutani} garantiert Existenz von
    \[
        \ol s\in S \qq \text{mit} \qq \ol s \in B(\ol s).
    \]
\end{proof}
   % 2.5 Der Fixpunktsatz von Kakutani und eine Anwendung in der Spieltheorie
\section{Der Leray-Schauder Grad und der Schauder'sche Fixpunktsatz}

Das Ziel ist es, den Fixpunktsatz von Brouwer auf $\infty$-dimensionale Räume zu verallgemeinern.
So einfach ist das nicht.

\begin{theorem}[Kakutani]\label{3.1}
    Sei $H$ ein $\infty$-dimensionaler separabler Hilbertraum. Dann existiert eine stetige Abbildung
    \[
        f:H\ra H,
    \]
    welche die Einheitenkugel in sich selbst abbildet und die keinen Fixpunkt besitzt.
\end{theorem}

\begin{proof}
    Sei $(e_n)_{n\in \Z}$ eine Orhonormalbasis von $H$. Für $x=\sum_{j=-\infty}^\infty\alpha_je_j$
    definieren wir
    \[
        U(x):=\sum_{j=-\infty}^\infty \alpha _je_{j+1};\qq f(x):=\frac12 (1-\|x\|)e_0+U(x)
    \]
    Zu zeigen:
    \begin{description}
    \item{-}
    $U$ ist stetig.
    \item{-}
    $\|U(x)\|<1$ für $\|x\|<1$
    \item{-}
    $f$ besitzt keinen Fixpunkt in $\ol{B_1(0)}$
    \begin{itemize}
        \item bei $x=0$.
        \item auf dem Rand.
        \item im Innern.
    \end{itemize}
    \end{description}
    \[  \]
\end{proof}

\subsection{Der Abbildungsgrad auf endlichen Banachräumen}

\begin{defi}\label{3.2}
    Sei $X$ ein (reeller) Banachraum von Dimension $n$, und sei $\Phi$ ein Isomorphismus von $X$ nach 
    $\R^n$. Für $f\in D_y(\ol \Omega,X)$, $\Omega \subset X$ offen, beschränkt, $y\in X$ definieren wir
    \[
        \deg(f,\Omega,y)=\deg(\Phi\circ f \circ \Phi^{-1},\Phi(\Omega),\Phi(y))
    \]
\end{defi}

\begin{prop}\label{3.3}
    Der in Definition \ref{3.2} definierte Abbildungsgrad ist unabhängig von der Wahl von $\Phi$.
\end{prop}

\begin{proof}
    Sei $\Psi$ ein zweiter Isomorphismus. Es gilt $A=\Psi\circ\Phi^{-1}\in \Gl(n)$. Wir schreiben
    \[
        f^\star:=\Phi \circ f \circ \Phi^{-1}, \qq y^\star\in \RV(\tilde f^\star)
    \]
    und wählen wie üblich $\tilde f^\star\in D^1(\Phi(\ol\Omega),\R^n)$ aus derselben Komponente von
    $D_y(\Phi(\ol \Omega),\R^n)$ wie $f^\star$, so dass $y^\star\in \RV(\tilde f^\star)$. Dann ist auch
    \[
        A\circ \tilde f^\star \circ A^{-1}\in D^1(\Psi (\ol \Omega),\R^n)
    \]
    in derselben Komponente von $D_y(\Psi(\ol \Omega),\R^n)$ wie $A\circ f^\star\circ A^{-1}=\Psi\circ
    f\circ \Psi^{-1}$.
    Nach der Kettenregel gilt
    \begin{align*}
        J_{A\circ \tilde f^\star\circ A^{-1}}(Ay^\star)&= \det(A) J_{\tilde f^\star}(y^\star)
        \det(A^{-1})=J_{\tilde f^\star} (y^\star).
    \end{align*}
    Und es folgt
    \[ 
        \deg(\Psi\circ f \circ\Psi^{-1},\Psi(\Omega),\Psi(y))
        = \deg(\Phi\circ f \circ\Phi^{-1},\Phi(\Omega),\Phi(y))
    \]
\end{proof}

Die Eigenschaften des Abbildungsgrades im $\R^n$ übertragen sich mit Definition \ref{3.2} auf den
Abbildungsgrad in endlichdimensionalen Banachräumen. Das gilt ebenfalls für die Reduktionseigenschaft aus
Satz \ref{2.14}, nachdem $\Phi:X\ra\R^n$ so gewählt werden kann, dass $\Phi(X_1)=\R^n$, wobei $X_1$ ein
$m$-dimensionaler Untervektorraum von $X$ ist.

\subsection{Kompakte Operatoren}

Es seien $X$ und $Y$ Banachräume, $\Omega\subset X$ (nicht unbedingt offen).

\begin{defi}
    Ein stetiger Operator $F:\Omega\ra Y$ heißt
    \begin{description}
        \item{•} \textit{endlichdimenisional}, falls ein endlichdimensionaler Untervektorraum $Y_1\subset Y$
        existiert, sodass $F(\Omega)\subset Y_1$.
        \item{•} \textit{kompakt}, falls $F$ beschränkte Teilmenge von $\Omega$ auf relativ kompakte
        Teilmengen von $Y$ abbildet.
    \end{description}
    Wir bezeichnen die Menge der endlichdimensionalen Operatoren von $\Omega$ nach $Y$ mit
    \[
        \ms F(\Omega,Y)
    \]
    und die kompakten mit
    \[
        \ms C(\Omega,Y)
    \]
\end{defi}

\begin{remark}
    \begin{description}
        \item{1)}
        Es gilt
        \[
            \ms F(\Omega,Y)\subset \ms C(\Omega,Y)\subset C(\Omega,Y)
        \]
        \item{2)}
        Falls $\Omega$ kompakt ist, so gilt
        \[
            \ms C(\Omega,Y)\subset C(\Omega,Y)
        \]
        \item{3)}
        Falls $\dim(Y)<\infty$, so gilt
        \[
            \ms F(\Omega,Y)=\ms (\Omega,Y)
        \]
        \item{4)}
        Falls $\Omega\subset\R^n$ beschränkt, so gilt
        \[
            \ms F (\ol \Omega,\R^n)=\ms C(\ol \Omega,\R^n)=C(\ol\Omega,\R^n)
        \]
    \end{description}
\end{remark}
 % 3 Der Leray-Schauder Grad und der Schauder'sche Fixpunktsatz \\ 
                   % 3.1 Der Abbildungsgrad auf endlichen Banachräumen \\ 3.2 Kompakte Operatoren
\subsubsection*{Korrekturen}

Bei Satz \ref{2.15}. Vorsicht vor der unbeschränkten Komponente $K_\infty$ von $\R^n\setminus f(\partial
\Omega)$! Diese taucht jedoch in der Produktformel ausschließend nicht auf, denn
\[
    f(\ol \Omega)\subset B_r(0), \, \text{mit}\, f(\ol \Omega) \text{ kompakt}
\]
\[
    \Ra \, \deg(f,\Omega,K_\infty\cap B_{2r}(0))=0 \qq \text{und} \qq f(\ol \Omega) \cap (\R^n\setminus
    B_{2r}(0))=\varnothing
\]
somit können wir vereinfacht schreiben
\[
    \deg(f,\Omega,K_\infty)=0
\]
Der in der Produktformel aufauchende Term
\[
    \deg(f,\Omega,K_\infty)\cdot \deg(g,K_\infty,y)
\]
ist also unproblematisch, da er verschwindet.\\[0.5cm]

\begin{lem}\label{3.5}
    Sei $K\subset X$ kompakt, $\eps>0$. Dann existiert ein endlichdimensionaler Untervektorraum $X_\eps
    \subset X$ und eine stetige Abbildung
    \[
        P_\eps:K\ra X_\eps,
    \]
    so dass
    \[
        \|P_\eps(x)-x\|\leq \eps \qq \forall x\in K
    \]
\end{lem}

\begin{proof}
    Es sei $\{x_j\}_{j=1}^N\subset K$, so dass $\bigcap _{j=1}^NB_\eps(x_j)$ eine Überdeckung von $K$
    darstellt. Sei $\{\Phi_j\}_{j=1}^N$ eine Zerleung der Eins, welche $\{B_\eps(x_j)\}_{j=1}^N$
    untergeordnet ist. Wir definieren
    \[
        P_\eps(x):= \sum_{j=1}^N \Phi_j(x)x_j
    \]
    Damit gilt
    \[
        \|P_\eps(x)-x\|=\|\sum_{j=1}^N\Phi_j(x)x_j-\sum_{j=1}^N \Phi_j(x)x_j\leq \sum_{j=1}^N \Phi_j(x)
        \|x-x_j\|\leq \eps\|
    \]
\end{proof}

\begin{theorem}\label{3.6}
    $X,Y$ Banachräume. Sei $\Omega\subset X$ beschränkt. Dann gilt
    \[
        \ol{\ms F(\Omega,Y)}^{C{\Omega,Y}}=\ms C(\Omega,Y)
    \]
    (D.h. der Abschluss von $\ms F(\Omega,Y)$ bzgl. der Norm der gleichmäßigen Konvergenz ist der Raum
     der kompakten Operatoren $\ms C(\Omega,Y)$.)
\end{theorem}

\begin{proof}
    Es sei $(F_n)_{n\in \N}$ eine Folge in $\ms C(\Omega,Y)$, sodass  $F_n\ra F$. Wir nehmen an, dass
    $F\nin \ms C(\Omega,Y)$. Dann existiert eine Folge $(x_n)_{n\in \N}$, so dass
    \[
        \|F(x_n)-F(x_l)\|\geq \rho>0 \qq\text{für}\qq k\neq l
    \]
    Es sei $n$ so groß, dass
    \[
        \|F_n-F\|_{\ms C(\Omega,Y)}<\frac{\rho}{4}
    \]
    Damit gilt
    \[
        \|F_n(x_k)-F_n(x_l)\|\geq \frac\rho2.
    \]
    Das ist ein Widerspruch zur Annahme, dass $F_n\in \ms C(\Omega,Y)$. Somit gilt auch
    \[
        \ol{\ms F(\Omega,Y)}^{\ms C(\Omega,Y)}\subset \ms C(\Omega,Y)
    \]
    Sei andererseits $F\in \ms C(\Omega,Y)$,und sei $K=\ol{F(\Omega)}$ kompakt. sei $P_\eps$ gewählt wie
    in Lemma \ref{3.5}. Es gilt
    \[
        F_\eps=P_\eps\circ F\in \ms F(\Omega,Y) \qq \text{und}\qq F_\eps\ra F
    \]
\end{proof}

Zur Definition des Abbildungsgrades auf unendlichdimensionalen Räumen betrachten wir insbesondere
kompakte Störungen der Identität (d.h. Operatoren der Form $\Id+F$, $F\in \ms C$). Wir bezeichnen deshalb
zunächst einige Interessante Eigenschaften von $\Id+F$.

\begin{lem}\label{3.7}
    Sei $X$ ein Banachraum und $\Omega\subset X$ beschränkt und abgeschlossen, $F\in \ms C(\Omega,Y)$.
    Dann ist $\Id + F$ eine eigentliche Abbildung (d.h. $(\Id+F)^{-1}(K)$ ist kompakt für $K\in X$ 
    kompakt), die abgeschlossene Teilmengen von $\Omega$ auf abgeschlossene Teilmengen von $X$ abbildet.
\end{lem}

\begin{proof}
    Sei $A\subset \Omega$ abgeschlossen, und sei
    \[
        y_n=(\Id+F)(x_n), \qq \text{mit } x_n \in A,
    \]
    so dass $y_n \ra y$ in $X$. Zu zeigen ist, dass $y\in (\Id+F)(A)$. Es gilt
    \[
        y_n-x_n=F(x_n),
    \]
    somit folgt $y_n-x_n\ra z$ in $X$, nach Extraktion einer Teilfolge. Es folgt $x_n\ra x \in A$, da
    $A$ abgeschlossen und $y_n$ konvergent nach Annahme. Wir haben aber $x=y-z\in A$. Nachdem aber
    $y=x+F(x)\in (Id+F)(A)$ folgt die Behauptung.

    Sei nun $\Omega$ abgeschlossen und $K\subset X$ kompakt. Sei $(x_n)_{n\in \N}$ eine Folge in
    $(\Id+F)^{-1}(K)$. Wir können wieder eine Teilfolge
    \[
        y_n=x_n+F(x_n)
    \]
    wählen, so dass $y_n\ra y$. Wie im vorhergehenden Teil folgt, folgt dass $x_n\ra x$, somit ist
    $(\Id+F)^{-1}(K)$ kompakt.
    \[ \]
\end{proof}

\subsection{Der Leray-Schauder Grad}

Sei $\Omega\subset X$, $X$ ein Banachraum. Wir setzten
\[
    \ms D_y(\ol \Omega,X)=\{ F\in \ms C(\ol\Omega,X)\, | \, y\nin(\Id + F)(\partial \Omega) \}
\]
und
\[
    \ms G_y(\ol\Omega,X):=\{F\in \ms F(\ol \Omega,X)\, |\, y\nin(\Id+F)(\partial\Omega)\}
\]
Es gilt für $F\in \ms D_y(\ol \Omega,X)$, dass $\dist(y,(\Id+F)(\partial\Omega))>0$, da $\Id+F$
abgeschlossene Mengen auf abgeschlossene Mengen abbildet.

\begin{defi}{\label{3.8}}
    Sei $\Omega\subset X$ offen, beschränkt, $y\in X$, $F\in \ms D_y(\ol \Omega,X)$. Sei $\rho:=
    \dist(y,(\Id+F)(\partial \Omega))$ und sei
    \[
        F_1\in \ms F(\ol \Omega, X),
    \]
    sodass
    \[
        \|F-F_1\|<\rho. \qq (\Ra \, F_1\in \ms G_y(\ol \Omega,X))
    \]
    Nun sei $X_1\subset X$ ein endlichdimensionaler Untervektorraum, so dass
    \[
        F_1(\Omega)\subset X_1,\, y\in X_1
    \] 
    und sei $\Omega_1=\Omega\cap X_1$. Damit ist $F_1\in \ms G_y(\ol\Omega_1,X_1)$ und wir definieren
    \[
        \deg(\Id+F,\Omega,y):=\deg(\Id+F_1,\Omega_1,y)
    \]
\end{defi}

\begin{prop}\label{3.9}
    die Obige Definiton ist unabhängig von der Wahl von $F_1$ und $X_1$.
\end{prop}

\begin{proof}
    Sei $F_2\in \ms F (\ol \Omega,X)$, $\|F_2-F\|<S$, $X_2$ entsprechend endlichdim UVR von $X$. Es sei
    $X_0=X_1+X_2$, $\Omega_0=\Omega\cap X_0$. Dann gilt
    \[
        F_i\in \ms F_y(\ol\Omega_0,X_0) \qq i=1,2
    \]
    und es folgt wegen der Reduktionseigenschaft des endlich dimenisionalen Abbildungsgrades, dass
    \[
        \deg(\Id + F_i,\Omega_0,y)=\deg(\Id+F_i,\Omega_i,y)\qq i=1,2
    \]
    Sei $H(t)=\Id + (1-t)F_1+F_2.$ Es folgt, dass
    \[
        H(t)\in D_y(\ol\Omega_0,X_0) \qq t\in [0,1],
    \]
    nachdem gilt $\|H(t)-(\Id+F)\|<\rho$ $\forall t\in [0,1]$.
    \[
        \Ra \, \deg(\Id+F_1,\Omega_0,y)=\deg(\Id+F_2,\Omega_0,y)
    \]
    Somit ergibt Definiton \ref{3.8} Sinn.
    \[ \]
\end{proof}

\begin{theorem}\label{3.10}
    Sei $\Omega\subset X$ beschränkt und offen und sei $F\in \ms D_y(\ol\Omega,X)$, $y\in X$. Dann gilt
    \begin{description}
        \item{i)}
        $\deg(\Id+F,\Omega,y)=\deg(\Id+F-y,\Omega,0)$
        \item{ii)}
        $\deg(\Id, \Omega,y)=1$, falls $y\in \Omega.$
        \item{iii)}
        Falls $\Omega_1,\Omega_2$ offene, disjunkte Teilmengen von $X$ sind, so dass
        $y\nin (\Id+F)(\ol\Omega\setminus (\Omega_1\cup\Omega_2)),$ dann gilt
        \[
            \deg(\Id+F,\Omega,y)=\deg(\Id+F,\Omega_1,y)+\deg(\Id+F,\Omega_2,y).
        \]
        \item{iv)}
        Falls $H:[0,1]\times\ol\Omega\ra X$ und $y:[0,1]\ra X$ beide stetig, so dass $H(t)\in 
        \ms D_{y(t)}(\ol\Omega,X) \qq \forall t\in [0,1]$. Dann gilt
        \[
            \deg(\Id+H(0),\Omega,y(0))=\deg(\Id+H(1), \Omega,y(1))
        \]
    \end{description}
\end{theorem}

\begin{proof}
    \textit{i)-iii)} folgen sofort aus der Definition und den entsprechenden Eigenschaften des
    endlichdimensionalen Abbildungsgrades.

    \noindent Zu \textit{iv)}: Mittels \textit{i)} betrachten wir $y(t)=0$ ($\tilde H(t)=H(t)-y(x)$).
    Nachdem $H([0,1],\partial \Omega)$ abgeschlossen und kompakt ist, gilt
    $\rho=\dist(y,H([0,1])\times\partial\Omega)>0$. Nach Satz \ref{3.6} existiert $H_1\in \ms F([0,1]
    \times\ol\Omega, X)$, so dass $\|H(t)-H_1(t)\|<\rho$ $t\in [0,1]$. Damit folgt
    \[
        \deg(\Id+H(t),\Omega,0)=\deg(\Id+H_1(t),\Omega,0) \qq \forall t\in [0,1].
    \]
    Die Behauptung folgt nun wieder aus den Eigenschaften des endlichen Abbildungsgrades.
    \[ \]
\end{proof}
   % Korrekturen, in der Vorlesung vom 23.05.11 \\ 3.3 Der Leray-Schauder Grad
\subsection{Das Leray-Schauder-Prinzip und der Schaude'sche Fixpunktsatz}

Fixpunktsätze in schneller Folge!

\begin{theorem}[Leray-Schauder-Prinzip]\label{3.12}
    Es sei $F\in \ms C(X,X)$ und es existiere $M>0$, so dass jede Lösung $x$ von
    \[
        x=tF(x) \qq t\in [0,1]
    \]
    die \textit{a priori} Abchätzung
    \[
        \|x\|\leq M
    \]
    erfüllt. Dann besitzt $F$ einen Fixpunkt.
\end{theorem}

\begin{proof}
    Sei $\rho>M$. Es gilt
    \[
        \deg(\Id-F,B_\rho(0),0)=\deg(\Id,B_\rho(0),0)=1
    \]
    Aus der \textit{a priori} Abschätzung folgt
    \[
        H(t)\in \ms D_0(B_\rho(0),X) \qq \forall t\in [0,1]
    \]
    Somit existiert eine Lösung von $x-F(x)=0$.
    \[ \]
\end{proof}

\begin{theorem}[Schauder'sche Fixpunktsatz]\label{3.13}
    Sei $K$ ein abgeschlossene, konvexe, beschränkte Teilmenge von $X$, ein Banachraum. Sei
    $F\in \ms C(K,K)$. Dann besitzt $F$ einen Fixpunkt.
\end{theorem}

\begin{proof}
    $K$ beschränkt, also sei $\rho>0$ mit $K\subset B_\rho(0)$. Es sei $R$ die stetige Fortsetzung von
    $\Id:K\ra K$ auf $X$ nach Satz \ref{2.12}. Es gilt $R:X\ra K$, $R(x)=x$ für $x\in K$.

    \noindent\textit{Anmerkung}: So eine Abbildung heißt „Retraktion“, $K$ ist dann ein Retrakt von $X$.

    \noindent Wir betrachten $\tilde F=F\circ R\in \ms C(\ol{B_\rho(0)},\ol{B_\rho(0)})$. Mit $H(t)
    -t\tilde F$ folgt
    \[
        \deg(\Id-\tilde F, B_\rho(0),0)=\deg(\Id,B_\rho(0),0)=1
    \]
    Somit existiert $x_0=\tilde F(x_0)\in K \, \Ra \, \tilde F(x_0)=F(x_0)=x_0$. 
    \[ \]
\end{proof}

\begin{remark}
    Existenz eines Fixpunktes gilt natürlich auch für $F:G\ra G$, wobei $G$ nur homöomorph zu einer
    konvexen, abgeschlossen, beschränkten Teilmenge eines Banachraumes $X$ ist.
\end{remark}

\begin{theorem}\label{3.14}
    Sei $\Omega\subset X$ offen und beschränkt, $F\in \ms C(\ol\Omega,X)$. Angenommen, es existiert ein
    $x_0\in \Omega$, so dass
    \[
        F(x)-x_0\neq \alpha \qq \forall x\in \partial \Omega, \, \forall \alpha\in(1,\infty)
    \]
    Dann besitzt $F$ einen Fixpunkt.
\end{theorem}

\begin{proof}
    Sei $H(t)(x)=x-x_0-t(F(x)-x_0)$. Dann gilt aber
    \[
        H(t)\neq 0 \qq \forall x \in \partial\Omega, \, t\in [0,1)
    \]
    Falls $H(1)(x)=0$ für ein $x\in \partial \Omega$, dann ist $x$ ein Fixpunkt von $F$ und wir
    sind fertig. Ansonsten gilt allerdings
    \[
        \deg(\Id-F,\Omega,0)=\deg(\Id-x_0,\Omega,0)=\deg(\Id,\Omega,x_0)=1
    \]
    und es existiert wieder ein Fixpunkt von $F$. 
    \[ \]
\end{proof}

\begin{cor}\label{3.15}
    Sei $\Omega\subset X$ offen und beschränkt, $F\in \ms C (\ol \Omega, X)$. Dann besitzt $F$ einen
    Fixpunkt falls \textit{eine} der folgenden Bedingungen erfüllt ist:
    \begin{description}
        \item{i)}
        $\Omega=B_\rho(0)$, $F(\partial \Omega)\subset \ol\Omega$ (Rohte)
        \item{ii)}
        $\Omega=B_\rho(0)$, $\|F(x)-x\|^2\geq \|F(x)\|^2-\|x\|^2$ $\forall x\in \partial \Omega$
        (Altmann)
        \item{iii)}
        $X$ ist ein Hilbertraum, $\Omega=B_\rho(0)$, $(F(x),x)\leq \|x\|^2 \qq \forall x 
        \in \partial \Omega$ (Krasnosel'shii).
    \end{description}
\end{cor}

\begin{proof}
    Übungsaufgabe. \[ \]
\end{proof}

\noindent \textbf{Beispiele:}

\begin{description}
    \item{1. nichtlineares 2-Punkt Randwertsproblem:}
    Wir betrachten auf $[0,T]$ das Randwertsproblem
    \begin{align}\label{11}
        \begin{split}
        u''(t)&=g(t,u(t),u'(t))\qq 0\leq t \leq T \\
        u(0)&=u(1)=0
        \end{split}
    \end{align}
    Es sei $C_0^2=\{u\in C^2\, | \, u(0)=u(T)=0\}$

    \textbf{Behauptung:} $L$ ist bistetig.
    \begin{proof}
        \begin{description}
            \item{Injektivität:} klar.
            \item{Surjektivität:} Beispielsweise mittels \textit{Green'scher Funktion}:
            \begin{align*}
                K(s,t)&=2T\cdot\sum_{n=1}^\infty \frac1{ \pi^2n^2 } \sin \frac{n\pi s}{T}
                \sin\frac{n\pi t}{T}\\
                u(s)&=L^{-1} (f)(s)=-\int_0^TK(s,t)f(t)\d t
            \end{align*}
            \item{Stetigkeit:} klar.
            \item{Stetigkeit:} \textit{Satz über die offene Abbildung}.
        \end{description}
        \[ \]
    \end{proof}
    \begin{theorem}
        Sei $g:[0,T]\times \R \times \R\ra \R$ eine stetige beschränkte Funktion. Dann bestizt
        (\ref{11}) eine Lösung $u\in C_0^2$.
    \end{theorem}
    \begin{proof}
        Es sei $G:C^1([0,1])\ra C([0,1])$ definiert durch $G(u)(t)=g(t,u(t),u'(t))$ und es sei
        \[
            J: C_0^2 \ra C^1([0,T])
        \]
        die Einbettung. Dann ist (\ref{11}) äquivalent zu
        \[
            L(u)=GJ(u) \, \LRa \, u = GJL^{-1}(u).
        \]
        Es gilt $JL^{-1}$ ist kompakt, das folgt aus \textit{Arzelá-Ascoli}.
        $G$ ist ein beschränkter Operator, denn $g$ ist stetig und beschränkt, 
        somit ist $GJL^{-1}$ kompakt und es existiert $\rho>0$ mit $GJL^{-1}(B_\rho(0))\subset 
        B_\rho(0)$. Damit besitzt $GJL^{-1}$ einen Fixpunket.
        \[ \]
    \end{proof}
   % 3.4 Das Leray-Schauder-Prinzip und der Schauder'sche Fixpunktsatz
\item[2. Existenz von Lösungen der stationären Navier-Stokes-Gleichung]
Es sei $\Omega\subset \R^3$ offen, zusammenkängend und beschränkt, und es sei $K:\Omega\ra
\R^3$ eine gegeben Funktion. Wir suchen eine Lösung $v:\Omega\ra\R^3$ zum Problem
\[
\begin{cases}
0= \eta \Delta v - v\cdot \nabla v - \nabla \rho + K & \text{ in } \Omega\\
   0= \div v & \text{ in } \Omega\\
   0= v & \text{ auf } \partial \Omega
   \end{cases}
   \]
   mit hinreichend glatten $v,\rho$.

   Eine sehr schöne Herleitung der Gleichung findet man in einem Vortrag auf der Website des Clay
   Mathematics Institute $\ra $ Link sehe Vorlesungssite.
   \end{description}

   \subsubsection*{Einschub: Sobolev-Räume}

   Wir betrachten die Menge $C^1(\Omega, \R)$ und das darauf definierte Skalarprodukt
   \[
   (u,v)_{H^1}=\int_\Omega u(x)v(x)\d x+ \int_\Omega\nabla u(x) \nabla v(x) \d x.
   \]
   Der normierte Vektorraum $(C^1(\Omega,\R),\sqrt{(\cdot,\cdot)_{H^1}})$ ist nicht vollstöndig. Die
   Vervollständigung bezeichnen wir mit $H^1(\Omega,\R)$. Die Vervollständigung von $C^1(\Omega,\R)$
   bezüglich der $H^1$-Norm bezeichnen wir mit $H_0^1(\Omega,\R)$. Die Räume $H^1$, $H^1_0$ heißen
   Sobolev-Räume und sind (separable) Hilberträume.

   Wir benötigen den folgenden wichtigen Satz aus der linearen Funktionalanalysis:

   \begin{theorem}[Rellich]\label{001}
   Die Einbettung
    \[
        H_0^1(\Omega,\R)\hookrightarrow L^2(\Omega,\R)
    \]
    ist kompakt.
    \end{theorem}

    Weiter gilt die Abschätzung

    \begin{lem}[Poincaré-Friedrichs] \label{002}
    Sei $u\in H_0^1(\Omega,\R)$. Es gilt
    \[
    \int_{\Omega}u^2\leq (d_j)^2\int_\Omega(\partial_ju)^2,
    \]
    wobei $d_j=\sup\{|x_j-y_j|\, | \, (x_1,…,x_n),(y_1,…y_n)\in \Omega\}$, bzw. gilt
    \[
    \|u\|_{L^2}\leq \frac{\diam \,  \Omega}{\sqrt{n}}\|x\|_{H^1}
    \]
    \end{lem}

    \begin{proof}
    Übungsaufgabe. \[ \]
    \end{proof}

    Wir verfeinern nun das Einbettungresultat von Rellich. Dazu benötugen wir das folgende

    \begin{lem}[Ladyzhenshaya]
    Sei $\Omega\subset\R^3$. Für $u\in H_0^1(\Omega,\R)$ gilt
    \[
    \|u\|_{L^4}\leq \sqrt[4]{8} \|u\|_{L^2}^{\nicefrac14}\|\nabla u \|_{L^2}^{\nicefrac34}
    \]    
    \end{lem}

    \begin{proof}
    Wir betrachten zuerst den Fall, dass
    \[
    u^2(x_1,x_2,x_3)=\int_{-\infty}^{x_1} \partial_1 u^2(\xi,x_1,x_3)\d \xi\leq 2\int_{-\infty}
    ^{\infty}|u(\xi,x_2,x_3)\partial_1u(\xi,x_2,x_3)|\d\xi.
    \]
    Es folgt
    \[
    \max_{x_1\in\Omega}u^2(x_1,x_2,x_3)\leq2\int_{-\infty}^\infty|u\cdot\partial_1u| \d x_1
    \]
    Nun lassen wir $x_3$ fest und integrieren über $x_1$ und $x_2$:
    \begin{align*}
    \iint u^4(x_1,x_2,x_3) &\leq \int \max_{x_1} u^2(x_1,x_2,x_3) \d x_2 \cdot \int\max_{x_2}
    u^2(x_1,x_2,x_3)\d x_1\\
        &\leq 4 \iint |u\partial_1 u|  \d x_1 \d x_2 \cdot \iint|u \partial _2 u| 
        \d x_1 \d x_2\\
        &\leq 4\left( \iint u^2 \d x_1 \d x_2\right)^{\nicefrac22}\cdot\left(\iint(\partial_1u)^2
                \d x_1\d x_2\right)^{\nicefrac 12}\cdot\left( \iint(\partial_2 u)^2 \d x_1 \d x_2
                \right)^{\nicefrac 12} \qq \text{(nach Cauchy-Schwarz)}
    \end{align*}
   Jetzt integrieren wir über $x_3$ und bekommen
   \begin{align*}
   \iiint u^4\d x_1\d x_2 \d x_3 &\leq 4 \int \d x_3 \left(\iint u ^2 \d x_1 \d x_2\right) 
    \left( \iint(\partial_1u)^2 + (\partial_2 u)^2 \d x_1 \d x_2 \right)\\
        &\leq 4\left( \iint\max_{x_3} u^2(x_1,x_2,x_3) \d x_1 \d x_2 \right)\left( \iiint 
        (\partial_1 u )^2 + (\partial _2 u)^2 \d x_1 \d x_2 \d x_3\right)\\
        &\leq 8 \iiint | u \partial_3 u | \d x_1 \d x_2 \d x_3 \cdot\left(\iiint (\partial_1u)^2+
        (\partial _2 u)^2 \d x_1 \d x_2 \d x_3\right)\\
        & \leq 8\cdot \left( \iiint u^2 \right)^{\nicefrac 12}\cdot \left( \iiint (\partial_3u)^2
        \right)^{\nicefrac12}\cdot \left( \iiint(\partial_1 u)^2 + (\partial_2u)^2 \right)^{\nicefrac 22}\\
        &\leq 8 \|u\|_{L^2}\cdot\|\nabla u \|_{L^2}^3
    \end{align*}
    Sei nun  $u\in H_0^1(\Omega,\R)$. Wir wählen $(u_j)_{j=1}^\infty$ eine Folge in $C_0^1(\Omega,\R)$,
    so dass
    \[
    u_j\ra u \begin{cases} \text{ in } H_0^1\\ \text{ in } L^2  \end{cases}.
    \]
    Dies ist wegen der Ladyzhenskaya-Ungleichung eine Cauchy-Folge in $L^4$ und konvergiert somit
    gegen $v\in L^4$. Mit Hölder gilt $u=v$ und wir können den Limes in der Ungleichung bilden.
    \[ \]
    \end{proof}

    Mit \textit{Poincaré-Friedrichs} und \textit{Ladyzhenshaya} folgt sofort, dass
    \[
    \|u\|_{L^4}\leq \left( \frac{8\diam\, \Omega}{\sqrt{n}} \right)^{\nicefrac14} \cdot\|u\|_{H^1}
    \]
    für alle $u\in H_0^1(\Omega,\R)$ und damit auch

    \begin{cor}\label{004}
    Die Einbettung
    \[
        H_0^1(\Omega,\R) \hookrightarrow L^4(\Omega,\R)
    \]
    ist kompakt.
    \end{cor}

    Zurück zu Navier-Stokes:

    Es sei 
    \[
        X=\{v\in C^2(\ol\Omega,\R^3) \, | \, \div \, v=0, v\big|_{\partial\Omega}=0\}
    \]
    und es sei
    \[
        \ms H:=\ol{X}^H_1 = \{v\in H_0^1(\Omega,\R^3) \, | \, \div \, v =0\}
    \]
    Wir betrachten $\ms H$ als Hilbertraum mit dem Skalarprodukt
    \[
        (u,v)_{\ms H}= \int_\Omega\nabla u \cdot \nabla v \d x
    \]
Dieses ist mittels \textit{Poincaré} äquivalent zum üblichen Skalarprodukt auf $H^1$:

Wir multiplizieren nun die Navier-Stokes-Gleichung \ref{12} mit $w\in X$ und erhalten

\begin{align*}
    \int_\Omega (\eta\Delta v - v\nabla v + K)w &= \int_\Omega \nabla \vp  \cdot w =0 \\
    \LRa \, \int_\Omega (\eta \underbrace{\nabla v \cdot \nabla w}_{=\sum_{k,j=1}^3 
    \partial_jv_k\partial_k w_j} -\underbrace{v\cdot v \cdot \nabla w}_{=\sum_{k,j=1}^3
    v_kv_j(\partial_kw_j)}- K\cdot w)&=0 
\end{align*}
mit
\[
    a(u,v,w):= \int_\Omega u\cdot(v\cdot \nabla w)
\]
folgt $u$ Lösung von Navier-Stokes
\begin{align}\label{13}
    \Ra \, \eta\cdot (v,w)_{\ms H} - a(v,v,w)- \int_\Omega Kw =0 \qq \forall x \in \ms H
\end{align}

Lösungen von (\ref{13}) heißen schwache Lösungen der Navier-Stokes-Gleichung.

\begin{remark}
    Es gilt
    \begin{align*}
        a(v,v,v)&=\int_\Omega v(v\cdot \nabla v)=\frac12 \int_\Omega\sum_{k,j} v_k\partial_k(v_j,v_j)
        =-\frac12 \int_\Omega\sum(v_j,v_j)\partial_kv_k=0 \qq \text{für } v\in \ms H 
    \end{align*}
\end{remark}
Für $K\in L^2 (\Omega,\R^3)$ ist $\int_\Omega K \cdot$ ein stetiges Funktional auf $\ms H$, somit
existiert nach \textit{Riesz} ein $\tilde K \in \ms H$, so dass
\[
    \int K\cdot w= (\tilde K, w)  \qq \forall w\in \ms H
\]
Dies gilt auch für $a(u,v,\cdot)$ mit $u,v\in \ms H$, somit existiert $B(u,v) \in \ms H$ mit
\[
    a(u,v,w)=(B(u,v),w)_{\ms H}.
\]
Es folgt, dass
\begin{align*}
    (\ref{13}) \, &\LRa \, (\eta v - B(v,v)- \tilde K, w)_{\ms H}=0 \qq \forall w\in \ms H\\
        &\LRa \, \eta v - B(u,v)=\tilde K.
\end{align*}
Es sei nun $Y=L^4(\Omega,\R^3)$. Dank \textit{Ladyzhenshaya-Ungleichung} ist die Einbettung
\[
    \ms H \hookrightarrow Y
\]
kompakt und es gilt mit zweimaliger Anwendung der \textit{Cauchy-Schwarz-Ungleichung}:
\[
    |a(u,v,w)|\leq\|u\|_{L^4}\|v\|_{L^4}\|w\|_{H^1}
\]


\begin{theorem}[Existenz schwacher Lösungen]\label{005}
Sei $\ms H$ ein Hilbertraum, $Y$ ein Banachraum und es sei die Einbettung $\ms H \hookrightarrow Y $
kompakt. Insbesondere sei
\[
    \|u\|_Y\leq \beta \|u\|_{\ms H} \qq \forall u \in \ms H
\]
Es sei $a: \ms H^3 \ra \R$ eine Multilinearform, so dass
\begin{align}\label{14}
    |a(u,v,w)|\leq \|u\|_Y\|v\|_Y\|w\|_{\ms H}
\end{align}
und $a(v,v,v)=0$ für alle $v\in \ms H$. Es sei $\tilde K\in \ms H$, $\eta >0$. Dann existiert
$v\in \ms H$ mit
\[
    \eta(v,w)- a(v,v,w)= (\tilde K, w)_{\ms H} \forall w \in \ms H
\]
\end{theorem}

\begin{proof}
    o.E. sei $\eta=1$. Wir suchen eine Lösung zu $v-B(v,v)+\tilde K=0$. Es gilt
    \begin{align*}
        \|B(v,v)\|_{\ms H}& \leq \alpha \|u\|_Y\|v\|_Y\leq \alpha \beta^2 \|u\|_{\ms H}\|v\|_{\ms H}
    \end{align*}
    Es sei $F(v)=B(v,v)$. $F$ ist lokal Lipschitz, denn für $\|u\|_Y< \rho,$ $\|v\|_Y<\rho$ gilt
    \begin{align*}
        \|F(u)-F(v)\|_{\ms H} &= \|B(u-v,u)-B(v,u-v)\|_{\ms H}\leq 2\alpha \rho \|u-v\|_Y\\
            &\leq 2\alpha\beta^2 \rho \|u-v\|_{\ms H}.
    \end{align*}
    Sei $(v_n)_{n\in \N}$ beschränkt in $\ms H$. Nach Wahl einer Teilfolge ist $v_n$ eine Cauchyfolge
    in $Y$. $\Ra$ $F(v_n)$ ist eine Cauchyfolge in $\ms H$, da gilt
    \begin{align*}
        \|F(u)-F(v)\|_{\ms H}\leq 2\alpha \rho\|u-v\|_Y
    \end{align*}
    Damit ist $F\in \ms C(\ms H,\ms H)$. 

    Nehmen wir nun an, dass $v$ eine Lösung ist von
    \[
        v=tF(v)-t\tilde K \qq \text{für } t\in [0,1]
    \]
    Damit gilt
    \[
        (v,v)_{\ms H} = t\underbrace{a(v,v,v)}_{=0} + t(\tilde K, v)_{\ms H} \, \Ra \,
        \|v\|_{\ms H}\leq \|\tilde K\|_{\ms H}.
    \]

    Die beiden Voraussetzungen des \textit{Leray-Schauder-Prinzips} sind damit erfüllt,
    somit existiert ein Fixpunkt $v$, der $ v=F(v)-\tilde K$ erfüllt. \[ \]
\end{proof}
   % Einschub: Sobolev-Räume \\
\section{Monotone Operatoren}

Man möchte folgendes Prinzip auf Banachräume verallgemeinern: $F:\R\ra\R$ genüge
\begin{description}
    \item{(a)}
    monoton
    \item{(b)}
    stetig
    \item{(c)}
    $F$ ist koerziv, d.h. $\lim_{v\ra \pm \infty} F(v)=\pm \infty$
\end{description}
Dann
\[
    \forall b\in \R \, \exists ! v\in \R: \, F(v)=b
\]

\begin{defi}\label{4.1}
    Sei $X$ ein reflexiver Banachraum und
    \[
        A:X\ra X^\star
    \]
    ein Operator. A heißt
    \begin{description}
        \item{(i)}
        monoton, falls für alle $u,v\in X$ gilt
        \[
            \lal Au-Av,u-v\ral\geq 0
        \]
        \item{(ii)}
        strikt monoton, falls für alle $ u, v\in X$ $u\neq v$ gilt
        \[
            \lal Au-Av,u-v\ral>0
        \]
        \item{(iii)}
        stark monoton, falls es ein $c >0$ gibt, sodass für alle $u,v\in X$ gilt
        \[
            \lal Au-Av,u-v\ral\geq c\|u-v\|^2
        \]
        \item{(iv)}
        koerziv, falls
        \[
            \lim_{\|u\|\ra \infty}\frac{ \lal Au,u\ral }{\|u\|}=\infty
        \]
    \end{description}
\end{defi}

\begin{remark}
    $A$ stark monoton $\Ra$ $A$ koerziv.
\end{remark}

\begin{proof}
    \[
        \frac{\lal Av,v\ral}{\|v\|}=\frac{ \lal Au-A(0), v\ral + \lal A(0),v\ral }{\|v\|}
        \geq \frac{c\|v\|^2 - \|A(0)\|\|v\|}{\|v\|} = c\|v\| - \|A(0)\|\ra \infty \qq (\|v\|\ra \infty)
    \]
\end{proof}

\subsubsection*{Beispiele:}

\begin{description}
    \item{1.}
    Sei $f:\R\ra\R$. Das Dualitätsproblem in $\R$ ist die Multiplikation. Im Spezialfall $X=\R=X^\star$
    stimmen die Monotoniebegriffe reeller Funktionen und für Operatoren überein, da
    \[
        \lal f(u)-f(v),u-v\ral=(f(v)-f(u))\cdot(u-v)\geq 0
    \]
    $f$ ist außerdem genau dann koerziv, d.h.
    \[
        \lim_{v\ra \pm \infty} \frac{f(u)\cdot u}{|u|}=\pm\infty \, \LRa \, f(v)\ra \pm \infty\qq(v\ra\pm
                \infty)
    \]
    \item{2.}
    $g:\R\ra \R$
    \[
        g(u):=\begin{cases} |u|^{p-2}u, & u\neq 0\\ 0, & u=0 \end{cases}
    \]
    $\Ra$ \begin{description}
        \item{(i)}
        $p>1 \, \Ra \, g$ strikt monoton.
        \item{(ii)}
        $p\geq 2 \, \Ra \, \lal g(u)-g(v),u-v\ral \geq c|u-v|^p$
        \item{(iii)}
        $p=2 \, \Ra \, g$ stark monoton
    \end{description}
    \begin{proof}
    Übungsaufgabe. \[ \]
    \end{proof}
\end{description}

\begin{defi}
    Sei $X$ reflexiver Banachraum und $A:X\ra X^\star$ ein Operator. $A$ heißt
    \begin{description}
    \item{(i)}
    \textit{demistetig} $\LRa$ $\left( u_n\ra u\, \Ra \, Au_n\rightharpoonup Au \right)$
    \item{(ii)}
    \textit{hemistetig} $\LRa \, \forall u,v,w \in X: \, t\mapsto \lal A(u+tv),w\ral$ ist stetig.
    \item{(iii)}
    \textit{stark stetig} $\LRa \, \left( u_n \rightharpoonup u\, \Ra \, Au_n\ra Au \right)$
    \item{(iv)}
    \textit{beschränkt} $\LRa \, A(V)$ beschränkt, falls $V$ beschränkt ist. 
    \end{description}
\end{defi}

\begin{remark}
    stark stetig $\Ra$ stetig $\Ra$ demistetig $\overset{(\star)}{\Ra}$ hemistetig
\end{remark}

\begin{proof}
    Zu $(\star)$. Seien $u,v,w\in X$, $t\ra t_0$ 
    \[
        \Ra\, u+tv \ra u+tv
    \]
    \[
        \overset{\text{demistetig}}{\Ra} \, A(u+tv) \rightharpoonup A(u+t_0v)
    \]
    \[
        \Ra \lal A(u+tv),w\ral\ra \lal A(u+t_0v),w\ral
    \]
\end{proof}

\begin{lem}\label{4.4}
    $X$ reflexiver Banachraum, $A:X\ra X^\star$. Dann gelten
    \begin{description}
        \item{(i)}
        $A$ stark stetig $\Ra$ $A$ kompakt.
        \item{(ii)}
        $A$ demistetig $\Ra$ $A$ lokal beschränkt.
        \item{(iii)}
        $A$ monoton $\Ra$ $A$ lokal beschränkt.
        \item{(iv)}
        $A$ monoton und hemistetig $\Ra$ $A$ demistetig.
    \end{description}
\end{lem}

\begin{proof}
    \begin{description}
    \item{(i)}
        Sei $(v_n)_{n\in \N}\subset X$ beschränkt, da $X$ reflexiv folgt:
        \[
            \exists u_{m_k} \rightharpoonup u \in X\overset{\text{stark stetig}}{\Ra} A(u_{m_k})\ra A(u)
        \]
        $\Ra$ $A$ kompakt.
    \item{(ii)}
    Wäre $A$ nicht lokal beschränkt, so existiert $v\in X$, $v_m\ra v$, sodass
    \[
        \|Av_m\|\ra \infty\, \lightning  \qq Av_m \rightharpoonup Av
    \]
    \item{(iii)}
    Angenommen $u_m\ra u: \, \|Au_m\|\ra \infty$
    \[
        a_n:=(1+\|Au_n\|\|u_n-u\|)^{-1} \qq ((a_n)_{n\in \N} \text{ beschränkt duch $1$})
    \]
    $A$ monoton
    \begin{align*}
        \Ra \, 0&\leq \lal Au_n-Av,u_n-v\ral= \lal Au_n-Av,(u_n-u)+(u-v)\ral\\
        \Ra\, a_n\lal Au_n,v-u\ral &\leq a_n(\lal Au_n,u_n-v\ral-\lal Av,u_n-v\ral)\\
        &\leq 1+ c\lal u,v\ral \qq \forall v\in X
    \end{align*}
    Einsetze $v\leadsto 2u-v$
    \begin{align*}
        -a_n \lal Au_n,v-u\ral&\leq 1+\tilde c\lal u,v\ral \qq \forall v \in X\\
        \Ra \, |\lal a_n A u_n,v-u\ral|&\leq 1+c\lal u,v\ral
    \end{align*}
    $\Ra \, (a_nAu_n)_{n\in \N}$ ist punktweise beschränkt. Mit \textit{Banach-Steinhaus}
    \[
        \sup_n\|a_nAu_n\|_{X^\star}\leq c(u)
    \]
    \[
        \Ra \, \|Au_n\|\leq \frac{c(u)}{a_n} = c(u)(1+\|Au_n\|\|u_n-v\|)
    \]
    Für
    \[
        \|u_n-u\|c(u)=\lambda <1: \, \|Au_n\|\leq \frac{c(u)}{1-\lambda}<\infty
    \]
    Ein Widerspruch zu $\|Au_n\|\ra \infty$.
    \item{(iv)}
    Gelte $u_n\ra u$. $A$ monoton $\Ra$ $(Au_n)$ beschränkt.
    $X^\star$ reflexiv
    \[
        \Ra \, Au_{n_k}\rightharpoonup b \in X^\star
    \]
    \[
        \Ra \, Au=b \qq \text{(siehe nachfolgendes Lemma)}
    \]
    Dies gilt für jede Teilfolge von $u_n, \Ra  Au_n \rightharpoonup b$.
    \end{description}
    \[ \]
\end{proof}

\begin{lem}[Minty]\label{4.5}
    Sei $X$ ein refelexiver Banachraum und $A:X\ra X^\star$, hemistetig und monoton. Dann gelten
    \begin{description}
        \item{(i)}
        $A$ maximal monoton, d.h.
        \begin{align*}
            u\in X, \, b\in X^\star, \text{ und } \lal b-Av,u-v\ral\geq 0 \qq \forall v\in X\, \Ra \,b=Au
        \end{align*}
        \item{(ii)}
        $A$ genügt der Bedingung (M), d.h.
        \begin{align}\label{15}
            Av=b \La \begin{cases} v_n \rightharpoonup v \\ Av_n\rightharpoonup b \\ \lal Av_n,v_n\ral
            \ra (b,v)
        \end{cases}
        \end{align}
        \item{(iii)}
        \[
            u_n \rightharpoonup u, \, Au_n\ra b \, \Ra \, Au=b
        \]
        \[
            u_n\ra u \,\wedge \, Au_n \rightharpoonup b \, \Ra \, Au=b
        \]
    \end{description}
\end{lem}

\begin{proof}
    \begin{description}
    \item{(i)}
    Sei $w\in X$ und $v=v-tw, \, t>0$.
    \[
        \Ra \,\lal b-A(u-tw),w\ral\geq 0
    \]
    \[
        \Ra\,\lal b-Au, w\ral\geq 0 \qq w\in X
    \]
    \[
        \Ra \,\lal b-Av,w\ral=0 \forall w
    \]
    \[
        b=Aw
    \]
    \item{(ii)}
    $A$ ist monoton.
    \begin{align*}
    \Ra \, 0&\leq \lal Au_n-Av,u_n-v\ral = \lal Au_n,u_n\ral-\lal Av,u_n\ral - \lal Au_n,Av,v\ral\\
    &\ra \lal b,v\ral-\lal Av,u\ral-\lal b-Av,v\ral
    \end{align*}
    \end{description}
\end{proof}

\begin{theorem}[Brouder, Minty]\label{4.6}
    Sei $X$ ein separabeler, reflexiver Banachraum mit Basis $(w_i)_{i\in \N}$. Sei $A:X\ra X^\star$
    monoton, koerziv und hemistetig. Dann existiert für alle $b\in X^\star$ eine Lösung $u\in X$ von
    \[ Au=b \]
    $\{A=b\}$ ist abgeschlossen, beschränkt, und konvex. Falls $A$ strikt monoton ist, ist die Lösung
    eindeutig.
\end{theorem}

\begin{idea} Der Beweis erfolgt per \textit{Galerkin-Approximation}.
\begin{description}
    \item{1.}
    $X$ separabel, d.h.
    \[
        X=\bigcup_{n=1}^\infty X^n, \qq X_n=\mr{span}(w_1,…,w_n)
    \]
    Approximiere „die Gleichung“ $Au=b$ durch „endlichdimensionale“ Probleme der Form
    \[
        (Au,z^n)=(b,z^n)    \forall z^n\in X_n
    \]
    \item{2. \textit{A-priori}-Abschätzung für Lösungen:}
    Wir zeigen, dass die Folge $(v_n)_{n\in \N}$ der Lösungen dieser Probleme beschränkt ist (mit Hilfe
            der Koerzivität).
    \item{3. Schwache Konvergenz}
    $X$ reflexiv $\Ra$ $(v_n)_{n\in \N}$ hat eine schwach konvergente Teilfolge
    \[
        u_{n_k}\rightharpoonup u \in X
    \]
    \item{4.}
    Lemma von Minty $\Ra$ Au=b.
\end{description}
\[ \]
\end{idea}
   % 4 Monotone Operatoren
\begin{proof}
    \begin{description}
    \item{(i)}
    Der Beweis erfolgt per \textit{Galerkin-Approximation}
    \[
        X_n=\mr{span}(w_1,…,w_n)
    \]
    Betrachte
    \begin{align}\label{15}
        (Au^n-b,w_k)=0 \qq \forall k=1,…,n
    \end{align}
    Definiere
    \[
        g:\R^n\ra \R^n; \qq g_k(c^n):= (A(\sum_mc_m^nw_m)-b,w_k)
    \]
    \[
        \ref{15} \, \Ra \, g(c^n)=0 \qq \text{für ein } c^n\in\R^n  
    \]
    O.B.d.A sei $|c^n|:= \|\sum_{k=1}^n c^n_kw_k\|_X$ als norm auf $\R^n$. Aus Lemma \ref{4.4} folgt,
    dass $A$ demistetig. 

    Außerdem gilt
    \[
        \sum_{k=1}^n g_k(c^n)c^n_k=(Au^n,u^n)-(b,u^n), \qq \text{wobei } u^n=\sum_kg_k(c^n)w_k^n
    \]
    $A$ ist koerziv
    \[
        \Ra \, R_0>0 \qq \forall \|u\|\geq R_0>0: \qq (Au,u)\geq (1+\|b\|)\|v\| =0
    \]
    \begin{align*}
        \Ra \, g(c^n)\cdot c^n &= \sum_{k=1}^ng_k(c^n)\cdot c_k^n \geq (1+\|b\|)\|u^n\| - \|b\|\|u^n\|
        =\|u^n\|>0 \qq \forall c^n \text{ mit } \|u^n\|\geq R_0
    \end{align*}
    Mit \textit{Brouwer} folgt (beachte: $|\cdot|_{\R^n}=\|\|_X$)
    \[
        \exists u^n\in X_n: \, g(c^n)=0,
    \]
    d.h. \ref{15} gilt für $u^n$.
    \[
        \Ra \, \|u^n\|\geq R_0
    \]
    \item{(ii)}
    $A$ monoton $\Ra$ $A$ lokal becshränkt
    \[
        \Ra \, \exists r,\delta>0: \, \|u\|\leq r \, \Ra \, \|Au\|\leq \delta
    \]
    und
    \[
        \lal Au^n-A^v,u^n-v\ral\geq 0 \qq \forall v\in X
    \]
    Wegen $u^n\in X_n$ und \ref{15} gilt
    \[
        \lal Au^n,u^n\ral=\lal b,u^n\ral\, \Ra \, |\lal Au^n,u^n\ral|=\|b\|R_0
    \]
    \begin{align*}
        \|Au^n\|&= \sup_{v\in X, \, \|v\|=r} \frac1r\lal Au^n,v\ral \overset{\text{Monotonie}}{\leq} 
        \sup_{v\in X, \|v\|=5} \frac1r (\lal Av,v\ral+\lal Au^n,u^n\ral-\lal Av,u^n\ral)
        \leq \frac1r (\delta r+ \|b\|R_0 +\delta R_0)
    \end{align*}
    $\Ra $ $(Au^n)_{n\in \N}$ ist beschränkt.
    \item{(iii)}
    Konvergenz des Galerkin-Verfahrens; $\|u^n\|\leq R_0$ und $X$ reflexiv.
    \[
        \Ra \, \exists \text{ Teilfolge } \, u^n\rightharpoonup u \,\in X
    \]
    Sei 
    \[
        w\in \bigcup_{m=1}^\infty X_m \, \Ra \, \exists n_0\in\N: \, w\in X_n\qq n\geq n_0
    \]
    \begin{align*}
        \Ra \, \forall n\geq n_0: \, \lal Au^n,w\ral&=(b,w)\\
        \Ra \, \lim_{n\ra \infty} \lal Au^n,w\ral=\lal b,w\ral \qq w\in \bigcup_{m=1}^\infty X_m
    \end{align*}
    $X^\star$ reflexiv und $(Au^n)$ beschränkt
    \begin{align*}
        \exists Au^{n_m}&\rightharpoonup c \in X^\star\\
        \Ra \, \forall w\in X: \, \lal Au^{n_k},w\ral&\ra (c,w)\\
        \Ra \, b=c, \qq \text{weil } \ol{\bigcup_{m=1}^\infty X_m}&=X \qq (X \text{ separabel})\\
        \Ra Au^{n_m} &\rightharpoonup b
    \end{align*}
    Wiederhole dieses Argument für jede Teilfolge von $Au^n$.
    \begin{align*}
        \Ra \, Au^n&\rightharpoonup b\\
        \Ra \, \lal Au^n,u^n\ral&=\lal b,u^n\ral\ra \lal b,u\ral\\
        \lim_{n\ra \infty} \lal Au^n,u^n\ral&= \lal b,u\ral
    \end{align*}
    Mit Lemma \ref{4.5} (ii) erhalten wir
    \[ 
        Au=b
    \]
    \item{(iv)}
    Eigenschaften der Lösungsmenge $S=\{A=b\}$:
    \begin{description}
    \item{(a) Beschränktheit:}
    \[
        \exists \, R_0>0 \, \forall \|V\|\geq R_0>0: \, \lal Au,u\ral\geq (1+\|b\|)\|u\|
    \]
    Sei $u$ eine Lösung mit $\|u\|\geq R$
    \begin{align*}
        \Ra \, 0&=\lal Au,u\ral -\lal b,u\ral\geq \|u\|>0 \, \lightning\\
        \Ra \, S&\subset B_{R_0}(0)
    \end{align*}
    \item{(b) Konvexität:}
    Seien $u_1,u_2\in S$, $w=tu_1+(1-t)u_2$, $v\in X$.
    \begin{align*}
        \lal b-Av,w-v\ral &=\lal b-Av,t(u_1-v)\ral+\lal b-Av,(1-t)(u_2-v)\ral\\
                 &= t\lal Au_1-Av,u_1-v\ral+(1-t)\lal Au_2-Av,u_2-v\ral\\
                 &\geq 0 \qq \forall v\in X\\
        \overset{\text{Lem. \ref{4.5}}}\Ra \, b&=Aw \, \Ra \, w\in S 
    \end{align*}
    \end{description}
    \end{description}
    \[ \]
\end{proof}

\begin{remark}
    Der Satz gilt auch, falls $X$ nicht separabel ist (ohne Beweis).
\end{remark}

\begin{cor}\label{4.7}
    Sei $X$ separabel, refl. $\R$-Banachraum und $A:X\ra X^\star$ strikt monoton, koerziv, hemstetig.
    Dann existiert
    \[
        A^{-1}:X^\star\ra X
    \]
    und $A^{-1}$ it trikt monoton und demistetig.
\end{cor}

\begin{proof}
    Übungsaufgabe. \[ \]
\end{proof}
   % 
\section{Der Nemyckii-Operator}

\begin{defi}\label{4.8}
    Sei $G\subset \R^n$, $f:G\times \R^n\ra \R^k$. Durch Anwendung auf $u:G\ra \R^n$ definieren
    wir den Operator $F$ als
    \begin{align}
        (Fu)(x)=f(x,u(x)), \qq x\in G
    \end{align}
\end{defi}

\begin{lem}\label{4.9}
    Die Funktion $f$ erfülle
    \begin{description}
        \item{1. Carathéodory-Bedingung:}
        \[
            f(\cdot,\eta): x \ra f(x,\eta) \qq \text{messbar } \forall \eta\in \R^n
        \]
        \[
            f(x,\cdot):\eta\ra f(x,\eta) \qq \text{stetig f.f.a. } x\in G
        \]
        \item{2.Wachstums-Bedingung:}
        \[
            |f^j(x,\eta)|\leq |a(x)|+ b \sum_{i=1}^n|\eta^i|^{\nicefrac{p_i}{q}}
        \]
        mit $b>0$, $a\in L^q(G)$ und $1\leq p_i$, $q<\infty$, $i=1,…,n$. 
    \end{description}
    Dann ist $F$ aus Definition \ref{4.8} ein Operator der Form
    \[
        F:\prod_{i=1}^n L^{p_i}(G) \ra (L^q(G))^k
    \]
    und $F$ ist stetig sowie beschränkt. Es gilt
    \[
        \|Fu\|_{L^q}\leq c\left( \|a\|_{L^q} + \sum_{i=1}^n \|u^i\|_{L^{p_i}}^{\nicefrac{p_i}{q}} \right)
    \]
    $F$ heißt in diesem Fall \textit{Nemyckii-Operator}.
\end{lem}

\begin{proof}
    (nur für $n=1$, $k=1$, $p=p_1$, $u=u^1$)
    \begin{description}
    \item{1. Messbarkeit von $Fu$:}
    Es sei $u\in L^p$, also messbar. Wir approximieren $u$ durch
    $(u_j)_{j=1}^n$, so dass
    \[
        u_j\ra u \text{ auf } G, \, u_j=\sum_{i=1}^{N_j} c_i^j \chi_{G_i^j}
    \]
    Es gilt
    \[
        (Fu)(x)= \lim_{j\ra \infty} f(x,u_j(x)) \text{für fast alle } x
    \]
    wegen Stetigkeit von in der 2. Variablen. Weiter gilt
    \[
        f(x,u_j(x))=\sum_{i=1}^{N_j} f(x,c_i^j) \chi_{G_i^j(x)}
    \]
    Damit ist $f(x,u_j(x))$ messbar für alle $j\in \N$ als Produkt von messbaren Funktionen.
    Somit ist $f(x,u(x))$ auch messbar als punktweiser Limes von messbaren Funktionen.

    \item{2. Beschränktheit von $F$:}
    \begin{align*}
        \|Fu\|_{L^q}^q &= \int_G|f(x,u(x))|^q= \int(|a(x)| + b |u(x)|^{\nicefrac pq})\\
            &\leq C\left( \|a(x)\|_{L^q}^q+\|u\|_{L^p} \right) \qq \text{nach Young}.
    \end{align*}
    \item{3. Stetigkeit von $F$:}
    \[u_j\ra u \qq \text{in } L^p\]
    \[
        \Ra \, \exists \text{ Teilfolge :} \, u_{j_k}\ra u \qq \text{f.ü.}
    \]
    \[
        \Ra \, (Fu_{j_k})(x) \ra (Fu)(x) \qq \text{f.f.a. } x\in G \text{ (wegen Stetigkeit)}
    \]
    Aber es gilt zudem, dass
    \[
        \int_G|f(x,u_{j_k}(x)) - f(x,u(x))|^q 
        \leq C \underbrace{(\|u\|^q_{L^q} + \|u_{j_k}\|_{L^p}^q+\|u\|_{L^p}^q)}_{\text{beschränkt.}}
    \]
    Es gilt
    \begin{align*}
        |f(x,u_{j_k}(x)) - f(x,u(x))|^q &\leq C(|f(x,u_{j_k}(x))|+ |f(x,u(x))|^q)\\
        &\leq C(|a(x)|^q + b^q|u_{j_k}|^q + |f(x,u(x))|^q) =: h_{j_k}(x)
    \end{align*}
    Es folgt
    \[
        \|Fu_{j_k}-Fu\|_{L^q}^q\leq \int_G h_{j_k}
    \]
    mit
    \[
        h_{j_k}\ra h \qq \text{in } G; \qq \int_{G}h_{j_k} \ra \int_G h,
    \]
    denn $\|u_{j_k}\|_{L^p}\ra \|u\|_{L^p}$. Jetzt folgt mit majorisierter Konvergenz, dass
    \[
        F(u_{j_k})(x)\ra (Fu)(x) \text{ in } L^p
    \]
    Die selbe Rechnung mit gleichem Limes folgt allerdings auch unter vorheriger
    Auswahl einer Teilfolge von $u_j$, somit gilt auch für die gesamte Folge, dass
    \[
        \|Fu_j-Fu\|\ra 0 \qq (j\ra \infty)
    \]
    \end{description}\[ \]
\end{proof}

\subsubsection*{Anwendung: $p$-Laplace}

Wir betrachten das Randwertproblem:
\begin{align}\label{23}
    \begin{split}
    -\div\left( |\nabla u|^{p-2} \nabla u \right) + su &=f \qq \text{auf } \Omega\\
    u&=0\qq \text{auf } \partial \Omega
    \end{split}
\end{align}
für $1<p<\infty$, $\Omega\subset\R^n$ offen, beschränkt, $s\geq 0$, $u=0$ auf $\partial\Omega$.

\noindent Schwache Formulierung: Es sei $f\in L^p (\Omega)$

Wir suchen $u\in X= W^{1,p}_0(\Omega)$, so dass
\begin{align}\label{22}
    \int_\Omega |\nabla u|^{p-2} \nabla u \nabla \vp + s\int_{\Omega}u\vp=\int f\vp \qq \forall \vp \in X
\end{align}

Wir definieren einen Operator $A$ durch
\[
    \lal Au,\vp\ral = (Au)(\vp) := \int_\Omega |\nabla u|^{p-2}\nabla u \nabla \vp +\int su\vp \qq 
    \forall u,\vp \in X
\]
und ein Funktional $b$ durch
\[
\lal b,\vp \ral = \int_\Omega f\cdot \vp
\]

\begin{lem}\label{4.10}
    Es sei $p\geq \frac{2n}{n+2}$ $f\in L^{p'}$, dann gilt
    \[
        A:X\ra X'
    \]
    beschränkt und $b\in X'$. Weiter gilt, dass (\ref{22}) äquivalent ist zu
    \[
        Au=b
    \]
\end{lem}

\begin{proof}
    Wir benutzen $\|u\|_X=\|\nabla u \|_{L^p}$ als äquivalente Norm auf $X$ mittels \textit{Poincaré}.
    \begin{description}
        \item{1)}
        Es gilt
        \begin{align*}
            |\lal Au,\vp \ral| &\leq \int_\Omega |\nabla u |^{p-1}|\nabla \vp| + s\int_\Omega |u\vp|\\
            &\overset{\text{Hölder}}{\leq} \left( \int_\Omega |\nabla u|^{(p-1)p'} \right)
            ^{\nicefrac1{p'}}\cdot \left( \int_\Omega |\nabla \vp|^p \right)^{\nicefrac 1p}
            +s \|u\|_{L^2}\|\vp\|_{L^2}\\
            &\overset{p'=\frac{p}{p-1}}{=} \|\nabla u\|_{L^p}^{p-1} \|\nabla \vp\|_{L^p} +
            s \|u\|_{L^2}\|\vp\|_{L^2}= ……
        \end{align*}
        Für $p\geq \frac{2n}{n+2}$ bettet $W^{1,p}$ genau stetig in $L^2$ ein. Es gilt somit, dass
        \[
            ……\leq C (\|\nabla u \|_{L^p}^{p-1}+s\|\nabla u\|_{L^p})\|\nabla\vp\|_{L^p}
        \]
        \[
            \Ra \, Au\in X'
        \]
        und
        \begin{align*}
            \|Au\|_{X'}&\leq \sup_{\|\vp\|_X\leq1} |\lal Au,\vp \ral| \leq C \|\nabla u\|_{L^p}^{p-1}
            =C\|\nabla u\|_X^{p-1}
        \end{align*}
        \item{2)}
        Gleichermaßen gilt
        \[
            \|b\|_{X'} \leq C\|f\|_{L^{p'}}
        \]
        \item{3)}
        Es folgt, dass
        \[
            (\ref{22}) \, \LRa \, \lal Au,\vp\ral= \lal b,\vp\ral \qq \vp \in X
        \]
        Das ist nichts anderes als
        \[
            Au=b
        \]
    \end{description}\[ \]
\end{proof}

\begin{lem}\label{4.11}
    Es seien die Voraussetzungen von Lemma \ref{4.10} gegeben, $A$ definiert nach (\ref{22}). Dann sind
    die Voraussetzungen von Theorem \ref{4.6} gegeben.
\end{lem}

\begin{proof}
    \begin{description}
    \item{1. Monotonie:}
    Es sei $g=(g^1,…,g^n):\R^n\ra \R^n$ und
    \[
        g:\xi \mapsto |\xi|^{p-2} \xi, \qq g(0)=0  
    \]
    Damit folgt, dass
    \[
        \frac{\partial g^i}{\partial \xi^j}(\xi)=|\xi|^{p-2} \delta_i^j + (p-2)|\xi|^{p-4}\xi_i\xi_j
        \qq i,j=1,…,n \qq \xi\neq 0.
    \]
    Somit gilt aber, dass
    \begin{align*}
        \sum_{i,j=1}^n \frac{\partial q^i}{\partial \xi _j}(\xi) \eta^i\eta^j&= |\xi|^{p-2}
        \left( |\eta|^2+ (p-2)\frac{(\xi \eta)^2}{|\xi|^2} \right)
        \geq \min(1,p-1) |\xi|^{p-2} |\eta|^2
    \end{align*}
    Nun sei $u\neq v \in X.$ es folgt, dass
    \begin{align*}
        \lal Au-Av,u-v\ral &= \int_\Omega (g(\nabla u)- g(\nabla v))(\nabla u - \nabla v) + s \int_\Omega
        |u-v|^2\\
            &\geq \int_\Omega \underbrace{\int_0^1 \frac{\d}{\d \tau} g(\nabla v + \tau(\nabla u -
                        \nabla v)) \d\tau\cdot(\nabla u-\nabla v)}_{:= I}  \d x\\
    \end{align*}
    Es gilt:
    \begin{align*}
        I&\geq \int_0^1 \sum_{i,j=1}^n\frac{\partial g^i}{\partial\xi_j}(\nabla v + \tau(\nabla u-\nabla v))
        \cdot (\partial_ju-\partial_jv)(\partial_iu-\partial_iv)\d \tau\\
        &\geq  c |\nabla u - \nabla v|^2\underbrace{\int_0^1 |\nabla v+\tau(\nabla u -\nabla v)|^{p-2}\d\tau}
        _{<\infty \text{ für }p>1} >0,
    \end{align*}
    denn $|\nabla v + \tau (\nabla u - \nabla v)|^{p-2}>0$ außer für max. 1 Punkt $\tau_0(x)$. Somit gilt
    aber
    \[
        \lal Au-Av, u-v \ral >0 \qq \Ra \, A \, \text{strikt monoton.}
    \]
    \item{2. Koerzitivität:}
    Es sei $u\in X$. Es gilt
    \[
        \lal Av,u\ral= \int_\Omega |\nabla u|^p + s|u|^2\geq \int_\Omega |\nabla u|^p
    \]
    \[
        \frac{\lal Au,u \ral}{\|u\|_X} \geq \|\nabla u\|^{p-1}_{L^p} \ra \infty \qq \text{für }\|u\|_X
        \ra \infty \, \text{und } p>1.
    \]
    \item{3. Stetigkeit:}
    Für die Funktion $g$ aus 1. gilt
    \[
        |g^i(\xi)|\leq c |\xi|^{p-1} = c |\xi|^{\nicefrac pq} \qq \text{für } q=\frac{p}{p-1}
    \]
    und $g$ ist stetig. Somit ist 
    \[
        F:(L^p)^n\ra (L^{p'})^n, \qq  F:\nabla u \mapsto g(\nabla u)
    \]
    ein $n$-dimenisonaler \textit{Nemyckii-Operator}. Die Stetigkeit von $A$ ist gegeben durch
    \[
        \lal Au,\vp\ral=\int_\Omega \underbrace{F(\nabla u)}_{\in L^{p'}} \cdot \underbrace{\nabla \vp}
        _{\in L^p} +\underbrace{s\int_\Omega u \vp}_{\text{beschränkt.}}
    \]
    folgt durch Normabschätzung. Damit ist $A$ stetig und somit hemistetig.
    \item{4.}
    $W_0^{1,p}$ ist separabler Banachraum, somit sind alle Voraussetzungen von \ref{4.6} erfüllt
    \end{description}
    \[ \]
\end{proof}
   % 4.1 Der Nemyckii-Operator
\newpage
Wir haben gezeigt, dass $A$ alle Voraussetzung von Satz \ref{4.6} erfüllet, somit folgt sofort:

\begin{theorem}\label{4.12}
    Sei $\Omega\subset\R^n$ offen, beschränkt, $s\geq 0$. Es sei $p\geq  \frac{2n}{n+2}$, $p>1$. Dann
    existiert zu jedem $f\in L^{p'}$ genau eine schwache Lösung des $p$-Laplace-Randwertproblems
    (\ref{23}).
\end{theorem}

\begin{proof}
    Lemma \ref{4.10}, \ref{4.11}, Satz \ref{4.6}.
    \[ \]
\end{proof}

\begin{remark}
    \begin{description}
        \item{1)}
        Der Satz \ref{4.12} gilt auch für Gleichungen der Form
        \begin{align*}
            -\div(A,(x,\nabla u))&=f \qq \text{in } \Omega\\
                        u&=0 \qq \text{auf }\partial\Omega
        \end{align*}
        falls $A$ die folgenden Bedingungen erfüllt
        \begin{itemize}
            \item \textit{Crathéodory} (Stetigkeit im 2. Argument, messbar im 1. Argument)
            \item  $|A(x,\eta)|\leq C\cdot (g(x)+ |\eta|^{p-1})$ $g\in L^{p'}(\Omega)$
                (Wachstumsbedingung)
            \item $(A(x,\eta)- A(x,\xi))\cdot (\eta-\xi)>0$ f.f.a. $x$, $\eta\neq \xi$
            \item $A(x,\eta)\cdot\eta\geq c|\eta|^p-h(x)$ $h\in L^1(\Omega)$ (Koerzitivität)
        \end{itemize}
        \item{2)}
        Man kann $f$ auch beliebig aus $(W_0^{1,p}(\Omega))'$ wählen. Das ist eine etwas größere Menge
        als $L^{p'}$.
        \item{3)} Aber was ist mit $s<0$?
    \end{description}
\end{remark}

\section{Pseudomonotone Operatoren}

In diesem Abschnitt wird die Theorie monotoner Operatoren und die Theorie kompakter Operatoren vereint.
Ein typischer pseudomonotoner Operator ist von der Form
\[
    A=A_1+A_2,
\] 
\[
    A_1: X\ra X' \qq \text{monoton, hemistetig.}
\]
\[
    A_2:X\ra X' \qq \text{strak stetig (d.h. kompakt in separablen, reflexiven Räumen $X$).}
\]

\subsubsection*{Beispiel}
\[
    \lal A_1u,\vp\ral = \int_\Omega |\nabla u|^{p-2} \nabla u \cdot \nabla \vp\qq
    \text{monoton, hemistetig.}
\]
\[
    \lal A_2u , \vp \ral =\int_\Omega su\cdot \vp\qq
    \text{kompakt als Abbildung von $X=W_0^{1,p}$ nach $X'$ für vernünftige $p$, auch für $s<0$.}
\]

\begin{defi} \label{4.13}
Wir sagen $A$ genüge Bedingung (M), falls gilt
\begin{align}\label{24}
    \left.\begin{array}{rl}
    u_n\rightharpoonup u & \text{in $X$} \\
    Au_n\rightharpoonup b & \text{in $X$} \\
    \limsup_{n\ra \infty} \lal Au_n,u_n\ral &\leq\, \lal b,u\ral
    \end{array}\right\}
    \,\Ra \, Au=b
\end{align}
\end{defi}

\begin{lem}\label{4.14}
    Es sei $X$ ein reflexiver, reeller Banachraum,
    \[
        A:X\ra X', \qq B: X\ra X'
    \]
    seien Operatoren.
    Es gilt
    \begin{description}
        \item{1)}
        $A$ monoton, hemistetig $\Ra$ $A$ genügt (M) (\ref{24})
        \item{2)}
        $A$ genügt (M) (\ref{24}), $B$ stark stetig $\Ra$ $A+B$ genügt (M) (\ref{24}).
    \end{description}
\end{lem}

\begin{proof}
    \begin{description}
        \item{1) (siehe \textit{Minty-Lemma} \ref{4.5})} Es  gilt $\forall v \in X$, dass
        \begin{align*}
            0&\leq \lal Au_n-Av, u_n-v \ral = \underbrace{\lal Au_n,u_n\ral}_{\leq \lal b,u\ral-
            \lal Au_n-Av,v\ral} - \lal Av,u_n\ral\\
            \Ra \, 0&\leq \lal b,u\ral- \lal Av,u\ral - \lal b-Av,v\ral= \lal b-Av,u-v\ral
        \end{align*}
        Sei $w \in X$, $v=u-tw$.
        \begin{align*}
            \Ra \, 0\leq \lal b-Av,u-v\ral &= \lal b-A(u-tw),tw\ral \qq t\ra 0 \qq
            \text{(Hemistetigkeit)}\\
            \Ra \, \lal b-Av, w \ral &\geq 0 \qq \forall w\\
            \lal b-Av, w \ral &\leq 0 \qq \text{(durch Einsetzen von $w$, $tw-w$ oben)}\\
            \Ra \, Av&= b .
        \end{align*}
        \item{2)}
        Es gelte $u_n \rightharpoonup u $ in $X$.
        \[
            Au_n+Bu_n \rightharpoonup b \qq \text{in }X'
        \]
        \[
            \limsup_{n\ra \infty} \lal Au_n+ Bu_n,u \ral\leq \lal b,u\ral
        \]
        Es gilt wegen starker Stetigkeit von $B$ , dass
        \[
            Bu_n\ra Bu \qq \text{in }X'
        \]
        somit folgt dass
        \begin{align*}
            \limsup _{n\ra \infty} \lal Au_n,u_n\ral & \leq \lal b-Bu,u\ral\\
            \Ra \, Au&= b - Bu\\
            \Ra \, Au + Bu &= b.
        \end{align*}
        $A$ erfüllt (M) \ref{24}.
    \end{description}
    \[ \]
\end{proof}

\begin{defi}[Pseudomonotonie] \label{4.15}
    Es sei $X$ ein reflexiver, reeller Banachraum
    \[
        A: X\ra X' \qq \text{ein Operator}
    \]
    $A$ heißt \textit{pseudomonoton}, falls gilt
    \begin{align*}
        \left.
        \begin{array}{ll}
            u_n \rightharpoonup u \qq \text{in } X\\ \limsup\lal Au_n,u_n-u\ral \leq 0
        \end{array}
        \right\}
        \, \Ra \lal Au, u-w\ral \leq \liminf \lal Au_n, u_n - w \ral \qq w\in X
    \end{align*}
\end{defi}

\begin{lem}\label{4.16}
    Sei $X$ ein reflexiver, reeller Banachraum, $A,B: X\ra X'$ zwei Operatoren
    \begin{description}
        \item{(i)}
        $A$ monoton, hemistetig $\Ra$ $A$ pseudomonoton.
        \item{(ii)}
        $A$ stark stetig $\Ra$ $A+ B$ pseudomonoton.
        \item{(iii)}
        $A, B$ beide pseudomonoton $\Ra$ $A+ B$ pseudomonoton.
        \item{(iv)}
        $A$ pseudomonoton $\Ra$ $A$ erfüllt (M) (\ref{24})
        \item{(v)}
        $A$ pseudomonoton, lokal beschränkt $\Ra$ $A$ demistetig.
    \end{description}
\end{lem}

\begin{proof}
    \begin{description}
        \item{1)}
        Sei
        \[
            u_n\rightharpoonup u \qq\text{in } X
        \]
        \[
            \limsup \lal Au_n , u_n - u \ral \leq 0.
        \]
        $A$ monoton
        \begin{align*}
            \lal Au_n- Au, u_n -u \ral & \geq 0 \\
            \Ra \, \liminf\lal Au_n , u_n - u \ral &\geq \liminf \lal Au, u_n - u \ral = 0,
        \end{align*}
        wegen schwacher Konvergenz von $u_n$.
        \[
            \lim \lal Au_n , u_n - u\ral =0
        \]
        Es sei $w\in X$, $z=u+t(w-u)$, $t>0$.
        \begin{align*}
            \Ra \, \lal Au_n - Az, u_n - z \ral &= \lal Au_n - Az , u_n -(u+t(w-u)) \ral \geq 0\\
            \Ra \, t\lal Au_n , u-w \ral &\geq \underbrace{-\lal Au_n , u_n - u\ral + \lal Az, u_n - u
                \ral}_{\ra 0} + t\lal Az, u-w\ral\\
            \Ra \, \liminf \lal Au_n , u_n - w \ral &\geq \liminf \lal Au_n, u-w \ral \geq \lal Au, v-w
            \ral
        \end{align*}
        und die Behauptung folgt.
        \item{2)}
        \begin{align*}
            u_n&\rightharpoonup u \qq \text{in } X\\
            \Ra \, Au_n&\ra Au \qq \text{$A$ stark stetig}\\
            \Ra \, \lal Au,u-w\ral &= \lim \lal Au_n ,u_n - w \ral,
        \end{align*}
        womit die Behauptung folgt.
        \item{3)}
        Es sei $u_n \rightharpoonup u$ Folge in $X$ mit
        \[
            \limsup \lal Au_n + Bu_n, u_n -u\ral \leq 0
        \]
        Wir zeigen, dass dann gilt
        \[
            \limsup\lal Au_n , u_n -u\ral \leq 0
        \]
        und
        \[
            \lal Bu_n, u_n - u \ral \leq 0
        \]
        Daraus folgt die Behauptung durch Addition der beiden Ungleichungen der Pseudomonotonität für
        $A$ und $B$ separat.
        
        \noindent Angenommen (zum Widerspruchsbeweis), dass
        \begin{align*}
            \limsup \lal Au_n , u_n - u \ral &> 0\\
            \Ra \, \text{Teilfolge} \, u_{n_k}:\, \lim \lal Au_{n_k}, u_{n_k}-u\ral &=a\\
            \Ra \, \limsup \lal Bu_{n_k},u_{n_k-u}&= \limsup \lal (B+A)u_{n_k} - A u_{n_k}, u_{n_k}-u\ral
            \\
            &\leq \underbrace{\limsup \lal (B+A) u_{n_k}, u_{n_k}-u\ral}_{\leq 0} +\underbrace{
            \lim \lal -Au_{n_k}, u_{n_k} - u\ral}_{=-a} \leq -a\\
            \Ra \, 0= \lal Bu,u-u\ral &\leq \liminf \lal Bu_{n_k} , u_{n_k}-u\ral\\
            &\leq \limsup \lal Bu_{n_k}, u_{n_k}-u\ral \leq -a <0 \qq \lightning,
        \end{align*}
        womit die Behauptung folgt.
        \item{4)}
        Es sei $u_n\rightharpoonup u$ Folge in $X$, so dass
        \[
            Au_n \rightharpoonup b
        \]
        und
        \begin{align*}
            \limsup \lal Au_n,u_n\ral &\leq \lal b,u\ral\\
            \Ra \,  \limsup \lal Au_n, u_n - u\ral &\leq \limsup \lal Au_n, u_n\ral - \underbrace{
            \lim\lal Au_n, u\ral }_{=\lal b,u\ral }\\
            & \leq \lal b,u\ral - \lal b,u\ral = 0.
        \end{align*}
        Damit gilt wegen der Pseudomonotonietät für alle $w\in X$:
        \begin{align*}
            \lal Au, u-w\ral &\leq \liminf \lal Au_n, u_n - w\ral \leq \lal b,u \ral - \lal b,w\ral
            = \lal b, u-w\ral
        \end{align*}
        mit $w'=2u-w$ folgt die umgekehrte Ungleichung, somit gilt
        \[
            \lal Au, u-w \ral = \lal b , u-w \ral \qq w\in X \, \Ra \, Au=b 
        \]
        und damit die Behauptung.
        \item{5)}
        Es gelte $u_n \ra u$ in $X$.
        \[
            \Ra \, (Au_n)_{n\in \N} \text{ beschränkt in $X'$, da $A$ lokal bechränkt.}
        \]
        \[
            \Ra \, \exists \text{ Teilfolge:} \, Au_{n_k}\rightharpoonup b
        \]
        \[
            \Ra \, \lim \lal Au_{n_k},u_{n_k}-u\ral =0
        \]
        Aus der Pseudomonotonität folgt
        \begin{align*}
            \lal Au, u-w \ral &\leq \liminf \lal Au_{n_k} , u_{n_k} - w\ral \qq \forall w \in X\\
            &= \lal b ,  u-w \ral \qq \text{$u_n$ stark-schwach konvergent.}\\
            \overset{\text{wie in 4)}}{\Ra}  Au&= b.\\
            \Ra \, Au_{n_k}&\rightharpoonup Au.
        \end{align*}
        Dieses Argument gilt auch bei vorheriger Auswahl einer weiteren Teilfolge, somit folgt die
        Behauptung, da $A$ demistetig.
    \end{description}
    \[ \]
\end{proof}
   % 4.2 Psoudomonotone Operatoren
\begin{theorem}[Brezis]\label{4.16}
    Es sei $A:X\ra X'$ ein pseudomonotoner, beschränkter koerzitiver Operator,
    $X$ ein separapbler, reflexiver reeller Banachraum. Dann existiert für alle $b\in X'$
    eine Lösung $u\in X$ von
    \[
        Au=b
    \]
\end{theorem}

\begin{proof}
    Nach Lemma \ref{4.16} ist $A$ demistetig und genüge Bedingung (M). Es sei $(w_i)_{i\in \N}$
    eine abzählbare Teilmenge von $X$, so dass $\mr{span}\{w_i\}_{i\in \N}$ dicht liegt in $X$ und
    $w_n$ linear unabhängig von $\{w_j\}_{j=1}^{n-1}$ für alle $n\in\N$ ist.

    Wie im Beweis von \textit{Browder-Minty} suchen wir approximative Lösungen $u^n\in \mr{span}\{w_j\}
    _{j=1}^n=:X^n$
    \[
        u_n=\sum_{j=1}^n c_j^nw_j
    \]
    das „Galerkin-Systems“
    \begin{align}\label{25}
        g_k(c^n)=g_k(u_n):=\lal Au_n - b, w_k\ral =0\qq \forall k=1,…,n.
    \end{align}
    Auf $\R^n$ betrachten wir die Norm $\|c^n\|_X=\|u_n\|_X$ mit
    \[
        u_n=\sum_{j=1}^nc_j^nw_j
    \]
    Es gilt
    \begin{description}
    \item{-} $g_k$ ($k=1,…,n$) stetig in $\R^n$, da $A$ demistetig
    \[
        \Ra \, \lal Au^j-b,w_n\ral\ra \lal Au -b,w_k\ral
    \]
    \item{-} $\exists R_0>0:$ $\sum_{k=1}^n g_k(\tilde c^n)\cdot \tilde c_k^n>0$,
    falls $\|\tilde u_n\|=R_0$ mit $\tilde u_n=\sum_{j=1}^n\tilde c_j^nw_j$. (Koerzitivität von $A$) 
    \end{description}
    Aus den Eigenschaften des \textit{Brouwer'schen Abbildungsgrades} folgt die Existenz einer Lösung
    $u_n\in X_n$ von \ref{25} für alle $n\in \N$ mit $\|u_n\|_X\leq R_0$.
    (Das war Theorem \ref{2.10}: Ein stetiges Vektorfeld, das auf einer Kugeloberfläche in jdm. Punkt
     nach außen zeigt muss an einem Punkt im Innern verschwinden.)
    Wegen der Beschränktheit von $(u_n)_{n\in\N}$ Lösungen von (\ref{25}) existiert eine (schwach)
    konvergente Teilfolge $u_{n_j}\rightharpoonup u$ in $X$. Wir zeigen, dass $u$ die gesuchte Lösung
    ist. Wegen Beschränktheit von $A$ und der schwachen Konvergenz von $u_{n_j}$ existiert eine weitere
    Teilfolge (wieder mit $n_j$ indiziert), so dass
    \[
        Au_{n_j}\rightharpoonup c \qq \text{in }X
    \]
    Weiter gilt aber
    \[
        \lal Au_{n+k}, w \ral = \lal b,w \ral \qq \forall w\in \bigcup_{j=1}^n X_j \qq k\geq 0.
    \]
    \begin{align*}
        \Ra \, \lim_{n\ra \infty} \lal Au_n, w\ral = \lal b,w\ral \qq \forall w\in \bigcup_{j=1}^\infty
        X_j
    \end{align*}
    Es folgt wie in \textit{Browder-Minty}, dass $c=b$, denn für alle $w\in \bigcup_{j=1}^\infty X_j$
    gilt
    \[
        \lal c-b , w\ral =0
    \]
    und $\bigcup_{j=1}^\infty X_j$ liegt dicht in $X$.

    \noindent Somit folgt aber, dass
    \[
        \lal Au_{n_j}, u_{n_j} \ral = \lal b,u_{n_j}\ral \ra \lal b,u\ral
    \]
    Von oben wissen wir noch dass $u_n\rightharpoonup u$ und $Au_{n_j}\rightharpoonup b$. Aufgrund von
    (M) folgt sofort, dass
    \[
        Au=b
    \]
\end{proof}

\subsubsection*{Anwendung: $p$-Laplace (Redux)}

Wir suchen $u:\Omega\ra \R$, so dass
\begin{align*}
    -\div (|\nabla u |^{p-2}\nabla u) + g(u)&= f \qq \text{in }\Omega,\\
     u&=0 \qq\text{auf }\partial \Omega,
\end{align*}
wobei $\Omega\subset\R^n$ offen, beschränkt, $p>1$. Wir betrachten wieder $x=W^{1,p}_0(\Omega)$ und die
Operatoren $A_1$, $A_2$ definiert durch
\begin{align*}
    \lal A_1u,\vp\ral &:= \int_{\Omega} |\nabla u |^{p-2} \nabla u \cdot \nabla\vp\\
    \lal A_2u,\vp\ral &:= \int_\Omega g(u)\cdot \vp
\end{align*}
und
\[
    \lal b,\vp \ral := \int_\Omega f\vp.
\]
Der Operator $A_1$ und das Funktional $b$ wurden bereits in $p$-Laplace Teil 1 behandelt. Zu $A_2$:

\begin{lem}\label{4.17}
    Es sei $\Omega\subset \R^n$ beschränkt, offen, $g:\R\ra \R$ stetig, so dass
    \[
        |g(y)|\leq c \cdot(1+|y|^{r-1})
    \]
    für  $1<r<\infty$.
    Falls $1\leq p < n$ und $r\leq \frac{np}{n-p}$, dann ist $A_2:X\ra X'$ beschränkt. Für
    $r<\frac{np}{n-p}$ ist $A_2$ stark stetig. Falls $p\geq n$, so gilt das für alle $r\in (1,\infty)$.
\end{lem}

\begin{proof}
    siehe Übungsblatt, folgt aus Sobolev-Einbettung.\[ \]
\end{proof}

\noindent Koezitivität: Achtung, beispielsweise für $g(y)=-sy$, $s$ hinreichend groß ist $A=A_1+A_2$
unter Umständen nicht koerzitiv. Es existiert dann auch nicht immer eine Lösung.

\begin{beispiel}
\[
    \Delta u - \lambda u=0
\]
besitzt nichttriviale Lösungen $u$ für $\lambda$ Eigenwert von $\Delta$.
$\mapsto$ Fredholm-Alternative: Nichtexistenz von Lösungen von $\Delta u - \lambda u = \alpha u$
$\alpha>0$.

Wir rechnen
\begin{align*}
    \frac{\lal A_1u+A_2u,u\ral}{\|u\|_X} &= \frac{\|\nabla u\|_{L^p}^p+\int g(u)u}{\|\nabla u\|_{L^p}}\\
    &\geq \|\nabla u\|_{L^p}^p - {\|A_2u\|_{X'}}\\
    \text{mit} \qq \|A_2u\|_{X'} &= \sup_{\vp \in X} \lal A_2u, \vp \ral \geq \frac{\lal A_2 u , u\ral}
    {\|\nabla u\|_{L^p}}
\end{align*}
\end{beispiel}

\begin{beispiel}
\[
    g(u)=-su, \qq s>0
\]
Es gilt $\|u\|_{L^2}\leq C\| \nabla u\|_{L^p}$ mit $C$ der Sobolev-Konstante.
\[
    \Ra \, \frac{-\int su^2}{\|\nabla u\|_{L^p}} \geq -s C^2 \cdot\|\nabla u\|_{L^p}
\]
und es folgt die Koerzitivität von $A=A_2+A_1$ für $p>2$ und jedes $s>0$, (dann gewinnt $\|\nabla u \|
_{L^p}^{p-1}$ immer) oder in der linearen Gleichung ($p=2$) für $s\cdot C^2<1$.
\end{beispiel}

\begin{beispiel}
Falls
\begin{align}\label{26}
    \inf_{y\in \R} g(y) y> -\infty,
\end{align}
dann ist $A=A_1+A_2$ immer koerzitiv, falls $p>1$.
\end{beispiel}

\begin{theorem}\label{4.18}
    Sei $\Omega \subset\R^n$ offen, beschränkt, $p>1$. Die Funktion $g$ erfülle die Voraussetzungen von
    Lemma \ref{4.17}. Weiter sei der Operator $A=A_1+A_2$ koerzitiv, also sei z.B. (\ref{26}) erfüllt.
    Dann existiert für alle $f\in L^{p'}(\Omega)$ ein $u\in W_0^{1,p}(\Omega)$, so dass gilt
    \[
        (A_1+A_2)u=b
    \]
    ($u$ ist schwache Lösung von 
     \[
        \div(|\nabla u|^{p-2}\nabla u)+ g(u)=f
     \]
     mit Nullrandbedingung.)
\end{theorem}

\begin{proof}
    Folgt aus dem \textit{Satz von Brezis} \ref{4.16}. \[ \]
\end{proof}

\subsection{Evolutionsprobleme}

Es wird notwendig sein, Funktionen der Art
\[
    f:S\ra X
\]
mit $S\subset \R$ und $X$ ein reflexiver reeller Banachraum und insbesondere deren Integrale zu
betrachten. Deshalb ein 

\subsubsection*{Einschub: Das Bochner-Integral}

\begin{defi}[Treppenfunktion]\label{4.19}
    Eine Funktion $f:S\ra X$ heißt \textit{Treppenfunktion}, falls gilt
    \[
        f(s)=\sum_{j=1}^n \chi_{B_j} (s) x_j
    \]
    mit $x_j\in X$, $j=1,…,n$; $B_j\subset S \subset \R$ Lebesgue-messbare Mengen mit
    $|B_i|\leq \infty$ und $B_i\cap B_j=\varnothing$ für $i\neq j$ Maß von $B_i$. $\chi_{B_i}$ ist
    die charakteristische Funktion von $B_i$.    
\end{defi}

\begin{defi}[Bochner-Integral]\label{4.20}
    Für eine Treppenfunktion $f$ definieren wir das \textit{Bochner-Integral}
    \[
        \int_S f(s)\d s:= \sum_{j=1}^n|B_j|\cdot f(x_j)
    \]
\end{defi}

\begin{remark}
    Es gilt natürlich
    \[
        \int_S f(s)\d s \in X
    \]
\end{remark}
\begin{remark}
    Verallgemeinerung auf andere Maßräume (z.B. Teilmengen des $\R^n$) ebenfalls möglich.
\end{remark}

   % 4.3 Evolutionsprobleme \\ Einschub: Das Bochner-Integral
\begin{defi}\label{4.22}
    Eine Funktion $f:S\ra X$ heißt \textit{Bochner-Messbar}, falls eine Folge $(f_n)_{n\in\N}$
    von Treppenfunktionen $f_n:S\ra X$ existiert, so dass für fast alle $s\in S$ gilt
    \[
        \lim_{n\ra \infty} \|f_n(s)-f(s)\| _X=0.
    \]
    Falls $f$ in eine solche Folge zudem gilt
    \[
        \lim_{n\ra \infty} \int_S\|f_n(s)-f(s)\|_X \d s =0,
    \]
    so heißt $f$ \textit{Bochner-Integrierbar} und schreiben
    \begin{align}\label{26}
        \int_S f(s)\d s=\lim_{n\ra\infty} \int_S f_n(s)\d s
    \end{align}
\end{defi}

\begin{prop}\label{4.23}
    \begin{description}
    \item{1)}
    Falls $f:S\ra X$ Bochner-messbar ist, so ist die Funktion $\|f(\cdot)\|:S\ra \R$ Lebesgue-messbar.
    \item{2)}
    Der Grenzwert in (\ref{26}) existiert für Bochner-integrable Funktionen $f$.
    \item{3)}
    Der Grenzwert in (\ref{26}) ist unabhängig von der Wahl der Folge $(f_n)_{n_\in\N}$. 
    \end{description}
\end{prop}

\begin{proof}
    \begin{description}
        \item{1)}
        Es sei $f$ Bochner-messbar, $(f_n)_{n\in \N}$ eine Folge von Treppenfunktionen, so dass
        $\lim_{n\ra \infty}\|f_n(s)-f(s)\|_X=0$ ffa $s\in S$. Dann ist $\|f_n(\cdot)\|:S\ra \R$ eine
        (reelle) Treppenfunktion. Es gilt
        \[
            \lim_{n\ra\infty} \|f_n(s)\| =\|f(s)\| \qq \text{ffa } s\in S 
        \]
        Damit folgt, dass $\|f(\cdot)\|$ Lebesgue-messbar ist als punktweiser Limes messbarer Funktionen.
        \item{2)}
        Es gilt für $n,k\in\N$ gilt dass
        \begin{align*}
            \left\|\int_Sf_n\d s - \int_Sf_k(s)\d s\right\|_X &=\left\| \int_S f_n(s)-f_k(s)\d s
                \right\| _X\leq \int_S\|f_n (s)- f_k(s)\|_X\d s\\ 
            &\leq \int_S\|f_n-f\|_X + \|f_k-f\|_X \d s \ra 0 \qq (n,k\ra \infty)
        \end{align*}
        Damit ist $\int_S f_n(s)\d s$ eine Cauchy-Folge und besitzt einen Grenzwert im Banachraum $X$.
        \item{3)}
        Es sei $(\tilde f _n)_{n\in\N}$ eine andere Folge von Treppenfunktionen, für die gilt
        \[
            \lim_{n\ra \infty} \int_S\|\tilde f_n(s)-f(s)\|_Y \d s=0
        \]
        Wir wählen
        \[
            \hat f_n= \begin{cases} \tilde f_{\nicefrac n2} & \text{$n$ gerade.}\\
                      f_{\nicefrac{(n+1)}{2}} & \text{$n$ ungerade.} \end{cases}
        \]
        Es gilt $\lim_{n\ra \infty}\int_S \hat f_n(s)$ existiert wegen 2) somit auch für TF. Somit ist
        gezeigt dass das Bochner-Integral für Bochner-integrierbare Funktionen definiert ist.

    \end{description}
    \[ \]
\end{proof}

\begin{theorem}[Pettis]\label{4.24}
    Es sei $X$ ein separabler Banachraum. Dann ist $f:S\ra X$ genau dann Bochner-messbar, wenn für alle
    $\vp \in X'$ die Funktion
    \[
        \lal \vp,f(\cdot)\ral_X:S\ra \R
    \]
    Lebesgue-messbar ist.
\end{theorem}

\begin{proof}
    \begin{description}
        \item{1)}
        Sei $f$ Bochner-messbar, $(f_n)_{n\in\N}$ die entsprechende Folge von Treppenfunktionen, sei
        $\vp \in X'$. Es folgt f.f.a. $s\in S$, dass
        \[
            \underbrace{\lal \vp, f_n(s)\ral}_{\text{Treppenfunktionen in $\R$}} \ra \lal \vp , f(s)\ral
        \]
        Somit ist $\lal \vp, f(\cdot)\ral$ als punktweiser Limes von reellwertiger Treppenfunktionen 
        Lebesgue-messbar.
        \item{2)}
        Wir zeigen zunächst dass $\|f(\cdot)\|_X: S\ra \R$ Lebesgue-messbar ist.

        \begin{lem}\label{27}
            Sei $X$ ein separabler Banachraum. Dann existiert $(\vp_n)_{n\in \N}$ eine Folge in $X'$ mit
            $\|\vp _n\|_{X'}\leq 1$, so dass $\forall \vp_0\in X'$, $\|\vp_0\|\leq 1$ eine Teilfolge
            $(\vp_{n_k})_{k\in\N}$ existiert, mit $\lim_{n\ra \infty}\vp_{n_k}(f)=\vp_0(f)$ für alle 
            $f\in X$.
        \end{lem}
        (Lemma nötig, weil $X'$ nicht unbedingt separabel ist nicht nötig für reflexiven Fall)

        \noindent Es sei $a\in \R$ und definieren
        \[
            A:=\{s: \, \|f(s)\|_X\leq a\}, 
        \]
        sowie
        \[
            A_\vp := \{s: \, |\vp(f(s))|\leq a  \} \qq \text{für } \vp \in X'
        \]
        Falls gilt $A=\bigcap_{j=1}^\infty A_{\vp_j}$, wobei $\vp _j$ die Folge aus dem Lemma oben ist,
        so ist $A$ als abzählbare Schnittmenge messbarer Mengen messbar und a) bewiesen. Es gilt
        natürlich, dass
        \[
            A\subset\bigcap_{\|\vp\|_{X'}\leq 1} A_{\vp},
        \]
        aber wir wissen auch, dass für festes $s$ ein $\vp_0\in X'$ existiert mit $\|\vp_0\|_{X'}=1$ und
        $\vp_0(f(s))=\|f(s)\|_X$ nach \textit{Hahn-Banach}. Damit gilt auch
        \[
            A\supset \bigcap_{\|\vp\|_{X'}\leq 1} A_\vp,
        \]
        somit folgt
        \[
            A=\bigcap _{\|\vp\|_{X'}=1} A_\vp
        \]
        Mit dem Lemma oben folgt aber dass gilt
        \[
            \bigcap _{\|\vp\|_{X'}\leq 1}A_\vp = \bigcap _{j=1}^\infty A_{\vp_j} \qq \text{($\vp_j$ aus
                    dem Lemma)}
        \]
        somit folgt $A=\bigcap_{j=1}^\infty A_{\vp_j}$ und a) ist bewiesen.
        \item{b)}
        Nachdem $X$ separabel ist, können wir für alle $n\in \N$ die Menge $X$ mit eine abzählbaren Menge
        offener Kugeln $\{S_{j,n}\}_{j\in\N}$,
        \[
            S_{j,n} = B_{\nicefrac1n} (f_{j,n})
        \]
        mit Radius $\frac1n$ überdecken. Nach Teil a) gilt, dass
        \[
            \|f(\cdot)- f_{j,n}\|:S\ra \R
        \]
        Lebesgue-Messbar ist für alle $j,n\in\N$. Damit ist aber auch die Menge
        \[
            B_{j,n}:=\{s\in S: \, f(s)\in S_{j,n}\}
        \]
        Lebesgue-messbar. Es gilt $S=\bigcup_{j=1}^\infty B_{j,n}$. Wir setzen 
        \[
            f_n(s)= f_{j,n}\qq \text{für } s\in B_{j,n}':= B_{j,n}\setminus \bigcup_{i=1}^{j-1} B_{i,n}.
        \]
        Es folgt, dass $\|f(s)- f_n(s)\|_X\leq \frac1n$ für alle $s\in S$. Nachdem $B_{j,n}'$
        Lebesgue-messbar ist für alle $j,n\in\N$ ist $f_n(\cdot)$ aber Bochner-messbar, denn auf dem
        endlichen Maßraum kann man $f_n$ durch endliche Treppenfunktionen approximieren, und damit ist
        auch $f$ Bochner-messbar.
    \end{description}
    \[\]
\end{proof}

\begin{remark}
    Zur Approximation $f_n$ abzählbarer Treppenfunktionen durch endliche Treppenfunktionen:
    Wir nehmen für festes $n\in \N$
    \[
        (f_{n,k})_{k\in\N} := \left(\sum_{j=1}^k f_{j,n} \cdot \chi_{B_{j,n}'}\right)
    \]
    und für alle $s\in S$ gilt:
    \[
        f_{n,k}(s) \ra f_n(s) \qq (k\ra \infty)
    \]
    Damit ist $f_n$ Bochner-messbar.
\end{remark}

Nun beweisen wir das Lemma \ref{27}:
\begin{proof}
    es sei $\{f_n\}_{n\in\N}$ dicht in $X$. Wir betrachten
    \[
        \eta_n\cdot \vp \ra \eta_n(\vp)= \begin{pmatrix} \vp(f_1)& \vp(f_2)& \cdots & \vp(f_n)\end{pmatrix}.
    \]
    für $\vp \in S'=\{\vp\in X': \, \|\vp\|_{X'}=1\}$. Der Zielraum ist endlichdimensional,
    also separabel. Somit existiert für festes $n$ eine Folge $\{\vp_{n,k}\}_{k\in\N}$ in $S'$, so dass
    $\{\eta_n(\vp_{n,k})\}_{k\in \N}$ dicht liegt in im Bild $\eta_n(S')$ von $S'$. Es folgt nun,
    dass für jedes $\vp _0\in S'$ eine Teilfolge $(\vp_{n,m_n})_{n\in\N}$ existiert, so dass
    \[
        |\vp_{n,m_n}(f_i)-\vp_0(f_i)|< \frac1n \qq \text{für alle } i\in \{1,…,n\}.
    \]
    damit gilt aber, dass
    \[
        \lim_{n\ra \infty} \vp_{n,m_n}(f_i) = \vp_0(f_i) \qq \forall i \in \N.
    \]
    Wegen der Dichtheit der $\{f_i\}_{i\in\N}$ und der gleichmäßigen Beschränktheit der $\vp_{n,k}$ gilt
    \[
        \lim_{n\ra \infty} \vp_{n,m_n}(f)= \vp _0(f) \qq \forall f\in X.
    \]
\end{proof}

\begin{cor}\label{4.25}
    Sei $X$ ein separabler Banachraum, $f:S\ra X$. Es seien $f_n: S\ra \N$, $n\in\N$ Bochner-messbare
    Funktionen, so dass f.f.a. $s\in S$ gilt
    \[
        f_n(s)\rightharpoonup f(s)
    \]
    Dann ist $f$ Bochner-messbar.
\end{cor}

\begin{proof}
    Sei $F\in X'$. Es gilt $\lal F,f_n(\cdot)\ral$ ist Lebesgue messbar, wegen Theorem \ref{4.24}
    ($\Ra$). und es gilt, dass
    \[
        \lal F,f_n (\cdot)\ral \ra \lal F,f(\cdot)\ral \qq \text{f.ü. auf } S
    \]
    Damit ist $\lal F,f(\cdot)\ral$ Lebesgue-messbar, und wegen Theorem \ref{4.24} ($\Ra$) ist $f$
    Bochner-messbar.
    \[ \]
\end{proof}
   %
\begin{theorem}\label{4.26}
    Eine Bochner-messbare Funktion $f: S\ra X$ ist genau dann Bochner-integrierbar, wenn die Funktion
    \[
        \|f(\cdot)\|_X :S \ra \R
    \]
    Lebesgue-integrierbar ist.
\end{theorem}

\begin{proof}
    \begin{description}
    \item{1)}
    Sei $f:S\ra X$ Bochner-integrierbar, $(f_n)_{n\in \N}$ ein Folge von Treppenfunktionen, so dass
    f.f.a. $s\in S$ gilt $f_n(s)\ra f(s)$ in $X$. Es folgt sofort, dass $\|f(\cdot)\|$ Lebegue-messbar
     ist. Punktweise gilt
     \[
        \|f(s)\|\leq \|f_n(s)\|+ \|f(s)-f_n(s)\|
     \]
     und somit
     \[
        \int_S\|f(s)\|\d s\leq \underbrace{\int_S\|f_n(s)\|}_{\text{Treppenfunktion}}
            + \underbrace{\int_S\|f(s)-f_n(s)\|<\infty}_{\text{Grenzwert existiert}}
     \]
    $\Ra$ $\|f(\cdot)\|$ Lebesgue-integrierbar.
    \item{2)}
    Sei $f$ Bochner-messbar und sei $\|f(\cdot)\|:S \ra \R$ Lebesgue-integrierbar. Es existiert eine
    Folge von Treppenfunktionen, so dass
    \[
        f_n(s) \ra f(s) \qq \text{f.f.a. } s\in S .
    \]
    Es sei
    \[
        g_n(s):= \begin{cases} f_n(s)& \text{falls } \|f_n(s)\|\leq \nicefrac32 \|f(s)\| \\
                 0&\text{falls } \|f_n(s)\|>\nicefrac32 \|f(s)\|\end{cases}
    \]
    Es gilt immer noch: $g_n$ Treppenfunktion und
    \[
        g_n(s)\ra f(s) \qq \text{f.f.a. } s\in S
    \]
    Gleichzeitig gilt aber, dass
    \[
        \|g_n(s)-f(s)\|\leq \|g_n(s)\|+\|f(s)\| \leq \frac52\|f(s)\|
    \]
    Mit dem \textit{Lebesgue'schen Konvergenzsatz} folgt, dass
    \[
        \lim_{n\ra\infty} \int_S \|g_n(s)-f(s)\|=0,
    \]
    womit $f$ Bochner-integrierbar ist. 
    \end{description}    
    \[ \]
\end{proof}

\begin{cor}\label{4.27}
    Es sei $f:S\ra X$ Bochner-integrierbar. Dann gilt
    \[
        \left\| \int_S f(s)\d s \right\|_X\leq\int_S\|f(s)\|_X \d s \qq  \text{und } \forall \vp \in X'
        \text{ gilt} \qq\lal \vp, \int f(s)\d s \ral _X = \int_S \lal \vp, f(s)\ral_X \d s
    \]
\end{cor}

\begin{proof}
    \begin{description}
    \item{1)}
    $(f_n)_{n\in \N}$ Folge von Treppenfunktionen, die $f$ punktweise f.ü. in $S$ approximieren.
    Es folgt, dass
    \begin{align*}
        \left\| \int_Sf(s)\d s\right\| &=\lim_{n\ra\infty} \left\| \int_Sf_n(s) \d s\right\|
        \leq \lim_{n\ra \infty} \int_S\|f_n(s)\| \d s\\
        &\leq \lim_{n\ra \infty} \int_S \|f_n(s)-f(s)\| +\|f(s)\|\d s = \int_S \|f(s)\| \d s.
    \end{align*}
    \item{2)}
    Wir können wie im Beweis von Theorem \ref{4.26} annehmen, dass
    \[
        \|f_n(s)\|_X \leq \frac32 \|f(s)\|.
    \]
    Damit gilt für $\vp \in X'$, dass
    \[
        \lal \vp, \int_S f(s)\d s\ral = \lim_{n\ra \infty} \lal \vp , \int_S f_n(s)\d s\ral 
        \underset{\text{endl. Summe}}{=} \lim_{n\ra \infty} \int_S \lal \vp, f_n(s)\ral \d s
        \underset{\text{Lebesgue}}{=} \int_S \lal \vp, f(s)\ral \d s
    \]
    \end{description}
    \[ \]
\end{proof}

\begin{remark}
    Es sei $I$ ein beschränktes Intervall, $f\in C(\ol I, X)$. Man kann in diesem Fall zeigen, dass
    Bochner- und Riemann-Integral von $f$ über $I$ übereinstimmen.
\end{remark}

\subsubsection*{$L^p$-Räume mit Werten in Banachräumen}

Es ist weiterhin $X$ ein separabler Banachraum, $S\subset \R$ offen.

\begin{defi}\label{4.28}
    Wir bezeichnen mit $L^p(S;X)$ $1\leq p<\infty$ die Menge aller Bochner-messbaren Funktionen
    $f:S\ra X$, für die gilt
    \[
        \int_S \|f(s)\|^p \d s < \infty
    \]
    Die Menge aller Bochner-messbaren Funktionen, für die ein $M>0$ existiert mit $\|f(s)\|_X\leq M$ 
    f.f.a. $s\in S$ bezeichnen wir mit $L^\infty(S;X)$.
\end{defi}

\begin{theorem}\label{4.29}
    Die Menge $L^p(S;X)$ $1\leq p \leq \infty$ bildet einen Banachraum bzgl. der Norm
    \[
        \|f\|_{L^p(S;X)} = \left( \int_S \|f(s)\|_X^p\d s \right)^{\nicefrac1p} \qq \text{für } 1\leq p 
        <\infty
    \]
    bzw. 
    \[
        \|f\|_{L^\infty(S;X)}=\esssup_{s\in S} \|f(s)\|_X
    \]
\end{theorem}

\begin{proof}
    Die Eigenschaften der Norm sind evident.
    \begin{description}
    \item{Vollständigkeit:} Sei $(f_n)_{n\in\N}$ eine Cauchy-Folge in $L^p(S;X)$, also
    \[
        \|f_n(s)-f_k(s)\|_{L^p(S;X)} = \left( \int_S \|f_k(s)-f_n(s)\|_X^p\right)^{\nicefrac1p} \ra 0
        \qq (n,k\ra \infty)
    \]
    Mit der Dreiecksungleichung in $X$ und dem Satz von Bochner folgt, dass $\|f_n(\cdot)\|_X$ eine
    Cauchy-Folge in $L^p(S;\R)$ ist, der Rest des Beweises folgt dann von reellen $L^p$-Räumen.
    \end{description}
    \[ \]
\end{proof}

\begin{lem}\label{4.30}
    Die Menge der Treppenfunktionen ist dicht in $L^p(S;X)$ $1\leq p <\infty$.
\end{lem}

\begin{proof}
    $p=1$: Siehe Beweis des Satzes von Bochner.\\
    $p>1$: Analog zu reellwertigen $L^p$. \[ \]
\end{proof}

\begin{theorem}[Hölder-Ungleichung]\label{4.31}
    Sei $f\in L^p(S;X)$, $g\in L^{p'}(S;X')$ ($X'$ Dualraum von $X$, $\nicefrac1p+\nicefrac1{p'}=1$, 
    $1\leq p \leq \infty$). Dann ist
    \[
        \lal g(\cdot), f(\cdot)\ral _X\in L^1(S;\R)
    \]
    und es gilt
    \[
        \left| \int_S\lal g(s),f(s)\ral_X\d s \right| \leq \|g\|_{L^{p'}(S;X')}\|f\|_{L^p(S;X)}
    \]
\end{theorem}

\begin{proof}
    Übungsaufgabe. \[ \]
\end{proof}

\begin{theorem}[Darstellungssatz]\label{4.32}
    Es sei $X$ ein separabler, reflexiver Banachraum, $1<p<\infty$. Dann besitzt jedes Funktional
    \[
        \vp \in \left( L^p(S;X)\right)'
    \]
    eine eindeutige Darstellung der Form:
    \[
        \vp (u) = \int_S\lal v(s),u(s)\ral_X \d s \qq \forall u\in L^p(S;X)
    \]
    mit $v\in L^{p'}(S;X')$ ($\nicefrac1p+\nicefrac1{p'}=1$).
\end{theorem}

\begin{proof}
    \textit{4 Schritte:}
    \begin{enumerate}
        \item % 1)
        Sei $t_0\in S$. Für $t\in S$, $x\in X$ setzen wir
        \begin{align*}
            u_{t,x}(s) &:= \begin{cases} x & \text{falls $t_0\leq t$ und $t_0\leq s \leq t$}\\
                        -x & \text{falls $t_0>t$ und $t\leq s \leq t_0$} \\
                        0 & \text{sonst}.\end{cases}\\
            \Ra \, u_{t,x} &\in L^p (S;X)
        \end{align*}
        Es sei $\vp\in (L^p(S;X))'$. Damit ist $\vp(u_{t,x})$ linear in $X$ und es gilt
        \[
            |\vp(u_{t,x})|\leq \|\vp\|_{(L^p(S;X))'} \cdot\|u_{t,x}\|_{L^p(S;X)} \leq \|\vp\|
            _{(L^p(S;X))'} \cdot \|x\|\cdot (t-t_0)^{\nicefrac1p}.
        \]
        $\vp(u_{t,x})$ ist also ein stetiges lineares Funktional auf $X$. Somit existiert für jedes
        $t\in S$ eine Darstellung
        \[
            \vp(u_{t,x})= \lal g(t),x\ral _X \qq \text{mit } g(t)\in X'
        \]
        Es gilt $g(t_0)=0$.
        \item % 2)
        Es seien $\{S_1,…,S_m\}$ nichtleere, disjunkte Teilintervalle der Form $S_L=[t_i,t_i+h_i]\subset
        S$ $i=1,…,m$. Für $x\in X$ gilt dann
        \begin{align*}
            \left|\sum_{i=1}^n\lal g(t_i+h_i),x\ral-\lal g(t_i),x\ral\right|&= \left| \vp
            \left( \sum_{i=1}^n (u_{t_i+h_i,x}-u_{t_i,x}) \right) \right|\\
                &\leq \|\vp\|_{L^{p'}(S;X')}\cdot\|x\|\cdot\left(\sum_{i=1}^m h_i\right)^{\nicefrac1p}.
        \end{align*}
        Wir setzen $z_i:= \frac{1}{n_i}(g(t_i+h_i)-g(t_i))$. Es existieren (wegen der Reflexivität von
        $X$), $y_i\in X$ mit $\|y_i\|_X=1$, $\lal z_i,y_i\ral = \|z_i\|_{X'}$. Wir wählen
        $x_i=\|z_i\|^{p'-1} y_i$ und es folgt
        \[
            \lal z_i,x_i \ral = \|z_i\|^{p'}_{X'}
        \]
        Nun sei
        \[
            u(t):= \begin{cases} x_i & \text{für $t\in S_i$, $i=1,…,m$} \\
                   0 & \text{für } t\in S\setminus\bigcup_{i=1}^m S_i\end{cases}
        \]
        Dann gilt
        \[
            u= \sum_{i=1}\left( u_{t_i+h_i,x_i}-u_{t_i,x_i} \right)
        \]
        Es folgt, dass
        \[
            \|u\|^p_{L^p(S;X)} = \sum_{n=1}^m \|x_i\|^p_X h_i = \sum_{i=1}^m \|z_i\|_{X'}^{p'} h _i
            \qq (p(p'-1)=p')
        \]
        Somit folgt
        \begin{align*}
            \|\vp\|_{(L^p(S;X))'} \cdot \|u\|_{L^p(S;X)} &\geq \vp (u) 
                = \sum_{i=1}^m \lal g(t_i+h_i)- g(t_i), x_i\ral = \sum_{i=1}^m \lal z_i,x_i\ral h_i
                \geq \sum_{i=1}^m \|z_i\|_{X'}^{p'} \cdot h_i
        \end{align*}
        Also gilt
        \[
            \|\vp\|_{(L^p(S;X))'} \geq \left( \sum_{i=1}^m \|z_i\|_{X'}^{p'} h_i \right)^{\nicefrac1{p'}}
        \]
        und es folgt, dass
        \begin{align}\label{28}
            \|\vp\|_{(L^p(S;X))'} &\geq \sup_{\text{Intervallteilung}} \left( \sum_{i=1}^m
            \underbrace{\left\|\frac1{n_i}(g(t_i+h_i)-g(t_i))\right\|_X^{p'}}_{z_i}\cdot h_i\right)
            ^{\nicefrac1{p'}}
        \end{align}
   % Einschub: L^p-Räume mit Werten in Banachräumen
\item % 3)
Wir suchen $v\in L^{p'}(S;X')$, so dass
\[
    g(t)=\int_{t_0}^t v(\tau) \d \tau
\]
Wir können $S$ bis auf eine Nullmenge als Vereinigung abgeschlossener Intervalle darstellen, somit genügt
es, die Aussage für $S=[t_0,t_1]$ zu beweisen. Es sei $T=t_1-t_0$ und wir setzen
\[
    g_{n,i} = g(\underbrace{t_0+2^{-n}T}_{(**)}i) \qq\text{für $i=0,1,…,2^n$ $n\in\N$},
\]
und
\[
    u_n(t)=\frac1{2^{-n}T} (g_{n,i+1}-g_{n,i})\qq \text{für $2^{-n}Ti\leq t-t_0<2^{-n}T(i+1)$}
\]
Es gilt wegen (\ref{28}), dass $2^n-1$
\begin{align}\label{29}
    \|v_n\|_{L^{p'}(S;X')} = \left( \sum_{i=0}^{2^n-1}\left\| \frac1{2^{-n}T}(g_{n,i+1}-g_{n,i}) \right\|
    ^{p'}_{X'}\cdot 2^{-n}T\right)^{\nicefrac1{p'}} \overset{(\ref{28})}{\leq} \|\vp\|_{(L^p(S,X))'}
\end{align}
Wir zeigen nun, dass f.f.a. $t\in S$ die Folge
\[
    (\|u_n(t)\|_{X'})_{n\in\N}
\]
beschränkt bleibt. Es sei $N_0:=\{t\in S: \, \sup_{n}\|u_n(t)\|_{X'}=\infty \}$ und
\[
    S_{r,K}:= \{t\in S: \, \sup_{n\leq r} \|u_n(t)\|_{X'}\geq K\} \qq r,K\in \N.
\]
Es gilt $S_{r,K}\subset S_{p,K}$ für $p\geq r$ und
\[
    N_0 = \bigcap_{K=1}^\infty \bigcup_{r=1}^\infty S_{r,K}
\]
Nachdem $u_n$ sthw. konstant ist, besteht $S_{r,k}$ aus disjunkten Intervallen der For $[t_i,t_i+n_i)$ 
$i=1, … , m$. Es gilt
\[
    K^{p'}|S_{r,K}|\leq \sum_{i=1}^m \left( \left\| \frac{g(t_i+h_i)-g(t_i)}{h_i}\right\|^{p'}_{X} h_i
            \right) \overset{(\ref{28})}{\leq} \|\vp\|
\]
Nachdem $S_{r,K}$ in $r$ wächst folgt
\begin{align*}
    K^{p'} \left|\bigcup_rS_{r,K}\right| &\leq \|\vp\|_{L^p(S;X)'}\\
    \Ra \, |N_0|&=\left| \bigcap_K\bigcup_r S_{r,K} \right| \leq \inf_{K}\left|\bigcup _r S_{r,K} \right|
    =0
\end{align*}
Nun sei $\{z_j\}_{j\in\N}$ dicht in $X$ und wir definieren
\[
    a_j(t)=\Lal g(t), z_j\Ral \qq t\in S
\]
\[
    b_{j,u}(t) =\Lal u_n(t), z_j \Ral \qq t\in S 
\]
Es gilt $b_{j,n}\in L^{p'}(S;\R)$ mit gleichmäßiger beschränkter Norm.
\[
    \|b_j^n\|_{L^{p'}(S;\R)}^{p'}=\int_S|\Lal u_n(\tau), z_j\Ral|^{p'}\leq \int_S \|u_n\|_X^{p'} \cdot
    \|z_j\|_X^{p'} \leq \|z_j\|_X^{p'} \|u_n\|_{L^{p'}(S;X')}
\]
Wir können somit sukzessive eine Teilfolge auswählen, so dass für alle $j\in \N$ gilt
\[
    b_j^n\rightharpoonup b_j \in L^{p'}(S;\R).
\]
\textbf{Behauptung 1:} $\forall t\in S$ gilt
\[
    \int_{t_0}^t b_j(\tau) \d \tau= a_j (t).
\]
\begin{proof}
    Es sei $\{\tilde S_n\}_{n\in\N}= \{t_{n,i}\}_{n\in\N}$, $i\in \{0, …, 2^n\}$, die Menge aller Punkte,
    welche durch die Intevallteilung ($**$) erreicht werden. Offensichtlich liegt $\{\tilde S_n\}
    _{n\in \N}$ dicht in $S$. Nun sei $\tilde t\in S_n$, damit gilt $\tilde t\in S_m$ $\forall m\geq n.$
    Es gilt für $m\geq n$
    \begin{align*}
        \int_{t_0}^{\tilde t} b_{m,j}(\tau) \d \tau &= \int_{t_0}^{\tilde t} \Lal u_m,z_j\Ral 
        = \sum_{i=1}^{2^m\frac{\tilde t}T} \frac1{2^{-m}T}\Lal g_{m,i+1} - g_{m,i}, z_j\Ral \cdot 2^{-m}T
        \\
        &= \Lal g_{m, 2^m\frac{\tilde t}T}, z_j\Ral = \Lal g(t_0+ \tilde t), z_j \Ral = a_j (\tilde t)
    \end{align*}
    Aber wegen $b_n^j\rightharpoonup b_j$ weil $\chi_{(0,\tilde t)}$ eine Testfunktion in $L^p(S;\R)$
    ist, gilt
    \[
        \int_{t_0}^{\tilde t} b_j(\tau) \d \tau= a_j(\tau)
    \]
    Mit Stetigkeit von $a_j(t)$ und von $\int_{t_0}^t b_j(\tau)\d\tau$ folgt die Behauptung.\[ \] 
\end{proof}
\noindent\textbf{Behauptung 2:} Es gilt in diesem Fall (d.h. $a_j(t)=\int_{t_0}^tb_j(\tau)\d\tau$ f.f.a
$t\in S$, und $b_j\in L^1([t_0,t_1);\R)$) dass $a_j$ fast überall differenzierbar ist und
\[
    b_j(t)=a'_j(t) \qq \forall t\in S\setminus N_j, \, |N_j|=0.
\]
\begin{proof}
    Übungsaufgabe. \[ \]
\end{proof}

Nun sei $t\in S\setminus \bigcup_{j=0}^\infty N_j$ und $i_n\in \N$ sei für $n\in \N$ so gewählt,
dass gilt $2^{-n}Ti_n\leq t-t_0<2^{-n}T(i_n+1)$. Dann folgt:
\[
    \lim_{n\ra \infty} b_{j,n}= \lim_{n\ra \infty} \Lal u_n(t),z_j\Ral = \lim_{n\ra \infty}
    \Lal \frac1{2^{-m}T} (g_{n,i_n+1}-g_{n,i_n}), z_j \Ral = a'_j(t).
\]
Damit folgt aber, dass $\forall t\in S\setminus \bigcup_{j=0}^\infty N_j$ gilt
\[
    u_n(t) \rightharpoonup v(t)  \qq  \text{in } X'. 
\]
Für $t\in \underbrace{\bigcup_{i=0}^\infty N_j}_{\text{Nullmenge}}$ setzen wir $v(t)=0$. Die so
definierte Funktion $v:S\ra X'$ ist Bochner-messbar (als Limes von Treppenfunktionen) und es gilt für
$t\in S$, dass
\[
    \|u(t)\|_{W'}\leq \liminf_{n\ra \infty} \|v_n(t)\|_{X'}
\]
Daraus und aus (\ref{29}) folgt mit Fatou, dass
\begin{align}\label{30}
    \|u\|_{L^{p'}(S;X')}^{p'} \leq \liminf\int_S\|u_n(t)\|_{X'}^{p'} \leq \|\vp\|_{(L^p(S;X))'}
\end{align}
Also ist $u\in L^{p'}(S;X')$. Wegen Behauptung 1 gilt für alle $j\in \N$, dass
\[
    \Lal g(t),z_j\Ral = a_j(t) = \int_{t_0}^t a_j'(\tau)\d\tau=\int_{t_0}^{t} \lim_{n\ra \infty}
    \Lal u_n(\tau) , z_j \Ral \d\tau= \int_{t_0}^t \Lal u(\tau), z_j\Ral \d\tau
    = \Lal \int_{t_0}^t u_(\tau) \d\tau, z_j\Ral.
\]
Also ist
\[
    g(t)= \int_{t_0}^tu(\tau) \d\tau
\]
\item % 4)
Es sei nun $u$ eine Treppenfunktion aus $L^p(S;X)$ auf einer endlichen Menge von Intevallen, also
\[
    u(t)=\begin{cases} x_i & \text{für } s_i<t\leq t_i, \, i=1,…,m \\
         0& \text{sonst}\end{cases}
\]
Damit folgt sofort
\begin{align*}
    \vp(u)&=\vp\left(\sum_{i=1}^m (u_{t_i,x_i}- u_{s_i,x_i} )\right)
    = \sum_{i=1}^m \Lal g(t_i)-g(s_i),x_i\Ral\\
    &= \sum_{i=1}^m \Lal \int_{s_i}^{t_i} u(t)\d t, x_i\Ral = \sum_{i=1}^m \int_{s_i}^{t_i}
    \Lal v(t) , u(t)\Ral
\end{align*}
Nachdem diese Treppenfunktionen dicht liegen in $L^p(S;X)$ und dank \textit{Hölder'scher Ungleichung}
die Abbildung
\[
    u\mapsto \int_S\Lal v(t),u(t)\Ral \d t
\]
stetig ist bzgl. $u\in L^p(S;X)$ gilt
\[
    \vp (u)=\int_S \Lal v(t),u(t)\Ral \d t \qq \forall u\in L^p(S;X)
\]
Wegen \textit{Hölder} gilt ebenfalls, dass
\[
    \|\vp\|_{(L^p(S;X))'} \leq \|v\|_{L^{p'}(S;X')}.
\]
Mit (\ref{30}) folgt also
\[
    \|\vp\|_{(L^p(S;X))'} = \|v\|_{L^{p'}(S;X')}.
\]
Eindeutigkeit der Darstellung folgt dann daraus, dass wegen Normgleichheit das Nullfunktional $\vp=0$
nur durch die Nullfunktion $v=0$ dargestellt werden kann.
\end{enumerate}
\[ \]
\end{proof}

\begin{remark}
    Die Zuordnung $T:\vp \mapsto v$.
    \[
        T:(L^p(S;X))'\ra L^{p'}(S;X')
    \]
    ist eine lineare, surjektive Isometrie.
\end{remark}

\begin{remark}
    Mit einigen kleinen Änderungen funktioniert der Beweis auch für $X$ nur separabel. Mit weiteren
    kleinen Änderungen funktioniert der Beweis auch für
    \[
        \vp \in (L^1(S;X))'.
    \]
    Die Funktion $v$ ist dann in $L^\infty(S;X)$.
\end{remark}

\begin{remark}
    Der Satz gilt auch für $X$  nur reflexiv. Der Beweis ist dann aufwendiger. Im Allgemeinen gilt der
    Satz, wenn $X$ die sog. \textit{Radon-Nilkodyn Eigenschaft} besitzt.
\end{remark}
   %
\begin{remark}
    $C_c^\infty((0,T);X)$ liegen dicht in $L^p((0,T); X)$ für $1\leq p < \infty$. Siehe Übungsaufgabe.
\end{remark}

\subsubsection*{Geltaud-Tripel}
Wir betrachten als Beispiel zunächst die Wäremleitungsgleichung auf $I\times \Omega$, $I=(0,T)$,
$\Omega\subset \R^n$ offen, beschränkt. Sei $u$ also ein eglatte Lösung von
\begin{align*}
    \partial_t u - \Delta u &= f\qq \text{in } I\times \Omega\\
    u&= 0 \qq \text{auf } I\times \partial \Omega\\
    u(0)&=u_0 \qq \text{in } \Omega
\end{align*}
Die Funktion $u$ erfüllt dann auch die schwache Formulierung, d.h. für alle $\vp \in L^2(F; W_0^{1,2}
(\Omega))$ gilt:
\begin{align}\label{31}
    \iint_{I,\Omega} \partial_t u \vp + \iint_{I,\Omega} \nabla u \cdot \nabla \vp = \iint_{I,\Omega}
    f\cdot \vp
\end{align}
Mit der Wahl $\vp=u$ erhalten wir
\begin{align*}
    \iint_{I,\Omega} \frac{\partial}{\partial t}\frac{u^2}2 + \underbrace{\|u\|
    _{L^2(I;W_0^{1,2}(\Omega))}^2}_{=\iint_{I\Omega |\nabla u |^2}} &\leq \|f\|_{L^2(I;L^2(\Omega))}
    \cdot \|u\|_{L^2(I;L^2(\Omega))}\\
    &\leq \frac 12 \|f\|^2_{L^2(I;L^2(\Omega))} + \frac12 \|u\|^2_{L^2(I;L^2(\Omega))}\\
    \Ra \, \int_I\frac{\d}{\d t} \frac{\|u\|_{L^2}^2} + \|U\|^2_{L^2(I,W_0^{1,2}(\Omega))}
    &\leq \frac 12 \left( \|f\|^2_{L^2(I;L^2(\Omega))} + \|u\|^2_{L^2(I,L^2(\Omega))}\right)
\end{align*}
\begin{align}\label{32}
    \Ra \, \|u\|^2_{L^\infty(I,L^2(\Omega))}+\|u\|_{L^2(I;W_0^{1,2}(\Omega))} &\leq \left( \|f\|^2
    _{L^2(I;L^2(\Omega))} + \|u_0\|^2_{L^2(I;L^2(\Omega))} \right)
\end{align}

Mit (\ref{31}) folgt aber auch, dass
\[
    \iint_{I,\Omega} u_t\cdot \vp = \int_I -(u,\vp)_{W_0^{1,2}} + \int_I(\vp,\vp)_{L^2} 
\]
Das heißt, dass $u_t$ kanonisch mit einem Element von $(W_0^{1,2}(\Omega))'$ identifiziert werden kann.
Die Abschätzung ergibt mit Hilfe von (\ref{32}), dass
\[
    \|u_t\|_{L^2(I;W_0^{1,2}(\Omega)')} \leq \tilde c (\|f\|^2_{L^2(I;L^2(\Omega))} + \|u_0\|^2
    _{L^2(I;L^2(\Omega))}   )  
\]
Wir schreiben damit dass $u_t=\Delta u + f \in (W^{1,2}_0(\Omega))'=:W^{-1,2}_0(\Omega)$
(punktweise f.ü.). Aber: $u_t$ ist als Funktion \textit{nicht} durch Anfangswert und $f$ in
$W_0^{1,2}(\Omega)$ beschränkt.
Es ist in diesem Fall günstig, die sog. \textit{Gelfand-Tripel} einzuführen.

Es sei $V$ ein reflexiver, separabler Banachraum, dicht und stetig eingebettet in einen Hilbertraum $H$.
Damit ist durch jedes stetige Funktional $f\in H'$ ein stetiges Funktional auf $V$ definiert:
\[
    \lal f,v \ral _V := (f,v)_H
\]
Damit gilt $H'\hookrightarrow V'$ stetig. Mittels dem \textit{Riesz'schen Darstellungssatz} identifizieren wir nun
$H'$ mit $H$ und erhalten
\[
    V \hookrightarrow H \hookrightarrow V'
\]
stetig.

\noindent$\leadsto$ Definiert einen Operator $T:H'\ra V'$, einfach als Einschränkung der Funktionale aus
$H'$. Es gilt:
\begin{description}
    \item{i)} $\|T\vp\|_{V'}\leq C\|\vp\|_{H'}$
    \item{ii)} $T$ ist injektiv.
    \item{iii)} $T(H)$ liegt dicht in $V'$, da $V$ reflexiv ist.
\end{description}
Die Skalarprodukte $\lal \cdot,\cdot \ral_V$ und $(\cdot,\cdot)_H$ sitmmen überein, falls beide Sinn
ergeben, d.h.
\[
    \lal f,v\ral_V=(f,v)_H \qq \forall f\in H, \forall v\in V,
\]
wobei wir von nun an $Tf$ mit $f$ identifizieren. Falls nun $V$ ebenfalls ein Hilbertraum ist, ergibt es
keinen Sinn, $V$ mit $V'$ zu identifizieren. Dazu ein einfaches

\begin{beispiel}
    Es sei
    \[
        H=l^2= \left\{ u=(u_n)_{n\in\N} : \, \sum_{i=1}^\infty a_i<\infty\right\}
    \]
    mit Skalarprodukt
    \[
        (u,v)_H=\sum_{n=1}^\infty u_nv_n \qq \forall u,v\in\mc l^2
    \]
    und sei
    \[
        V:=\left\{u=(u_n)_{n\in\N}:\, \sum_{n=1}^\infty n^2u_n^2<\infty\right\}
    \]
    mit Skalarprodukt
    \[
        ((u,v))_V:= \sum_{n=1}^\infty n^2u_nv_n.
    \]
    Es gilt $V\hookrightarrow H$ stetig, dicht. Wir identifizieren, $H$ und $H'$, und identifizieren
    damit aber $V'$ mit dem Raum 
    \[
        V'=\left\{ f=(f_n)_{n\in\N}: \, \sum_{n=1}^\infty \frac1{n^2} f_n^2 <\infty\right\}
    \]
    Das Dualitätsprodukt $\lal\cdot,\cdot\ral_V$ ist damit gegeben durch
    \[
        \lal f,v\ral_V=\sum_{n=1}^\infty f_nv_n.
    \]
    Der Isomorphismus $T:V\ra V'$ ist damit gegeben durch die Abbildung
    \[
        u=(u_n)_{n\in\N}\mapsto Tu=(n^2u_n)_{n\in\N}
    \]
\end{beispiel}

\subsubsection*{Die verallgemeinerte Zeitableitung und ihre Eigenschaften}
($\ra$ Zeidler, Nonlinear Funktional Analysis and its Applications II A)\\[.5cm]
\noindent Wir erinnern uns: $C^1(I;X)$, $I=[0,T]$ ist der Raum der stetigen Fréchet-differenzierbaren
Abbildungen $u:I\ra X$, $X$ ein Banachraum. Die Ableitung $u': I\ra \ms L(\R,X)$ identifizieren wir für
jedes $t\in I$ mit einem Element in $X$.


\begin{defi}[Verallgemeinerte Zeitableitung]\label{4.34}
    Es sei $u\in L^p(I,V)$, $1<p<\infty$ und $(V,H,V')$ ein Gelfand-Triplet. Eine Funktion
    \[
        \frac{\d u}{\d t} \in L^{p'} (I;V')
    \]
    heißt \textit{verallgemeinerte Zeitableitung} von $u$, falls gilt
    \[
        \int_0^T\lal \frac{\d u}{\d t} (t), v\ral _V \vp (t)\d t = -\int_0^T(u(t),v)_H \vp'(t)\d t
        \qq \forall v\in V,\, \forall \vp \in C_c^\infty (I,\R).
    \]
\end{defi}

\begin{remark}
    Im Allgemeinen ist die verallgemeinerte Zeitableitung nicht mit der schwachen Zeitableitung einer
    Funktion $u:I\times \Omega\ra \R$ überein. Identität gilt jedoch durchaus, falls 
    $C^\infty_c(\Omega,\R)$ dicht in $V$ liegt.
\end{remark}

\begin{remark}
    Für $u\in C^1([0,T], X)$ stimmen starke und verallgemeinerte Zeitableitung überein. 
\end{remark}

\noindent \textbf{Notation:} Im Folgenden bezeichnen wir mit $W$ den Raum:
\[
    W:=\left\{u\in L^p(I,V): \, \frac{\d u}{\d t}\in L^{p'}(I; V')\right\}
\]
mit
\[
    \|u\|_W:= \|u\|_{L^p(I;V)} + \left\|\frac{\d u}{\d t}\right\|_{L^{p'}(I;V)}
\]

\begin{prop}\label{4.34}
    $W$ ist ein Banachraum.
\end{prop}

\begin{proof}
    Alle Eigenschaften außer der Vollständigkeit sind klar. Diese folgt mit:

    \begin{prop}\label{4.35}
        Es sei $(V,H,V')$ ein Gelfand-Tripel und
        \[
            \frac{\d}{\d t} u_k = v_k \qq \forall k\in \N
        \]
        mit
        \begin{align*}
            u_k&\rightharpoonup u \qq L^p(I;V) \qq (k\ra \infty)\\
            v_k&\rightharpoonup v \qq L^{p'}(I;V') \qq (k\ra \infty)
        \end{align*}
        Dann gilt
        \[
            \frac{\d }{\d t} u =v,
        \]
        wobei $\frac{\d}{\d t} u$ die verallgemeinerte Zeitableitung bezeichne.
    \end{prop}
    \begin{proof}
        Falls $A$, $B$ Banachräume, $A\hookrightarrow B$ stetig, so gilt 
        \[
            a_k\rightharpoonup a \text{ in }  A \qq \Ra \qq a_k\rightharpoonup a \text{ in } B.
        \]
        Es folgt in unserem Fall, dass
        \[
            u_k\rightharpoonup u, \qq v_k\rightharpoonup v \qq \text{in }L^1(I;V').
        \]
        Offensichtlich ist aber $\vp \cdot v$, $\vp \in C_c^\infty (I;\R)$ $v\in V$ eine geeignete
        Testfunktion in $L^1(I,V')$. Im Limes folgt die Behauptung.\[ \]
    \end{proof}
    \[ \]
\end{proof}
   %
\begin{prop}\label{4.36}
    Der Raum $C^1(\ol I; V)$ liegt dicht in $W$.
\end{prop}

\begin{proof}
    Folgt wie globale Approximation von Sobolev-Funktionen. Es sei 
    \[
        \tilde u(t)=\begin{cases} u(t) & t\in I \\ u(-t) & t\in(-\tau,0) \\ u(T-t) & t\in (T,T+\tau) \\
        0 & \text{sonst} \end{cases}
    \]
    und es sei $\eta_n(t)$ eine Standard-Diracfolge. Wir betrachten
    \[
        u_n(t)=\int_\R \eta_n(t-s) \tilde u(s) \d s.
    \]
    Es gilt $u_n(t)\in C^1(\ol I;V)$ und es gilt
    \[
        \frac{\d}{\d t} \tilde u_n = \int_\R \eta'_n(t-s) \tilde u(s) \d s
        =\int_\R \eta_n(t-s) \frac{\d}{\d t} u(s) \d s \qq \forall t\in I
    \]
    Die Konvergenz von $u_n\ra u$ in $L^p(I;V)$ und von $\frac{\d}{\d t}u_n:=v_n\ra v$ in $L^{p'}(I;V')$
    folgt wie im skalarwertigen Fall. Mit Benutzung von Proposition \ref{4.35} folgt $v=\frac{\d}{\d t}u$
    und somit die Behauptung. \[ \]
\end{proof}

\begin{prop}\label{4.37}
    Es gilt
    \[
        W\hookrightarrow C(\ol I;V') 
    \]
    stetig.
\end{prop}

\begin{proof}
    Es sei $u\in W$. Damit ist $\frac{\d}{\d t}u \in L^{p'}(I;V')$ eine lokal integrierbare Funktion.
    Wir definieren
    \[
        v(t)=\int_0^t\frac{\d}{\d t} u(\tau)\d \tau,
    \]
    mit $\|v(t)-v(s)\|_{V'}\leq \int_0^t \|\frac{\d}{\d t} u(\tau)\|_{V'}$.
    Somit ist $v(t)\in C(\ol I; V')$. Wie im skalaren Fall folgt, dass
    \[
        v(t)=u(t)+c \qq \text{mit} \, c\in V' \, \text{f.f.a.} \, t\in I.
    \]
    Damit ist auch $u$ stetig. Mit Hölder folgt
    \begin{align*}
        \|v(t)\|&\leq T^{\nicefrac1p}\left\|\frac{\d}{\d t}u\right\|_{L^{p'}(I;V')} \\
        \|c\|_{V'}&\leq T^{-\nicefrac{1}{p}}\|u-v\|_{L^p(I;V)}
    \end{align*}
    Damit gilt aber
    \[
        \max \|u(t)\|_{V'}\leq C\cdot \left( \|u'\|_{L^{p'}(I;V)} + \|u-v\|_{L^{p'}(I;V)}\right)
        \leq \tilde C \|u\|_W
    \]
\end{proof}

Es folgt das wichtige Einbettungslemma

\begin{lem}\label{4.38}
    Es sei $(V,H,V')$ ein Gelfand-Tripel, $W$ der Funktionenraum wie oben. Dann gilt
    \[
        W\hookrightarrow C(\ol C; H)
    \]
    und für alle $u,v\in W$, $s.t \in \ol I$ gilt
    \begin{align}\label{33}
        \int_s^t \lal  \frac{\d u}{\d t}(\tau), v(t)\ral_V+ \lal \frac{\d v}{\d t} (\tau) , u(\tau)\ral_V
        \d\tau =(u(t),v(t))_H - (u(s),v(s))_H
    \end{align}
\end{lem}

\begin{proof}
    \begin{description}
    \item{1)}
    Es seies zunächst $u,v \in C(\ol I; V)$. Es gilt
    \begin{align*}
        (u(t),v(t))_H'&=(u'(t),v(t))_H + (u(t),v'(t))_H\\
        \Ra \, (u(t),v(t))_H-(u(s),v(s))_H &= \int_s^t (u',v)_H + (u,v')_H \overset{\text{Gelfand-Tri.}}
        {=} \int_s^t\lal \frac{\d}{\d t} u, v\ral_V+\lal u,\frac{\d}{\d t} v \ral _V
    \end{align*}
    \item{2)}
    Es sei nun $\vp\in C^1(\ol I, \R)$, mit $\vp(a)=0$, $\vp(b)=1$ mit $0<a<b<T$ und $u\in
    C^1(\ol I;V)$. Wir setzen $v=\vp u$, $w=u-\vp u$. Es gilt
    \begin{align*}
        v'&= \vp'u + \vp u' \\
        w'&=u'-\vp'u-\vp u'
    \end{align*}
    Nachdem \ref{33} bereits für $u,v,w\in C^1(\ol I; V)$ gilt, (Schritt 1) folgt
    \begin{align*}
        (v(t),u(t))_H&= \int_a^t\left\{ \vp'(s)(u(s), u(s))_H+ 2\vp\lal u'(s),u(s)\ral_V\right\} \d s\\
        (-w(t),u(t))_H &= \int_t^b\{ -\vp (s)(u(s),u(s))_H + 2 (1-\vp(s))\lal u'(s),u(s)\ral\} \d s
    \end{align*}
    Durch Subtraktion der beiden Gleichungen erhalten wir
    \[
        \|u(t)\|_H^2=\int_a^b \{\vp'(s)(u(s),u(s))+ 2 \vp(s)\lal u'(s),u(s)\ral_V\}\d s
        =2\int_a^b\lal u'(s),u(s)\ral _V\d s.
    \]
    Somit folgt
    \[
        \|u(t)\|_H^2\leq K\left( \|u\|_{C(\ol I; V)} \cdot \|u\|_{L^p(I;V)} + \left\| \frac{\d}{\d t}
        u\right\|_{L^{p'}(I;V')}\cdot\|u\|_{L^p(I;V)}\right)  \leq K\|u\|^2_W,
    \]
    wegen stetiger Einbettung von $W$ in $C(\ol I;V')$.
    \item{3)}
    Nun sei $u_n$ eine Cauchy-Folge in $C(I,V)$, diegegen $u\in W$ konvergiert. Es gilt
    \[
        \|u_n-u_k\|_{C(\ol I;H)} \leq K\|u_n-u_k\|_W \qq \text{(Schritt 2)}
    \]
    Damit konvergiert die Folge auch in $C(\ol I; H)$ gegen $u$. Die Normabschätzung bleibt im Limes
    erhalten, ebenso die Formel für die prtielle Integration \ref{33}.
    \end{description}
    \[ \]
\end{proof}

\begin{remark}
    \begin{description}
    \item{i)}
    \ref{33} ist analog zur reellwertigen partiellen Integrartion, $u,v:I\ra \R$,
    \[
        \int_s^tu'(\tau)v(\tau) + u(\tau)v'(\tau) \d \tau = u(t)v(t)-u(s)v(s)
    \]
    \item{ii)}
    Für $u=v\in W$ folgt
    \[
        \int_s^t \lal \frac{\d u}{\d t}(\tau), u(\tau)\ral_V \d\tau= \frac12 \|u(t)\|_H^2-\frac12\|u(s)\|
        _H^2
    \]
    \end{description}
\end{remark}

\subsubsection*{Das Anfangswertproblem, Existenz}

Wir betrachten das Problem
\begin{align*}
    \frac{\d u}{\d t} + Au &= b\\
    u(0)&= u_0
\end{align*}
und wollen eine Variante des \textit{Satzes von Brézis} beweisen.

\noindent Es sei also $A:V\ra V'$ ein Operator auf einen separablen, reflexiven Banachraum $V$. Es sei
$u\in L^p(I;V)$, $1<p<\infty$, $I=(0,T)$, $T<\infty$. Wir betrachten den induzierten Operator
\[
    (\tilde Au)(t) := A(u(t)), \qq t\in I.
\]
Unter bestimmeten Voraussetzungen gilt für den induzierten Operator, dass
\[
    \tilde A: L^p(I;V) \ra L^{p'}(I;V') = \left(L^p(I;V))\right)'    
\]
ist. In diesem Fall identifizieren wir $A$ mit $\tilde A$. Im Folgenden sei nun $X:= L^p(I;V)$,
$I=(0,T)$, $T<\infty$, $1<p<\infty$, $W'=L^{p'}(I;V)=(L^p(I;V))'$, $(V,H,V')$ sei ein Gelfand-Tripel.

\begin{theorem}\label{4.39} [Existenz von Lösungen für das Cauchy-Problem]
    Es sei $A:V\ra V'$ ein Operator, so dass der induzierte Operator $A:X\ra X'$ pseudomonoton und
    beschränkt ist, sowie der Koerzitivitätsbedingung
    \[
        \lal Au,u\ral _X\geq C_0\cdot \|u\|_x^p
    \]
    genügt. Dann existiert für alle $u_0\in H$, $b\in X'$ eine Lösung $u\in W$ von
    \[
        \frac{\d u}{\d t}+ Au=b
    \]
\end{theorem}

\begin{remark}
    Die Anfangswertbedingug ergibt wegen der Einbettung
    \[
        W\hookrightarrow C(\ol I;H)
    \]
    Sinn.
\end{remark}

\begin{proof}
    Dieser Beweis verläuft ähnlich dem Existenzbeweis für schwache Lösungen der Wärmeleitungsgleichung.
    Siehe z.B. \textit{Evans: Partial Differential Equations, Chapter 7}.

    \noindent Wir benutzen wieder das Galerkin-Verfahren. Es sei also $\{x_i\}_{i\in \N}\subset V$, so
    dass $\{w_i\}_{i=1}^n$ linear unabhängig sind und $\mr{span}\{w_i\}_{i\in\N}$ dicht liegt in $V$. Wir
    setzen $V_n=\mr{span}\{w_i\}_{i=1}^n$ und suchen eine Lösung $u_n$ der Form
    \[
        u_n(t)=\sum_{k=1}^n C^n_k(t) w_k,
    \]
    die für alle $t\in [0,T]$ das Galerkin-System.
    \begin{align}\label{34}
        \begin{split}
        \lal \frac{\d u_n}{\d t}(t), w_j\ral _V + \lal Au_n(t),w_j\ral _V &= \lal b_n(t), w_j\ral\\
        u_n(0)&= u_0^n
        \end{split}
    \end{align}
    für alle $j=1,…,n$ löst. Hier ist
    \[
        \begin{cases} b_n&\in C(\ol I; V') \\ b_n\ra b \qq \text{in} \, X'\end{cases}
    \]
    Eine solche Folge exisiert wegen Dichtheit von $C_c^\infty(\ol I;V')$ und
    \[
        u_0^n=\sum_{k=1}^n C_{0k}^nw_i \in V_n, 
    \]
    so dass $u_0^n\ra u_0$ in $H$.
    \begin{description}
    \item{1)}
        Lösbarkeit des Galerkin-Systems. Aufgrund der linearen Unabhängigkeit der $\{ w_i \}_{i=1}^n$ ist
        die Matrix $D=(d_{ij})_{i,j=1}^n=((w_i,w_j)_H)_{i,j=1}^n$ invertierbar. Somit können wir \ref{34}
        schreiben als System gewöhnlicher Differentialgleichungen für die Funktion
        \begin{align*}
            t\mapsto C^n(t)&=\begin{pmatrix} C_1^n(t) & \cdots & C_n^n(t)\end{pmatrix}\\
            \frac{\d C^n(t)}{\d t} &= f^n(t,C^n(t))\\
            C^n(0)&= C_0^n
        \end{align*}
        Wobei
        \[
            f_j^n(t,C)=\sum_{k=1}^n d_{jk}^{-1} (\lal b_n(t), w_k\ral_V-\lal A(\sum_{l=1}^nC_l^nw_l), w_k
            \ral_V)
        \]
        für $j=1,…,n$ und $C_0^n=(C_{01}^n, … , C_{0n}^n)$
    \end{description}\[ \]
\end{proof}
   % 

\end{document}
